\chapter{卡尔曼滤波与误差状态滤波}
\label{ch:eskf}

% ════════════════════════════════════════════════════════════════
%  P1 第二章（全吸收重建）：吸收
%  Barfoot SER2 第4章（非线性非高斯估计：Bayes滤波/EKF/广义高斯滤波/IEKF/
%    MC·线性化·sigma-point/UKF·SPKF/ISPKF（逐式算法盒 algo:spkf-*·ispkf）/
%    粒子滤波（SIR+proposal+N_eff判据+Madow重采样 algo:pf）/滤波器谱系/
%    批量MAP·GN/Laplace协方差/ML偏差/滑窗滤波SWF/KF≡批量MAP）
%  + Barfoot SER2 第2章补遗（稳态KF·离散代数Riccati DARE eq:dare 四条件/
%    连续时间 Kalman-Bucy eq:kb-* §3.5.2）
%  + Barfoot SER2 第5章（估计非理想因素：无偏与一致/NEES·NIS/偏置对KF的影响/
%    状态增广与可观测性/RANSAC/M估计·IRLS/协方差估计·逆Wishart=鲁棒核MAP）
%  + 十四讲第9章（贝叶斯滤波/KF的MAP配方法推导/信息形式/EKF/EKF-SLAM四局限/
%    滤波vs优化/H稀疏·Schur·边缘化·鲁棒核）
%  + Solà 1711.02508（四元数运动学 + ESKF：三态/真态·名义态·误差态运动学/
%    238b·238c逐步证明/离散雅可比F_x·F_i/校正·注入·重置雅可比 G/局部vs全局角误差）
%  + Thrun《Probabilistic Robotics》（直方图/网格滤波 eq:histogram-filter/
%    粒子滤波 N_eff·重要性采样·proposal·RBPF/信息滤波 IF prop:info-filter·SEIF）
%  + Särkkä《Bayesian Filtering and Smoothing》（粒子退化与 N_eff 判据）
%  + Barrau–Bonnabel 2017（不变卡尔曼滤波 InEKF：群仿射·(右)不变误差·误差自治/
%    log-linear => 一致性；扩展位姿群 SE_2(3) 结构要点，正式定义留 ch:lie）paper:inekf
%  + Mahony 2008 / Madgwick 2011（互补滤波：陀螺高通+加计低通的确定性姿态估计
%    eq:complementary-filter，作 ESKF 的轻量对照）
%  RTS 平滑取自线性高斯章脉络（块三对角 Cholesky 前向=KF、前向+后向=RTS）。
%  右扰动（局部角误差）为主；左扰动（全局角误差）并列。教学层遵循 v5.0 规范。
%  范围边界：IMU 误差模型/预积分的具体雅可比留 ch:imu/ch:vio；批量 GN/LM/BA 的
%  完整解算与 √SAM 留 ch:nlopt（本章只给与滤波对偶所需的骨架）。
% ════════════════════════════════════════════════════════════════

\begin{practice}[前置自测]
开始之前，先试答下列问题。若已能流畅作答，可只速读本章表格与速查卡；若卡住，正是本章要为你解决的。\;答错某题，可按括号指引回看对应小节。
\begin{enumerate}
  \item 卡尔曼滤波与上一章的“批量最小二乘”是两套不同的东西，还是同一个问题的两种算法？（\cref{sec:kf}）
  \item 卡尔曼增益 $\mathbf{K}$ 在“相信预测”与“相信观测”之间如何权衡？观测噪声趋于零时 $\mathbf{K}$ 会怎样？（\cref{sec:kf}）
  \item “一个高斯分布经过一个非线性函数之后还是高斯吗？”如果不是，EKF 凭什么把它当高斯处理？（\cref{sec:bayes-filter,sec:nonlin-pdf}）
  \item EKF 把非线性模型“线性化一次”就直接算后验，这一步可能出什么问题？怎么诊断它是否“过自信”？（\cref{sec:ekf,sec:consistency}）
  \item 为什么估计相机姿态时，常不把四元数直接塞进卡尔曼滤波的状态向量，而要用“误差状态”？（\cref{sec:eskf}）
  \item “滤波”与“BA/因子图优化”各自适合什么场景？为什么现代视觉 SLAM 后端多用后者？（\cref{sec:filter-vs-opt}）
\end{enumerate}
\end{practice}

\section*{本章目标}
学完本章，你应当能够：
\begin{enumerate}
  \item \textbf{推导}贝叶斯滤波的预测–校正递归，理解它为何“精确但一般不可实现”；
  \item \textbf{写出}线性高斯卡尔曼滤波（KF）的预测步与校正步，并能用 MAP 配方法与高斯条件化两条路各推一遍，解释卡尔曼增益、新息与信息形式；
  \item \textbf{证明} KF 是上一章批量 MAP 在马尔可夫结构下的\emph{精确递归等价}（非近似），并能说清滤波 / 平滑 / 滑窗三者关系；
  \item \textbf{推导} EKF 的线性化预测–校正与广义高斯滤波，说清它的三宗罪（工作点偏差、单次线性化、不一致），并能比较“一个高斯过非线性”的三种近似（蒙特卡洛 / 线性化 / sigma 点）；
  \item \textbf{定位}直方图（离散贝叶斯）、IEKF、UKF/SPKF、ISPKF、粒子滤波在滤波器谱系中的位置（IEKF$\to$众数，迭代 SPKF$\to$均值），并能写出 SPKF/ISPKF/粒子滤波（含 $N_{\text{eff}}$ 重采样判据）的逐式算法；
  \item \textbf{诊断}估计器的无偏与一致性（NEES/NIS），并理解偏置如何使 KF 系统性过自信、何时可用状态增广估出偏置（可观测性）；
  \item \textbf{实现}误差状态卡尔曼滤波（ESKF）：标称/误差/真值三态、预测、校正、注入、复位，能独立推出线速度误差与角误差动力学、复位雅可比，且用\emph{右扰动}（并能切换到左扰动）；并能说清互补滤波（轻量确定性姿态）与不变卡尔曼滤波 InEKF（群仿射 $+$ 不变误差 $\Rightarrow$ 一致性）相对 ESKF 的取舍；
  \item \textbf{比较}滤波与优化（平滑），用鲁棒核 / IRLS / 逆 Wishart 先验把“离群点”纳入估计，判断工程上何时用哪个。
\end{enumerate}

\subsection*{本章知识导航}
本章主线：把上一章“批量 MAP”\emph{递归化}，并逐步放宽到非线性（EKF 及其家族）与流形（ESKF），最后回到“滤波 vs 优化”的工程抉择。脉络层上，\textbf{骨干}是 \cref{sec:eskf-why,sec:bayes-filter,sec:kf,sec:ekf,sec:eskf}（递归动机$\to$贝叶斯$\to$KF$\to$EKF$\to$ESKF）与 \cref{sec:filter-vs-opt}（工程取舍）；\textbf{次优}是 \cref{sec:nonlin-pdf}（sigma 点 / UKF / 粒子滤波谱系）与 \cref{sec:consistency}（一致性诊断与偏置）；带（选读）标记的滑窗滤波、ML 偏差闭式、逆 Wishart 鲁棒核为\textbf{前沿/选读}。

\begin{center}
\small
\begin{tikzpicture}[node distance=7mm and 9mm, >=Stealth,
  blk/.style={draw, rounded corners, align=center, font=\small, inner sep=3pt, fill=main!6, draw=main!60!black},
  grp/.style={draw, rounded corners, align=center, font=\small, inner sep=3pt, fill=second!6, draw=second!60!black},
  fwd/.style={draw, rounded corners, align=center, font=\small, inner sep=3pt, fill=third!8, draw=third!60!black},
  ln/.style={->, thick, gray!60!black}]
\node[blk] (batch) {批量 MAP\\{\scriptsize（\cref{ch:slam_est}）}};
\node[blk, right=of batch] (bayes) {贝叶斯滤波\\{\scriptsize 预测–校正}};
\node[blk, right=of bayes] (kf) {卡尔曼滤波 KF\\{\scriptsize 线性高斯·最优}};
\node[grp, below=of kf] (ekf) {EKF / 谱系\\{\scriptsize 线性化·sigma点}};
\node[grp, left=of ekf] (eskf) {ESKF\\{\scriptsize 误差状态·右扰动}};
\node[fwd, left=of eskf] (cmp) {滤波 vs 优化\\{\scriptsize 工程取舍·鲁棒}};
\draw[ln] (batch)--(bayes); \draw[ln] (bayes)--(kf);
\draw[ln] (kf)--(ekf) node[midway, right, font=\scriptsize]{非线性};
\draw[ln] (ekf)--(eskf) node[midway, above, font=\scriptsize]{流形};
\draw[ln] (eskf)--(cmp);
\end{tikzpicture}
\end{center}

\noindent\textbf{推荐路径}：第 1--3 节（递归动机$\to$贝叶斯滤波$\to$KF）是主干，其中 \cref{thm:kf-map}（KF$\equiv$批量 MAP）是与上一章缝合的关键；第 4--5 节把高斯过非线性的三种办法讲透，给出 EKF 与整个滤波器谱系；第 6 节诊断一致性与偏置；第 7 节 ESKF 面向位姿/姿态估计，是 VIO、惯导融合（\cref{ch:imu}、\cref{ch:vio}）的工作母机；第 8 节给“滤波 vs 优化”的工程结论与鲁棒化。带 IMU 噪声标定/预积分的细节属 \cref{ch:imu}/\cref{ch:vio}，本章只用抽象运动模型。

\subsection*{前置知识桥接}
\cref{ch:slam_est} 把 SLAM 钉成“批量 MAP $=$ 最小二乘”，并给了正规方程 $(\mathbf{H}^\top\mathbf{W}^{-1}\mathbf{H})\hat{\mathbf{x}}=\mathbf{H}^\top\mathbf{W}^{-1}\mathbf{z}$（其非线性增量形式 $(\mathbf{H}^\top\mathbf{W}^{-1}\mathbf{H})\delta\mathbf{x}=\mathbf{H}^\top\mathbf{W}^{-1}\mathbf{e}$ 见 \cref{ch:nlopt} 的高斯牛顿）、信息矩阵 $\boldsymbol{\Lambda}=\boldsymbol{\Sigma}^{-1}$、协方差$\leftrightarrow$信息对偶、舒尔补条件化与边缘化、矩阵求逆引理（SMW）——本章反复使用这些（“回顾一句：把联合高斯按舒尔补条件化，就得到条件均值与条件协方差，这正是 KF 校正步的代数核心”）。\cref{ch:lie}、\cref{ch:rigid_body} 提供 ESKF 所需的右扰动 $\mathbf{R}\,\mathrm{Exp}(\delta\boldsymbol{\phi})$、右雅可比 $\mathbf{J}_r$、Hamilton 四元数 $\mathbf{q}$（实部在前）、$\otimes$ 乘法与 $\mathrm{Exp}/\mathrm{Log}$。\cref{ch:nlopt} 的高斯牛顿（GN）与 Levenberg–Marquardt（LM）是本章“批量 = 优化”与 IRLS 的承接。

\subsection*{如果跳过本章会怎样}
\begin{itemize}
  \item 你会把 KF 的“预测–更新”当孤立公式背，而看不出它就是 \cref{ch:slam_est} 批量 MAP 的递归重排（\cref{thm:kf-map}）；
  \item 在 \cref{ch:vio} 写 VIO 时，Ceres/GTSAM 的“误差状态 / 局部参数化”会显得凭空冒出——其实就是本章 ESKF 的右扰动（\cref{sec:eskf}）；
  \item 你将不会诊断一个滤波器是否“过自信”（不一致），也不知道为何一个未补偿的加速度计偏置会随时间把协方差越压越乐观直至发散（\cref{sec:consistency}）；
  \item 你将无法判断一个工程该上滤波还是上因子图优化，也不会用鲁棒核把误匹配挡在估计之外（\cref{sec:filter-vs-opt}）。
\end{itemize}

\begin{table}[H]\centering\small
\caption{本章阅读方式建议}
\begin{tabular}{lll}
\toprule
阅读方式 & 时间 & 适合谁 \\
\midrule
精读（含全部推导与练习） & 6--8 小时 & 要做滤波/VIO 推导与实现 \\
速读（跳过谱系、SWF、ML 偏差细节） & 2 小时 & 已懂 KF、想抓 ESKF 与右扰动 \\
速查（只看小结三张表 + ESKF 算法卡） & 15 分钟 & 写代码时回来查公式 \\
\bottomrule
\end{tabular}
\end{table}

\begin{note}
\textbf{本章记号与约定（兼记号约定区）.}\;
状态 $\mathbf{x}_k$；线性系统转移矩阵 $\mathbf{A}_{k-1}$、观测矩阵 $\mathbf{C}_k$；非线性时运动 $\mathbf{f}$、观测 $\mathbf{h}$（部分源记观测为 $\mathbf{g}$，本书统一 $\mathbf{h}$），其雅可比 $\mathbf{F}=\partial\mathbf{f}/\partial\mathbf{x}|_{\hat{\mathbf{x}}_{k-1}}$、$\mathbf{H}=\partial\mathbf{h}/\partial\mathbf{x}|_{\check{\mathbf{x}}_k}$。\textbf{输入}记 $\mathbf{u}_k$，\textbf{观测噪声}记 $\mathbf{v}_k\sim\mathcal{N}(\mathbf{0},\boldsymbol{\Sigma}_{\mathbf{v}})$、\textbf{过程噪声} $\mathbf{w}_k\sim\mathcal{N}(\mathbf{0},\boldsymbol{\Sigma}_{\mathbf{w}})$。先验/预测 $\check{(\cdot)}$、后验/校正 $\hat{(\cdot)}$；状态协方差 $\mathbf{P}$（先验 $\check{\mathbf{P}}$、后验 $\hat{\mathbf{P}}$）；信息矩阵 $\boldsymbol{\Lambda}=\boldsymbol{\Sigma}^{-1}$；卡尔曼增益 $\mathbf{K}_k$；新息协方差 $\mathbf{S}_k$。
\emph{字母隔离声明}：本章用 $\mathbf{F},\mathbf{H}$ 作\emph{雅可比}（与 \cref{ch:nlopt} 的 BA 雅可比块同名但语义不同，见各处就地说明）；$\mathbf{R}$ 一律留给\emph{旋转矩阵}，绝不用作观测噪声协方差。

\textbf{跨书雷区（务必记牢）}：Barfoot\cite{barfoot2024state} 用 $\mathbf{v}_k$ 表示\emph{输入}、$\mathbf{n}_k$ 表示\emph{观测噪声}、$\mathbf{Q}/\mathbf{R}$ 为过程/观测协方差、$\mathbf{C}$ 为观测矩阵、$\mathbf{A}$ 为转移矩阵；十四讲\cite{gaoxiang2019slam14} 的过程/观测协方差字母与 Barfoot \emph{相反}（过程 $\mathbf{R}$、观测 $\mathbf{Q}$）；Solà\cite{sola2017quaternion} 用 $\mathbf{V}$ 表观测噪声协方差（因 $\mathbf{R}$ 被旋转占用）。本书统一：输入 $\mathbf{u}$、观测噪声 $\mathbf{v}$、协方差 $\boldsymbol{\Sigma}_{\mathbf{w}}/\boldsymbol{\Sigma}_{\mathbf{v}}$，把字母 $\mathbf{R}$ 留给旋转矩阵。引用三书公式时按此映射，不再每处声明。
\end{note}

% ════════════════════════════════════════════════════════════════
\section{从批量到递归：滤波的动机}
\label{sec:eskf-why}

\paragraph{动机：数据是流式到来的。}
上一章的批量 MAP 把所有时刻一次解出，适合事后处理，却\emph{不能实时}：每来一帧都重解整条轨迹，代价随轨迹长度增长。但 SLAM 的数据是\emph{逐时刻}到来的，我们希望持有“当前估计”，用新数据\emph{增量}更新它。这就是\emph{递归/滤波}。十四讲\cite{gaoxiang2019slam14} 把这两条路明确为：把状态分布“从 $k-1$ 时刻推到 $k$ 时刻、只维护均值与协方差”的\emph{渐进（incremental）}做法，对应滤波；把“不仅用过去、也用未来信息更新过去状态”的\emph{批量（batch）}做法，对应平滑/优化。

\paragraph{反面：不递归会怎样。}
若坚持每来一帧就重解整条 $\{\mathbf{x}_0,\dots,\mathbf{x}_k\}$，则第 $k$ 步要解一个维度 $\propto k$ 的线性系统，朴素代价 $O(k^3)$，累计到第 $K$ 步是 $O(K^4)$。一个跑几分钟的里程计就会卡死。\textbf{递归的全部意义}就是把“重解整条轨迹”换成“在固定大小的状态上做一次常数代价的更新”。下文会证明，在线性高斯（LG）下这个“省下来”是\emph{免费的}——递归与批量给出\emph{完全相同}的答案。

\paragraph{蒙眼走路：预测增大不确定度，观测减小它。}
一个贯穿全章的直觉\cite{gaoxiang2019slam14}：
\begin{quote}
只有运动方程时，像\emph{蒙着眼}在未知地方走路——尽管知道每步走多远，但随时间流逝越来越不确定自己在哪，位置方差越来越大。\emph{睁开眼}不断观测外部场景，位置估计的不确定性就会变小，越来越自信。
\end{quote}
若用协方差椭圆表示不确定度：无观测时它随运动\emph{膨胀}；有正确观测时它\emph{收缩}并稳定。滤波就是这“预测（膨胀）—校正（收缩）”的不断交替。

\begin{figure}[ht]
\centering
\begin{tikzpicture}[>=Stealth, scale=1.0]
  % trajectory
  \draw[->, thick, gray!70!black] (-0.3,0) -- (8.3,0);
  \foreach \x/\lab in {0/0,2/1,4/2,6/3} { \fill (\x,0) circle (1.2pt) node[below=2pt, font=\scriptsize]{$k=\lab$}; }
  % predict grows ellipse, correct shrinks
  \draw[second, thick] (0,0) ellipse (0.35 and 0.3);
  \draw[main!70!black, thick, dashed] (2,0) ellipse (0.75 and 0.55);
  \draw[second, thick] (2,0) ellipse (0.4 and 0.32);
  \draw[main!70!black, thick, dashed] (4,0) ellipse (0.8 and 0.58);
  \draw[second, thick] (4,0) ellipse (0.42 and 0.33);
  \draw[main!70!black, thick, dashed] (6,0) ellipse (0.82 and 0.6);
  \draw[second, thick] (6,0) ellipse (0.43 and 0.34);
  \node[main!70!black, font=\scriptsize] at (3,1.0) {预测：椭圆膨胀};
  \node[second, font=\scriptsize] at (5.2,-1.0) {校正：椭圆收缩};
\end{tikzpicture}
\caption{滤波的“预测（虚线，膨胀）—校正（实线，收缩）”交替：协方差椭圆随运动放大、随观测缩小并趋于稳定。（仿 \cite{gaoxiang2019slam14} 重绘）}
\label{fig:eskf-predict-correct}
\end{figure}

\paragraph{为什么递归是可能的：马尔可夫性 $\Rightarrow$ 稀疏 $\Rightarrow$ 前向消元。}
上一章已指出：因运动只耦合相邻时刻（一阶马尔可夫）、观测只耦合同时刻，批量信息矩阵 $\mathbf{H}^\top\mathbf{W}^{-1}\mathbf{H}$ 是\emph{块三对角}的。对块三对角矩阵做 Cholesky 分解，其因子保持\emph{下双对角}，于是“前代—回代”各只需 $O(K)$ 步\cite{barfoot2024state}。
\begin{insight}
\textbf{关键洞见（本章灵魂）}：对块三对角批量正规方程做稀疏 Cholesky，其\emph{前向}消元逐式等于\textbf{卡尔曼滤波}，\emph{前向 $+$ 后向}等于 \textbf{RTS 平滑器}。所以“滤波”不是另起炉灶，而是“批量 MAP”在因果约束下的精确递归算法。\cref{thm:kf-map} 将把这句话坐实。
\end{insight}

\paragraph{滤波 vs 平滑：因果性之别。}
\emph{滤波器}（filter）在时刻 $k$ 只用 $\le k$ 的数据估计 $\mathbf{x}_k$——\emph{因果}、可在线。\emph{平滑器}（smoother）允许用\emph{未来}数据估计过去状态（“站在未来回看你该在哪”）——\emph{非因果}，需离线或滞后。批量/RTS 是平滑，KF 是滤波；二者在 $k=K$、用满全部数据时给出同一估计。Barfoot\cite{barfoot2024state} 把线性高斯的完整后验写成“前向 $\times$ 后向”两个因子之积
\begin{equation}
  p(\mathbf{x}_k\mid\mathbf{u},\mathbf{z})=\eta\,
  \underbrace{p(\mathbf{x}_k\mid\check{\mathbf{x}}_0,\mathbf{u}_{1:k},\mathbf{z}_{0:k})}_{\text{前向（滤波）}}\,
  \underbrace{p(\mathbf{x}_k\mid\mathbf{u}_{k+1:K},\mathbf{z}_{k+1:K})}_{\text{后向（未来信息）}},
  \label{eq:eskf-fwd-bwd}
\end{equation}
本章先把“前向”因子做成递归滤波器（\cref{sec:bayes-filter}），再在 \cref{sec:kf} 用后向修正补成平滑器。这一因式分解并非凭空断言：把批量 $\mathbf{H}$ 按时刻分 $12$ 块、用一次行操作隔离 $\mathbf{x}_k$，可\emph{显式}化简出 $\hat{\mathbf{P}}_k^{-1}=\hat{\mathbf{P}}_{k,f}^{-1}+\hat{\mathbf{P}}_{k,b}^{-1}$（前向 $+$ 后向两估计器的信息相加），逐式见 \cref{der:factor-batch}。

% ════════════════════════════════════════════════════════════════
\section{贝叶斯滤波：递归的总框架}
\label{sec:bayes-filter}

\paragraph{动机：把“估计现在的分布”写成对“上一时刻分布”的递推。}
设系统满足一阶马尔可夫的运动/观测模型
\begin{equation}
  \mathbf{x}_k=\mathbf{f}(\mathbf{x}_{k-1},\mathbf{u}_k,\mathbf{w}_k),\qquad
  \mathbf{z}_k=\mathbf{h}(\mathbf{x}_k,\mathbf{v}_k).
  \label{eq:bf-model}
\end{equation}
注意噪声\emph{可以进入非线性内部}（不必加性），这是 Barfoot\cite{barfoot2024state} 的一般写法；加性是其特例 $\mathbf{x}_k=\mathbf{f}(\mathbf{x}_{k-1},\mathbf{u}_k)+\mathbf{w}_k$、$\mathbf{z}_k=\mathbf{h}(\mathbf{x}_k)+\mathbf{v}_k$，可经线性化（近似）恢复。目标是\emph{置信度}（belief）$p(\mathbf{x}_k\mid\check{\mathbf{x}}_0,\mathbf{u}_{1:k},\mathbf{z}_{0:k})$——即在已知初始先验、全部输入与全部观测下、当前状态的\emph{整个} PDF（不仅是点估计）。这是 Jazwinski 与 Bayes 滤波的目标。

\begin{theorem}[贝叶斯滤波]\label{thm:bayes-filter}
置信度满足递归
\begin{equation}
  \underbrace{p(\mathbf{x}_k\mid\mathbf{u}_{1:k},\mathbf{z}_{0:k})}_{\text{后验置信度}}
  =\eta\,\underbrace{p(\mathbf{z}_k\mid\mathbf{x}_k)}_{\text{观测}}
  \int \underbrace{p(\mathbf{x}_k\mid\mathbf{x}_{k-1},\mathbf{u}_k)}_{\text{运动}}\,
  \underbrace{p(\mathbf{x}_{k-1}\mid\mathbf{u}_{1:k-1},\mathbf{z}_{0:k-1})}_{\text{先验置信度}}\,\mathrm{d}\mathbf{x}_{k-1},
  \label{eq:bayes-filter}
\end{equation}
$\eta$ 为归一化常数（为简洁，记号上省去对初始先验 $\check{\mathbf{x}}_0$ 的条件依赖，它始终被携带）。
\end{theorem}
\begin{proof}
四步，每步不跳（Barfoot\cite{barfoot2024state} §4.2.2、十四讲\cite{gaoxiang2019slam14} §9.1.1）。

\emph{(1) 分出最新观测.} 因各观测在给定状态后条件独立，用贝叶斯规则把 $\mathbf{z}_k$ 与 $\mathbf{x}_k$ 互换：
\[
  p(\mathbf{x}_k\mid\check{\mathbf{x}}_0,\mathbf{u}_{1:k},\mathbf{z}_{0:k})
  =\eta\,p(\mathbf{z}_k\mid\mathbf{x}_k)\,p(\mathbf{x}_k\mid\check{\mathbf{x}}_0,\mathbf{u}_{1:k},\mathbf{z}_{0:k-1}).
\]

\emph{(2) 引入隐状态 $\mathbf{x}_{k-1}$ 并积分（边缘化的逆，无近似）.}
\[
  p(\mathbf{x}_k\mid\cdots,\mathbf{z}_{0:k-1})
  =\int p(\mathbf{x}_k\mid\mathbf{x}_{k-1},\check{\mathbf{x}}_0,\mathbf{u}_{1:k},\mathbf{z}_{0:k-1})\,
   p(\mathbf{x}_{k-1}\mid\check{\mathbf{x}}_0,\mathbf{u}_{1:k},\mathbf{z}_{0:k-1})\,\mathrm{d}\mathbf{x}_{k-1}.
\]

\emph{(3) 用马尔可夫性化简（此步微妙，是递归诸多局限之源）.} 给定 $\mathbf{x}_{k-1}$ 与最新输入 $\mathbf{u}_k$，$\mathbf{x}_k$ 与更早历史无关；而 $\mathbf{x}_{k-1}$ 又与未来输入 $\mathbf{u}_k$ 无关：
\[
  p(\mathbf{x}_k\mid\mathbf{x}_{k-1},\dots)=p(\mathbf{x}_k\mid\mathbf{x}_{k-1},\mathbf{u}_k),\qquad
  p(\mathbf{x}_{k-1}\mid\mathbf{u}_{1:k},\mathbf{z}_{0:k-1})=p(\mathbf{x}_{k-1}\mid\mathbf{u}_{1:k-1},\mathbf{z}_{0:k-1}).
\]

\emph{(4) 回代}即得 \cref{eq:bayes-filter}。首步 $k=0$ 是只含 $\mathbf{z}_0$ 校正的特例，约定滤波器以 $p(\mathbf{x}_0\mid\check{\mathbf{x}}_0,\mathbf{z}_0)$ 初始化。
\end{proof}

\noindent \cref{eq:bayes-filter} 是\emph{预测–校正}结构：积分项用运动模型把先验置信度\emph{前推}成预测；外面乘观测似然 $p(\mathbf{z}_k\mid\mathbf{x}_k)$ 完成\emph{校正}。核心难点被它精确地“暴露”出来了——\textbf{如何让一个 PDF 穿过非线性函数 $\mathbf{f},\mathbf{h}$}。这正是 \cref{sec:nonlin-pdf} 的主题。

\begin{pitfall}[概念误区：贝叶斯滤波“精确但不可实现”]
\cref{thm:bayes-filter} 是精确的，但只是数学产物，一般情形\emph{不能直接算}\cite{barfoot2024state}：
\begin{enumerate}
  \item \textbf{内存问题}：置信度是\emph{无穷维} PDF（活在所有连续函数的空间），需无穷参数才能完整表示。对策：近似表示——近似为高斯（只存前两阶矩，KF/EKF/UKF），或用有限随机样本表示（粒子滤波）。
  \item \textbf{计算问题}：那个积分极昂贵，精确求需无穷算力。对策：近似求积分——线性化模型后闭式求积，或蒙特卡洛积分。
\end{enumerate}
\textbf{而一旦近似，严格的马尔可夫性就被破坏了}——这正是后面理解 EKF 短板、以及“批量优化是更优起点”的伏笔（\cref{sec:filter-vs-opt}）。一句话：所有实用滤波器都是 \cref{eq:bayes-filter} 的某种近似，差别只在“近似 PDF 的方式”与“近似积分的方式”。
\end{pitfall}

\paragraph{最朴素的可实现版：直方图（离散贝叶斯）滤波。}
高斯近似（KF/EKF/UKF，下两节）与样本近似（粒子滤波，\cref{sec:nonlin-pdf}）之间，还夹着一支\emph{最朴素}却最先“能算”的实现——\textbf{直方图滤波}（histogram filter，又称网格滤波 grid filter，或离散贝叶斯滤波）\cite{thrun2005probabilistic}。它的近似不在“压成高斯”也不在“撒样本”，而是把\emph{状态空间本身离散化}成有限个互不重叠的小格 $\{\mathcal{X}_1,\dots,\mathcal{X}_n\}$，每格 $i$ 配一个概率质量 $p_{k,i}\approx\Pr(\mathbf{x}_k\in\mathcal{X}_i)$。于是 \cref{eq:bayes-filter} 的无穷维 PDF 退化成一张\emph{有限概率表}，积分退化成\emph{有限求和}，预测与校正都成了对这张表的逐格更新（运动模型给格间转移概率 $p(\mathcal{X}_i\mid\mathcal{X}_j,\mathbf{u}_k)$，观测模型给每格似然 $p(\mathbf{z}_k\mid\mathcal{X}_i)$）：
\begin{equation}
  \text{预测：}\ \check{p}_{k,i}=\sum_{j}p(\mathcal{X}_i\mid\mathcal{X}_j,\mathbf{u}_k)\,\hat{p}_{k-1,j};\qquad
  \text{校正：}\ \hat{p}_{k,i}=\eta\,p(\mathbf{z}_k\mid\mathcal{X}_i)\,\check{p}_{k,i}.
  \label{eq:histogram-filter}
\end{equation}
\textbf{它是“非参数”的}（不假设单峰/高斯），故能表示\emph{多峰}置信度——这正是\emph{栅格定位}（Markov localization，\cref{ch:intro} 占据栅格地图的天然搭档）能处理“机器人可能在好几个长得一样的走廊口之一”这类全局歧义的原因：每个候选位置都保有一份概率质量，直到观测把假峰压下去。代价是维度灾难：格数随状态维\emph{指数}增长，故直方图滤波只在低维（如 $2$\,D 平面位置、或 $(x,y,\theta)$）实用，高维要么改自适应分辨率网格、要么让位给粒子滤波（用样本代替规则格、把算力花在高概率区）。

\begin{insight}
\textbf{“离散$\to$高斯$\to$样本”的下降链}：贝叶斯滤波 \cref{eq:bayes-filter} 精确但不可算，三种实用近似按“如何表示那个 PDF”排成一条谱（见 \cref{tab:filter-taxo}）：\emph{离散直方图}（有限格、任意形状、维度灾难）$\to$\emph{高斯}（前两阶矩、单峰、最省、KF/EKF/UKF）$\to$\emph{样本}（粒子、任意形状、维度灾难但靠重要性采样把算力集中到高概率区）。直方图是“能算”的起点与多峰非参数的最简范例，高斯是工程主力，粒子是高维非高斯的归宿——三者都不过是同一条 \cref{eq:bayes-filter} 的不同有限化。
\end{insight}

\begin{practice}
\begin{enumerate}
  \item 写出 \cref{thm:bayes-filter} 证明第 (1) 步用到的贝叶斯规则，并说明 $\eta$ 为什么可以“先不管、最后归一化”。
  \item （概念）解释为什么第 (3) 步的两条马尔可夫化简在“近似实现”后不再严格成立。答案要点：近似把后验从精确 PDF 换成高斯/样本，预测—校正不再是精确的条件分布，历史信息被“压扁”进有限参数。答错回看本节 pitfall 与 \cref{sec:filter-vs-opt}。
\end{enumerate}
\end{practice}

% ════════════════════════════════════════════════════════════════
\section{卡尔曼滤波：线性高斯的最优解}
\label{sec:kf}

线性高斯（LG）系统是 \cref{eq:bf-model} 的特例：
\begin{equation}
  \mathbf{x}_k=\mathbf{A}_{k-1}\mathbf{x}_{k-1}+\mathbf{u}_k+\mathbf{w}_k,\qquad
  \mathbf{z}_k=\mathbf{C}_k\mathbf{x}_k+\mathbf{v}_k,
  \label{eq:lg-model}
\end{equation}
$\mathbf{w}_k\sim\mathcal{N}(\mathbf{0},\boldsymbol{\Sigma}_{\mathbf{w}})$，$\mathbf{v}_k\sim\mathcal{N}(\mathbf{0},\boldsymbol{\Sigma}_{\mathbf{v}})$。\textbf{承重假设（三问）}：高斯先验 + 线性模型 + 高斯噪声。用在哪一步？——用在“线性变换把高斯映成高斯、两个高斯相乘还是高斯”这两处（下证），从而置信度\emph{精确}保持高斯、贝叶斯滤波闭式可解。删了崩在哪？——若模型非线性，预测/校正后置信度不再高斯，\cref{eq:bayes-filter} 的积分与乘积都没有闭式（这正是 EKF 要近似的地方，\cref{sec:ekf}）。可否弱化？——可去掉高斯假设：只要噪声零均值、不相关、方差已知，KF 仍是\emph{最优线性无偏估计}（BLUE，下详），但不再是“完整后验”。

\paragraph{招牌推导一：MAP 配方法（十四讲路线，最直观）。}\label{par:kf-map-route}
十四讲\cite{gaoxiang2019slam14} 把 KF 推导本身就建立在 MAP（最大后验）上：后验 $\propto$ 似然 $\times$ 先验，两侧都是高斯，\emph{只比较指数里的二次型}并配方。先做\emph{预测}：线性变换下高斯仍是高斯，
\begin{equation}
  \check{\mathbf{x}}_k=\mathbf{A}_{k-1}\hat{\mathbf{x}}_{k-1}+\mathbf{u}_k,\qquad
  \check{\mathbf{P}}_k=\mathbf{A}_{k-1}\hat{\mathbf{P}}_{k-1}\mathbf{A}_{k-1}^\top+\boldsymbol{\Sigma}_{\mathbf{w}}.
  \label{eq:kf-map-predict}
\end{equation}
观测似然 $p(\mathbf{z}_k\mid\mathbf{x}_k)=\mathcal{N}(\mathbf{C}_k\mathbf{x}_k,\boldsymbol{\Sigma}_{\mathbf{v}})$。设后验 $\mathcal{N}(\hat{\mathbf{x}}_k,\hat{\mathbf{P}}_k)\propto\mathcal{N}(\mathbf{C}_k\mathbf{x}_k,\boldsymbol{\Sigma}_{\mathbf{v}})\cdot\mathcal{N}(\check{\mathbf{x}}_k,\check{\mathbf{P}}_k)$，对齐两侧关于 $\mathbf{x}_k$ 的指数二次型：
\begin{equation}
  (\mathbf{x}_k-\hat{\mathbf{x}}_k)^\top\hat{\mathbf{P}}_k^{-1}(\mathbf{x}_k-\hat{\mathbf{x}}_k)
  =(\mathbf{z}_k-\mathbf{C}_k\mathbf{x}_k)^\top\boldsymbol{\Sigma}_{\mathbf{v}}^{-1}(\mathbf{z}_k-\mathbf{C}_k\mathbf{x}_k)
  +(\mathbf{x}_k-\check{\mathbf{x}}_k)^\top\check{\mathbf{P}}_k^{-1}(\mathbf{x}_k-\check{\mathbf{x}}_k).
  \label{eq:kf-quadform}
\end{equation}
\emph{比较 $\mathbf{x}_k$ 的二次项系数}（左右各取 $\mathbf{x}_k^\top(\cdot)\mathbf{x}_k$ 的中间块）：
\begin{equation}
  \hat{\mathbf{P}}_k^{-1}=\mathbf{C}_k^\top\boldsymbol{\Sigma}_{\mathbf{v}}^{-1}\mathbf{C}_k+\check{\mathbf{P}}_k^{-1}.
  \label{eq:kf-info-P}
\end{equation}
这是\textbf{信息形式的协方差更新}：后验信息 $=$ 先验信息 $+$ 观测信息 $\mathbf{C}_k^\top\boldsymbol{\Sigma}_{\mathbf{v}}^{-1}\mathbf{C}_k$。定义增益雏形 $\mathbf{K}=\hat{\mathbf{P}}_k\mathbf{C}_k^\top\boldsymbol{\Sigma}_{\mathbf{v}}^{-1}$；在 \cref{eq:kf-info-P} 两边右乘 $\check{\mathbf{P}}_k$ 并左乘 $\hat{\mathbf{P}}_k$：
\begin{equation}
  \mathbf{I}=\hat{\mathbf{P}}_k\mathbf{C}_k^\top\boldsymbol{\Sigma}_{\mathbf{v}}^{-1}\mathbf{C}_k+\hat{\mathbf{P}}_k\check{\mathbf{P}}_k^{-1}
  =\mathbf{K}\mathbf{C}_k+\hat{\mathbf{P}}_k\check{\mathbf{P}}_k^{-1}
  \;\Longrightarrow\;\hat{\mathbf{P}}_k=(\mathbf{I}-\mathbf{K}\mathbf{C}_k)\check{\mathbf{P}}_k.
  \label{eq:kf-Phat}
\end{equation}
\emph{比较 $\mathbf{x}_k$ 的一次项系数} $-2\hat{\mathbf{x}}_k^\top\hat{\mathbf{P}}_k^{-1}\mathbf{x}_k=-2\mathbf{z}_k^\top\boldsymbol{\Sigma}_{\mathbf{v}}^{-1}\mathbf{C}_k\mathbf{x}_k-2\check{\mathbf{x}}_k^\top\check{\mathbf{P}}_k^{-1}\mathbf{x}_k$，转置整理得\textbf{信息形式的均值更新}：
\begin{equation}
  \hat{\mathbf{P}}_k^{-1}\hat{\mathbf{x}}_k=\mathbf{C}_k^\top\boldsymbol{\Sigma}_{\mathbf{v}}^{-1}\mathbf{z}_k+\check{\mathbf{P}}_k^{-1}\check{\mathbf{x}}_k.
  \label{eq:kf-info-x}
\end{equation}
两边左乘 $\hat{\mathbf{P}}_k$，代入 $\mathbf{K}=\hat{\mathbf{P}}_k\mathbf{C}_k^\top\boldsymbol{\Sigma}_{\mathbf{v}}^{-1}$ 与 \cref{eq:kf-Phat} 的 $\hat{\mathbf{P}}_k\check{\mathbf{P}}_k^{-1}=\mathbf{I}-\mathbf{K}\mathbf{C}_k$：
\begin{equation}
  \hat{\mathbf{x}}_k=\mathbf{K}\mathbf{z}_k+(\mathbf{I}-\mathbf{K}\mathbf{C}_k)\check{\mathbf{x}}_k=\check{\mathbf{x}}_k+\mathbf{K}(\mathbf{z}_k-\mathbf{C}_k\check{\mathbf{x}}_k).
  \label{eq:kf-xhat}
\end{equation}
这就是后验均值 $=$ 预测 $+$ 增益 $\times$ 新息。

\begin{theorem}[卡尔曼滤波]\label{thm:kf}
设后验 $p(\mathbf{x}_{k-1}\mid\cdot)=\mathcal{N}(\hat{\mathbf{x}}_{k-1},\hat{\mathbf{P}}_{k-1})$，则时刻 $k$ 的预测与校正为
\begin{align}
  \text{预测：}\quad
  &\check{\mathbf{x}}_k=\mathbf{A}_{k-1}\hat{\mathbf{x}}_{k-1}+\mathbf{u}_k, &
  &\check{\mathbf{P}}_k=\mathbf{A}_{k-1}\hat{\mathbf{P}}_{k-1}\mathbf{A}_{k-1}^\top+\boldsymbol{\Sigma}_{\mathbf{w}}; \label{eq:kf-predict}\\
  \text{增益：}\quad
  &\mathbf{K}_k=\check{\mathbf{P}}_k\mathbf{C}_k^\top\big(\underbrace{\mathbf{C}_k\check{\mathbf{P}}_k\mathbf{C}_k^\top+\boldsymbol{\Sigma}_{\mathbf{v}}}_{=\,\mathbf{S}_k\ \text{新息协方差}}\big)^{-1}; \label{eq:kf-gain}\\
  \text{校正：}\quad
  &\hat{\mathbf{x}}_k=\check{\mathbf{x}}_k+\mathbf{K}_k\big(\mathbf{z}_k-\mathbf{C}_k\check{\mathbf{x}}_k\big), &
  &\hat{\mathbf{P}}_k=(\mathbf{I}-\mathbf{K}_k\mathbf{C}_k)\check{\mathbf{P}}_k. \label{eq:kf-correct}
\end{align}
其中 $\mathbf{z}_k-\mathbf{C}_k\check{\mathbf{x}}_k$ 称\emph{新息}（innovation）。\cref{eq:kf-gain} 的增益（协方差形式）与上文 $\mathbf{K}=\hat{\mathbf{P}}_k\mathbf{C}_k^\top\boldsymbol{\Sigma}_{\mathbf{v}}^{-1}$（信息形式）经矩阵求逆引理（SMW）等价。
\end{theorem}
\begin{proof}[招牌推导二：高斯条件化（Barfoot 路线）]
Barfoot\cite{barfoot2024state} 用“一开始就假设估计为高斯”的\emph{广义高斯滤波}更干净（其一般版即 \cref{prop:ggf}，并在 \cref{sec:nonlin-pdf} 用 sigma 点实例化）。预测同 \cref{eq:kf-predict}（对 \cref{eq:lg-model} 取期望与外积，交叉项因 $\mathbf{w}_k$ 独立于 $\mathbf{x}_{k-1}$ 而为零）。校正写状态与观测的\emph{联合高斯}：
\[
  p(\mathbf{x}_k,\mathbf{z}_k\mid\cdot)=\mathcal{N}\!\left(\begin{bmatrix}\check{\mathbf{x}}_k\\\mathbf{C}_k\check{\mathbf{x}}_k\end{bmatrix},
  \begin{bmatrix}\check{\mathbf{P}}_k & \check{\mathbf{P}}_k\mathbf{C}_k^\top\\ \mathbf{C}_k\check{\mathbf{P}}_k & \mathbf{C}_k\check{\mathbf{P}}_k\mathbf{C}_k^\top+\boldsymbol{\Sigma}_{\mathbf{v}}\end{bmatrix}\right).
\]
代入 \cref{ch:slam_est} 的舒尔补条件化公式 $\boldsymbol{\mu}_{x}+\boldsymbol{\Sigma}_{xy}\boldsymbol{\Sigma}_{yy}^{-1}(\mathbf{z}_k-\boldsymbol{\mu}_y)$、$\boldsymbol{\Sigma}_{xx}-\boldsymbol{\Sigma}_{xy}\boldsymbol{\Sigma}_{yy}^{-1}\boldsymbol{\Sigma}_{yx}$，令 $\mathbf{K}_k=\boldsymbol{\Sigma}_{xy}\boldsymbol{\Sigma}_{yy}^{-1}=\check{\mathbf{P}}_k\mathbf{C}_k^\top(\mathbf{C}_k\check{\mathbf{P}}_k\mathbf{C}_k^\top+\boldsymbol{\Sigma}_{\mathbf{v}})^{-1}$，即得 \cref{eq:kf-gain,eq:kf-correct}。两条路（MAP 配方与高斯条件化）殊途同归——这\emph{不是巧合}，因 LG 后验精确高斯、均值 $=$ 众数，故 MAP（求众数）与贝叶斯（求均值）同解。
\end{proof}

\paragraph{增益的物理意义：先验与观测按精度折中。}
由 \cref{eq:kf-correct}，后验 $=$ 预测 $+$ 增益 $\times$ 新息。两个极端：
\begin{itemize}
  \item 观测极准 $\boldsymbol{\Sigma}_{\mathbf{v}}\to\mathbf{0}$：$\mathbf{K}_k\to\mathbf{C}_k^{-1}$（方阵可逆时），$\hat{\mathbf{x}}_k$ 几乎完全采纳观测，$\hat{\mathbf{P}}_k\to\mathbf{0}$；
  \item 观测极差 $\boldsymbol{\Sigma}_{\mathbf{v}}\to\infty$ 或先验极准 $\check{\mathbf{P}}_k\to\mathbf{0}$：$\mathbf{K}_k\to\mathbf{0}$，$\hat{\mathbf{x}}_k\to\check{\mathbf{x}}_k$，忽略新息。
\end{itemize}
增益的角色是“按可信度给新息配权重”——与上一章加权最小二乘里 $\boldsymbol{\Sigma}^{-1}$ 给残差配权重是同一思想\cite{barfoot2024state}。

\begin{insight}
\textbf{协方差/信息两形式的对偶取舍}：校正步\emph{信息形式}（\cref{eq:kf-info-P,eq:kf-info-x}）只需\emph{相加}——多传感器逐个叠加，无需求 $M\times M$ 逆；但预测步\emph{协方差形式} $\mathbf{A}\hat{\mathbf{P}}\mathbf{A}^\top+\boldsymbol{\Sigma}_{\mathbf{w}}$ 更省（信息形式预测要求两次逆）。故两形式在预测/校正的代价\emph{互补}：信息滤波（IF）适合“多观测、少预测”，协方差滤波适合“多预测、少观测”。这正是 \cref{ch:slam_est} 中“独立高斯相乘 $=$ 信息相加”的递归落地。
\end{insight}

上文把信息形式散在 \cref{eq:kf-info-P,eq:kf-info-x} 与对偶取舍里；这里把它\emph{收成一组完整的递归}，作为 KF 的对偶算法独立命名。记信息矩阵 $\boldsymbol{\Lambda}=\mathbf{P}^{-1}$ 与\emph{信息向量} $\boldsymbol{\eta}=\boldsymbol{\Lambda}\,\mathbf{x}$（注意它不是状态本身，要状态时须解 $\mathbf{x}=\boldsymbol{\Lambda}^{-1}\boldsymbol{\eta}$）。

\begin{proposition}[信息滤波（IF）递归]\label{prop:info-filter}
线性高斯系统 \cref{eq:lg-model} 的卡尔曼滤波在\emph{信息形式}下为\cite{thrun2005probabilistic,barfoot2024state}：
\begin{align}
  \text{预测：}\quad
  &\check{\boldsymbol{\Lambda}}_k=\big(\mathbf{A}_{k-1}\hat{\boldsymbol{\Lambda}}_{k-1}^{-1}\mathbf{A}_{k-1}^\top+\boldsymbol{\Sigma}_{\mathbf{w}}\big)^{-1},
  &&\check{\boldsymbol{\eta}}_k=\check{\boldsymbol{\Lambda}}_k\big(\mathbf{A}_{k-1}\hat{\boldsymbol{\Lambda}}_{k-1}^{-1}\hat{\boldsymbol{\eta}}_{k-1}+\mathbf{u}_k\big);
  \label{eq:if-predict}\\
  \text{校正：}\quad
  &\hat{\boldsymbol{\Lambda}}_k=\check{\boldsymbol{\Lambda}}_k+\mathbf{C}_k^\top\boldsymbol{\Sigma}_{\mathbf{v}}^{-1}\mathbf{C}_k,
  &&\hat{\boldsymbol{\eta}}_k=\check{\boldsymbol{\eta}}_k+\mathbf{C}_k^\top\boldsymbol{\Sigma}_{\mathbf{v}}^{-1}\mathbf{z}_k.
  \label{eq:if-correct}
\end{align}
它与 KF（\cref{thm:kf}）给出\emph{完全相同}的估计，只是\textbf{代价结构对偶}：校正步只是\emph{两项相加}（多个独立观测可逐个累加 $\mathbf{C}^\top\boldsymbol{\Sigma}_{\mathbf{v}}^{-1}\mathbf{C}$、$\mathbf{C}^\top\boldsymbol{\Sigma}_{\mathbf{v}}^{-1}\mathbf{z}$，无需求 $M\times M$ 新息逆），而预测步要\emph{两次求逆}（$\hat{\boldsymbol{\Lambda}}_{k-1}^{-1}$ 与外层再求逆）。
\end{proposition}
\begin{proof}[要点]
校正即 \cref{eq:kf-info-P,eq:kf-info-x} 两式（信息相加），改记 $\boldsymbol{\Lambda},\boldsymbol{\eta}$ 即得。预测把协方差形式 $\check{\mathbf{P}}_k=\mathbf{A}\hat{\mathbf{P}}_{k-1}\mathbf{A}^\top+\boldsymbol{\Sigma}_{\mathbf{w}}$ 两边求逆即 \cref{eq:if-predict} 的 $\check{\boldsymbol{\Lambda}}_k$，均值式由 $\check{\boldsymbol{\eta}}_k=\check{\boldsymbol{\Lambda}}_k\check{\mathbf{x}}_k$、$\check{\mathbf{x}}_k=\mathbf{A}\hat{\mathbf{x}}_{k-1}+\mathbf{u}_k$、$\hat{\mathbf{x}}_{k-1}=\hat{\boldsymbol{\Lambda}}_{k-1}^{-1}\hat{\boldsymbol{\eta}}_{k-1}$ 代入得到。\hfill$\square$
\end{proof}

\begin{insight}
\textbf{IF 的两个独到之处与 SEIF}：(1) 信息形式能表达“\emph{全局无知}”——令 $\boldsymbol{\Lambda}\to\mathbf{0}$ 即“对状态一无所知”的弥散先验（协方差形式要写 $\mathbf{P}\to\infty$ 无法数值表示），故 IF 适合无先验启动与多机融合（各方信息直接相加）。(2) 在 SLAM 里，信息矩阵 $\boldsymbol{\Lambda}$ 天然\emph{近似稀疏}（远处路标间几乎无直接信息耦合），\emph{稀疏扩展信息滤波}（SEIF\cite{thrun2005probabilistic}）正是\textbf{主动把 $\boldsymbol{\Lambda}$ 的弱元素置零}（稀疏化）以换取近常数时间更新，是 EKF-SLAM 因协方差稠密而 $O(n^2)$ 的一种早期破解。但稀疏化是\emph{近似}（丢弱耦合），其一致性争议也促成了后来向因子图\emph{精确}稀疏（\cref{sec:filter-vs-opt}）的转向。
\end{insight}

\begin{theorem}[KF $\equiv$ 批量 MAP 的递归]\label{thm:kf-map}
在线性高斯下，卡尔曼滤波在每个时刻给出的 $\hat{\mathbf{x}}_k,\hat{\mathbf{P}}_k$，与\emph{用至 $k$ 时刻全部数据}的批量 MAP（\cref{ch:slam_est}）的边缘后验\textbf{完全相同}，并非近似。
\end{theorem}
\begin{proof}[证明思路]
把“至 $k-1$ 的全部历史数据”用一个等价先验 $\{\hat{\mathbf{x}}_{k-1},\hat{\mathbf{P}}_{k-1}\}$ 代替（马尔可夫性保证 LG 下这一替代精确），构造仅含 $\{\hat{\mathbf{x}}_{k-1},\mathbf{u}_k,\mathbf{z}_k\}$ 的小批量最小二乘，解其正规方程并用舒尔补消去 $\mathbf{x}_{k-1}$，再用矩阵求逆引理（SMW）化简，恰好得到 \cref{eq:kf-predict,eq:kf-gain,eq:kf-correct}。完整代数见 \cref{sec:eskf-appendix}。根因是 LG 后验\emph{精确高斯}，其均值 $=$ 众数，故 MAP（求众数）与贝叶斯（求均值）同解。
\end{proof}

\begin{insight}
\cref{thm:kf-map} 是本章与上一章的缝合点：\textbf{滤波不是另一套理论，而是批量 MAP 的因果递归实现}。它也预告了非线性的分水岭——一旦模型非线性，后验不再高斯、均值 $\neq$ 众数，KF 的各种推广（EKF 取众数附近、迭代 SPKF 趋均值）开始分化（\cref{sec:ekf,sec:nonlin-pdf}）。
\end{insight}

\paragraph{KF 的性质：无偏、一致、最优（BLUE 与 CRLB）。}
KF 估计\emph{无偏}（$\mathbb{E}[\hat{\mathbf{x}}_k-\mathbf{x}_k]=\mathbf{0}$）且\emph{一致}（其报告的 $\hat{\mathbf{P}}_k$ 确\emph{等于}真实误差协方差 $\mathbb{E}[\hat{\mathbf{e}}_k\hat{\mathbf{e}}_k^\top]$，下节把“一致”量化为 NEES）——二者均可由误差动力学用\emph{数学归纳法}两步证得（\cref{der:kf-unbiased-consistent}，它正是带偏置的 \cref{thm:bias-overconfident} 令偏置 $=\mathbf{0}$ 的洁净特例）。正因报告协方差\emph{恰好}等于真实误差协方差，KF 才称\emph{最优}：即便去掉高斯假设，只要噪声零均值、不相关、方差已知，KF 仍是\textbf{最优线性无偏估计}（BLUE，高斯–马尔可夫定理），其协方差达到克拉默–拉奥下界（CRLB）\cite{barfoot2024state}。这里的期望是\emph{对随机变量全体实现取（集合）平均}、非单条轨迹的时间平均：单次实现的误差并不会真的为零或随时间衰减到零。

\begin{proposition}[协方差等比缩放不变性]\label{prop:kf-cov-scaling}
给三个协方差 $\boldsymbol{\Sigma}_{\mathbf{w}}$、$\boldsymbol{\Sigma}_{\mathbf{v}}$、$\hat{\mathbf{P}}_0$ 同乘一个正系数 $a$，则 KF 的\emph{点估计} $\{\hat{\mathbf{x}}_k\}$ 与\emph{增益} $\{\mathbf{K}_k\}$ 完全不变，而所有协方差 $\{\check{\mathbf{P}}_k,\hat{\mathbf{P}}_k\}$ 一律放大 $a$ 倍\cite{unitree2023quadruped}。
\end{proposition}
\begin{proof}
归纳。设 $\hat{\mathbf{P}}_{k-1}\to a\hat{\mathbf{P}}_{k-1}$ 而 $\hat{\mathbf{x}}_{k-1}$ 不变。预测 $\check{\mathbf{P}}_k'=\mathbf{A}_{k-1}(a\hat{\mathbf{P}}_{k-1})\mathbf{A}_{k-1}^\top+a\boldsymbol{\Sigma}_{\mathbf{w}}=a\check{\mathbf{P}}_k$，而 $\check{\mathbf{x}}_k$ 不含协方差故不变；增益的系数 $a$ 上下抵消
\[
  \mathbf{K}_k'=(a\check{\mathbf{P}}_k)\mathbf{C}_k^\top\big(\mathbf{C}_k(a\check{\mathbf{P}}_k)\mathbf{C}_k^\top+a\boldsymbol{\Sigma}_{\mathbf{v}}\big)^{-1}
  =a\check{\mathbf{P}}_k\mathbf{C}_k^\top\cdot\tfrac1a\big(\mathbf{C}_k\check{\mathbf{P}}_k\mathbf{C}_k^\top+\boldsymbol{\Sigma}_{\mathbf{v}}\big)^{-1}=\mathbf{K}_k,
\]
故 $\hat{\mathbf{x}}_k=\check{\mathbf{x}}_k+\mathbf{K}_k(\mathbf{z}_k-\mathbf{C}_k\check{\mathbf{x}}_k)$ 不变，而 $\hat{\mathbf{P}}_k'=(\mathbf{I}-\mathbf{K}_k\mathbf{C}_k)(a\check{\mathbf{P}}_k)=a\hat{\mathbf{P}}_k$，缩放向下一步传递；$k=0$ 由 $\hat{\mathbf{P}}_0\to a\hat{\mathbf{P}}_0$ 成立，归纳完成。\hfill$\square$
\end{proof}
\noindent 两条推论：(i) \textbf{可把全部协方差整体缩放到便于数值计算的量级而不动估计}；(ii) \emph{只有 $\boldsymbol{\Sigma}_{\mathbf{w}}$ 与 $\boldsymbol{\Sigma}_{\mathbf{v}}$ 的\textbf{比值}（而非绝对大小）影响点估计}——调参真正要调的是“信任运动还是信任观测”的比值，把 $\boldsymbol{\Sigma}_{\mathbf{w}}$ 设得过大等价于变相缩小 $\boldsymbol{\Sigma}_{\mathbf{v}}$、令滤波过度信任观测（表现为状态高频抖动）。但绝对大小仍影响 NEES/NIS 判据（$\hat{\mathbf{P}}_k$ 随 $a$ 缩放使 $\epsilon_{\text{nees}}$ 随 $1/a$ 变，\cref{def:nees-nis}），故“缩放不改点估计”\emph{不等于}“协方差可乱设”。

\paragraph{稳态卡尔曼滤波与离散代数 Riccati 方程（DARE）。}
\cref{eq:kf-predict,eq:kf-gain,eq:kf-correct} 的\emph{协方差}递推可\emph{独立于均值}先跑（它不依赖观测数据）。对\textbf{时不变}系统（$\mathbf{A}_k\equiv\mathbf{A}$、$\mathbf{C}_k\equiv\mathbf{C}$、$\boldsymbol{\Sigma}_{\mathbf{w}},\boldsymbol{\Sigma}_{\mathbf{v}}$ 常值），一个自然的问题是：$\check{\mathbf{P}}_k$ 会不会收敛到一个常值 $\mathbf{P}$？若会，便可预先把协方差迭代到收敛、此后只用\emph{常值增益} $\mathbf{K}$ 传播均值，省去每步的协方差更新与求逆——这就是\emph{稳态卡尔曼滤波}，工程上极常用。把预测与校正协方差式合并（代入 $\hat{\mathbf{P}}=(\mathbf{I}-\mathbf{K}\mathbf{C})\check{\mathbf{P}}$、$\mathbf{K}=\check{\mathbf{P}}\mathbf{C}^\top(\mathbf{C}\check{\mathbf{P}}\mathbf{C}^\top+\boldsymbol{\Sigma}_{\mathbf{v}})^{-1}$），稳态 $\mathbf{P}$ 须满足\textbf{离散代数 Riccati 方程}（DARE）\cite{barfoot2024state}：
\begin{equation}
  \mathbf{P}=\mathbf{A}(\mathbf{I}-\mathbf{K}\mathbf{C})\mathbf{P}(\mathbf{I}-\mathbf{K}\mathbf{C})^\top\mathbf{A}^\top+\mathbf{A}\mathbf{K}\boldsymbol{\Sigma}_{\mathbf{v}}\mathbf{K}^\top\mathbf{A}^\top+\boldsymbol{\Sigma}_{\mathbf{w}},\qquad
  \mathbf{K}=\mathbf{P}\mathbf{C}^\top(\mathbf{C}\mathbf{P}\mathbf{C}^\top+\boldsymbol{\Sigma}_{\mathbf{v}})^{-1}.
  \label{eq:dare}
\end{equation}
（这与 \cref{ch:control} LQR 的 Riccati 方程是\emph{估计--控制对偶}的两面。）DARE 存在\emph{唯一半正定}解 $\mathbf{P}$ 的充要条件是下列四条\cite{barfoot2024state}：
\begin{enumerate}
  \item $\boldsymbol{\Sigma}_{\mathbf{v}}>\mathbf{0}$（观测噪声正定，批量 LG 已默认）；
  \item $\boldsymbol{\Sigma}_{\mathbf{w}}\ge\mathbf{0}$（过程噪声半正定）；
  \item $(\mathbf{A},\mathbf{V})$ \emph{可镇定}（stabilizable），其中 $\mathbf{V}\mathbf{V}^\top=\boldsymbol{\Sigma}_{\mathbf{w}}$——当 $\boldsymbol{\Sigma}_{\mathbf{w}}>\mathbf{0}$ 时此条冗余；
  \item $(\mathbf{A},\mathbf{C})$ \emph{可检测}（detectable）——即“可观，或不可观模态都已稳定”，与 \cref{ch:slam_est} 批量 LG 的可观测性秩条件同源。
\end{enumerate}
此时滤波\emph{误差动力学稳定}：$\mathbb{E}[\check{\mathbf{e}}_k]=\mathbf{A}(\mathbf{I}-\mathbf{K}\mathbf{C})\,\mathbb{E}[\check{\mathbf{e}}_{k-1}]$，其转移矩阵 $\mathbf{A}(\mathbf{I}-\mathbf{K}\mathbf{C})$ 谱半径 $<1$（\cref{eq:dare} 对任一特征对取二次型可证 $|\lambda|^2<1$，可检测性保证右端的可观测性 Gramian 非退化）\cite{barfoot2024state}。\emph{可镇定/可检测}是比\emph{可控/可观}更弱的条件——只要求“坏模态可控/可观”，已稳定的模态可放任，这与 \cref{ch:control} 的能控性/能观性梯度一脉相承。

\paragraph{RTS 平滑器：前向 KF $+$ 后向修正。}
若允许离线，用未来数据可得更优的\emph{平滑}估计（即 \cref{eq:eskf-fwd-bwd} 的后向因子）。Rauch–Tung–Striebel（RTS）平滑器 $=$ 前向跑一遍 KF（记滤波量 $\hat{\mathbf{x}}_{k,f},\hat{\mathbf{P}}_{k,f}$），再后向递推：
\begin{align}
  \mathbf{G}_{k-1}&=\hat{\mathbf{P}}_{k-1,f}\mathbf{A}_{k-1}^\top\check{\mathbf{P}}_{k,f}^{-1}, \label{eq:rts}\\
  \hat{\mathbf{x}}_{k-1}&=\hat{\mathbf{x}}_{k-1,f}+\mathbf{G}_{k-1}\big(\hat{\mathbf{x}}_k-\check{\mathbf{x}}_{k,f}\big),\qquad
  \hat{\mathbf{P}}_{k-1}=\hat{\mathbf{P}}_{k-1,f}+\mathbf{G}_{k-1}\big(\hat{\mathbf{P}}_k-\check{\mathbf{P}}_{k,f}\big)\mathbf{G}_{k-1}^\top, \label{eq:rts-cov}
\end{align}
$\mathbf{G}_{k-1}$ 为\emph{平滑增益}（协方差后向式的完整化简见 \cref{sec:eskf-appendix}）。RTS 用整段数据精确求解原批量问题、无近似，复杂度仍 $O(K)$——再次印证“滤波/平滑 $=$ 批量 MAP 的稀疏递归”\cite{barfoot2024state}。平滑协方差恒不大于对应滤波协方差（未来数据只会增加信息）。

\paragraph{（选读）连续时间极限：Kalman--Bucy 滤波}
以上全是\emph{离散}时刻递推。若运动与观测都\emph{连续}进行（传感器以无穷高频采样的理想极限），KF 的连续时间对应即 Kalman--Bucy 滤波（Kalman \& Bucy 1961），是估计理论的经典地标\cite{barfoot2024state}。取线性时变（LTV）随机微分方程模型与连续观测
\begin{equation}
  \dot{\mathbf{x}}(t)=\mathbf{A}(t)\mathbf{x}(t)+\mathbf{B}(t)\mathbf{u}(t)+\mathbf{L}(t)\mathbf{w}(t),\qquad
  \mathbf{z}(t)=\mathbf{C}(t)\mathbf{x}(t)+\mathbf{v}(t),
  \label{eq:kb-model}
\end{equation}
$\mathbf{w}(t),\mathbf{v}(t)$ 为零均值高斯过程，功率谱密度（连续时间“协方差”）分别为 $\mathbf{Q}(t),\mathbf{R}(t)$（对称正定）。则估计均值 $\hat{\mathbf{x}}(t)$ 与协方差 $\hat{\mathbf{P}}(t)$ 满足\emph{两条耦合常微分方程}\cite{barfoot2024state}：
\begin{align}
  \dot{\hat{\mathbf{x}}}(t)&=\mathbf{A}(t)\hat{\mathbf{x}}(t)+\mathbf{B}(t)\mathbf{u}(t)+\mathbf{K}(t)\big(\mathbf{z}(t)-\mathbf{C}(t)\hat{\mathbf{x}}(t)\big),
  \label{eq:kb-mean}\\
  \dot{\hat{\mathbf{P}}}(t)&=\mathbf{A}(t)\hat{\mathbf{P}}(t)+\hat{\mathbf{P}}(t)\mathbf{A}(t)^\top+\mathbf{L}(t)\mathbf{Q}(t)\mathbf{L}(t)^\top-\mathbf{K}(t)\mathbf{R}(t)\mathbf{K}(t)^\top,
  \label{eq:kb-cov}
\end{align}
\emph{连续时间卡尔曼增益} $\mathbf{K}(t)=\hat{\mathbf{P}}(t)\mathbf{C}(t)^\top\mathbf{R}(t)^{-1}$，初值 $\hat{\mathbf{x}}(0),\hat{\mathbf{P}}(0)$，须数值或解析积分。\textbf{结构对照}：\cref{eq:kb-mean} 仍是“开环传播 $+$ 增益乘新息”，与离散 \cref{eq:kf-correct} 同构；\cref{eq:kb-cov} 是一条\emph{微分} Riccati 方程，正是 \cref{eq:dare} 在 $\Delta t\to0$ 的连续极限（其稳态 $\dot{\hat{\mathbf{P}}}=\mathbf{0}$ 给\emph{连续}代数 Riccati 方程，与 \cref{ch:control} LQR 的连续 Riccati 同形）。多数传感器实为离散采样，故工程上以离散 KF 为主；Kalman--Bucy 的价值在于它把“预测使不确定度增长（$+\mathbf{L}\mathbf{Q}\mathbf{L}^\top$）、观测使其减小（$-\mathbf{K}\mathbf{R}\mathbf{K}^\top$）”这一对作用写成同一条 ODE 里的两项，直观至极。

\begin{pitfall}[数值陷阱：协方差更新用 $(\mathbf{I}-\mathbf{K}\mathbf{C})\check{\mathbf{P}}$ 失对称]
\cref{eq:kf-correct} 的 $\hat{\mathbf{P}}_k=(\mathbf{I}-\mathbf{K}_k\mathbf{C}_k)\check{\mathbf{P}}_k$ 在浮点下\emph{不保证对称正定}（减法相消、$\mathbf{K}$ 略有误差就破坏对称），长跑后 $\hat{\mathbf{P}}$ 可能出现负特征值致滤波发散。对策（Solà\cite{sola2017quaternion} 脚注 26）：用对称形式 $\hat{\mathbf{P}}_k=\check{\mathbf{P}}_k-\mathbf{K}_k\mathbf{S}_k\mathbf{K}_k^\top$，或对称正定的 \textbf{Joseph 形式}
\[
  \hat{\mathbf{P}}_k=(\mathbf{I}-\mathbf{K}_k\mathbf{C}_k)\check{\mathbf{P}}_k(\mathbf{I}-\mathbf{K}_k\mathbf{C}_k)^\top+\mathbf{K}_k\boldsymbol{\Sigma}_{\mathbf{v}}\mathbf{K}_k^\top,
\]
后者对任意（哪怕次优的）$\mathbf{K}$ 都给出对称正定结果，工程实现强烈推荐。\textbf{更彻底的数值稳健化是“平方根滤波”}：直接传播协方差的\emph{平方根因子}（Cholesky 因子或其变体）而非 $\mathbf{P}$ 本身。理由与优化侧 \cref{sec:nlopt-solve} 的 $\sqrt{}$SAM 完全同源——直接对信息矩阵 $\boldsymbol{\Lambda}=\mathbf{P}^{-1}$ 操作会把条件数\emph{平方}（$\kappa(\boldsymbol{\Lambda})=\kappa(\mathbf{A})^2$，$\mathbf{A}$ 为白化雅可比），而对其平方根因子操作把条件数\emph{开方}回 $\kappa(\mathbf{A})$，数值精度大为改善。滤波侧的两个经典实现是\emph{平方根信息滤波}（SRIF，传播信息阵的平方根，与 $\sqrt{}$SAM 是同一对象的递归版）与 \textbf{UD 分解}（Bierman，把 $\mathbf{P}=\mathbf{U}\mathbf{D}\mathbf{U}^\top$ 分解为单位上三角 $\mathbf{U}$ 与对角 $\mathbf{D}$，更新中\emph{避免显式开方}、保正定）。Maybeck 甚至主张“平方根滤波应总是优先于标准 KF 递推”。本书把平方根方法的完整数值论证集中放在 \cref{sec:nlopt-solve}（优化侧），此处只点明滤波侧同理、实现见 SRIF/UD。
\end{pitfall}

\begin{example}[一维定位的两步 KF 手算]\label{ex:kf-1d}
一维常速模型 $x_k=x_{k-1}+u_k+w_k$、$z_k=x_k+v_k$，$\boldsymbol{\Sigma}_{\mathbf{w}}=Q$、$\boldsymbol{\Sigma}_{\mathbf{v}}=R$（标量）。由 \cref{thm:kf}：预测 $\check{x}_k=\hat{x}_{k-1}+u_k$、$\check{P}_k=\hat{P}_{k-1}+Q$；增益 $K_k=\check{P}_k/(\check{P}_k+R)\in(0,1)$；校正 $\hat{x}_k=\check{x}_k+K_k(z_k-\check{x}_k)=(1-K_k)\check{x}_k+K_k z_k$（预测与观测的凸组合）、$\hat{P}_k=(1-K_k)\check{P}_k=\dfrac{R}{\check{P}_k+R}\check{P}_k$。可见 $\hat{P}_k<\check{P}_k$（观测总在收缩不确定度），且 $\hat{P}_k<R$、$\hat{P}_k<\check{P}_k$——后验比先验和单次观测都更确定。这正是 \cref{fig:eskf-predict-correct} 椭圆“先胀后缩”的标量版。（对应 Barfoot 习题 4.7。）
\end{example}

\begin{practice}
\begin{enumerate}
  \item 由 \cref{eq:kf-info-P,eq:kf-info-x} 与矩阵求逆引理推出 \cref{eq:kf-gain} 的两种等价增益表达 $\mathbf{K}_k=\hat{\mathbf{P}}_k\mathbf{C}_k^\top\boldsymbol{\Sigma}_{\mathbf{v}}^{-1}=\check{\mathbf{P}}_k\mathbf{C}_k^\top(\mathbf{C}_k\check{\mathbf{P}}_k\mathbf{C}_k^\top+\boldsymbol{\Sigma}_{\mathbf{v}})^{-1}$。答案要点：对 $\hat{\mathbf{P}}_k^{-1}=\check{\mathbf{P}}_k^{-1}+\mathbf{C}^\top\boldsymbol{\Sigma}_{\mathbf{v}}^{-1}\mathbf{C}$ 用 SMW 求 $\hat{\mathbf{P}}_k$，再代入。答错回看 \cref{eq:kf-info-P} 与 \cref{ch:slam_est} 的 SMW。
  \item 接 \cref{ex:kf-1d}：取 $\hat{x}_0=0,\hat{P}_0=1,Q=1,R=1/2$，手算 $k=1$（$u_1=1,z_1=2$）的 $\check{P}_1,K_1,\hat{P}_1,\hat{x}_1$，并观察 $\hat{P}_1$ 如何先因预测变大、再因观测变小。
  \item （概念）解释为什么 LG 下 KF 与批量 MAP 同解，而非线性时不再如此。答案要点：LG 后验精确高斯$\Rightarrow$均值=众数$\Rightarrow$MAP=贝叶斯；非线性使后验偏斜，均值$\ne$众数。答错回看 \cref{thm:kf-map} 与 \cref{sec:ekf} 引言。
\end{enumerate}
\end{practice}

% ════════════════════════════════════════════════════════════════
\section{扩展卡尔曼滤波：线性化非线性系统}
\label{sec:ekf}

真实 SLAM 的 $\mathbf{f},\mathbf{h}$ 非线性（相机投影、李群位姿）。\textbf{核心反面}：一个高斯过非线性\emph{不再}高斯——下面这个一维例子把它讲透，也顺带说明“先验、测量都是高斯，后验却偏斜”的事实。

\paragraph{动机例：立体相机测距，后验为何偏斜。}
Barfoot\cite{barfoot2024state} §4.1 的一维立体相机模型：路标深度 $x$、视差 $y$，$y=\dfrac{fb}{x}+n$，$n\sim\mathcal{N}(0,R)$（$f$ 焦距、$b$ 基线）。先验高斯 $p(x)=\mathcal{N}(\check{x},\check{P})$。贝叶斯后验
\begin{equation}
  p(x\mid y)=\frac{p(y\mid x)\,p(x)}{\int p(y\mid x)\,p(x)\,\mathrm{d}x},\qquad
  p(y\mid x)=\mathcal{N}\!\Big(\frac{fb}{x},R\Big).
  \label{eq:stereo-bayes}
\end{equation}
即使先验与似然都是高斯，因 $1/x$ 非线性，\textbf{后验是非对称（偏斜）的}（被非线性观测拉向一侧），但仍单峰。这立刻引出两个不同的“点估计”目标：

\begin{itemize}
  \item \textbf{完全贝叶斯}：报告后验\emph{均值} $\bar{x}=\mathbb{E}[x\mid y]$（需算那个一般不可解的分母积分）；
  \item \textbf{最大后验（MAP）}：只找后验\emph{众数}（峰），等价最小化负对数似然
  \begin{equation}
    \hat{x}_{\text{map}}=\arg\min_x\Big[\underbrace{\tfrac{1}{2R}\big(y-\tfrac{fb}{x}\big)^2}_{-\ln p(y\mid x)}+\underbrace{\tfrac{1}{2\check{P}}(\check{x}-x)^2}_{-\ln p(x)}\Big].
    \label{eq:stereo-map}
  \end{equation}
\end{itemize}
线性观测时众数 $=$ 均值、二者一致；非线性时\textbf{众数 $\ne$ 均值，故 MAP 是有偏的}。要量化“有多偏、有多准”须\emph{两个}度量、缺一不可\cite{barfoot2024state}。其一是\emph{均值误差}
\begin{equation}
  \hat{e}_{\text{mean}}(\hat{x})=\mathbb{E}_{XN}[\hat{x}-x_{\text{true}}]=\mathbb{E}_{XN}[\hat{x}]-\check{x}
  \label{eq:e-sq-def}
\end{equation}
（下标 $XN$ 表示\emph{同时}对先验抽真值 $x_{\text{true}}$、对噪声抽 $n$ 取期望；因 $x_{\text{true}}$ 独立于 $n$ 故 $\mathbb{E}_{XN}[x_{\text{true}}]=\mathbb{E}_X[x_{\text{true}}]=\check{x}$），它只测\emph{偏差}。但 $\hat{e}_{\text{mean}}$ \emph{单独不够}——平凡地令 $\hat{x}=\check{x}$（直接报先验均值）就得 $\hat{e}_{\text{mean}}(\check{x})=0$ 却毫无信息量。故须补第二度量\emph{均方误差} $\hat{e}_{\text{sq}}(\hat{x})=\mathbb{E}_{XN}[(\hat{x}-x_{\text{true}})^2]$，它\emph{合并了偏差与方差}（机器学习里的偏差–方差权衡，Bishop），把“估计塌成先验”这种零偏差但高方差的退化解暴露出来；\textbf{两个度量都好}才是合格的估计器。取参数 $\check{x}=20$m、$\check{P}=9$m$^2$、$f=400$px、$b=0.1$m、$R=0.09$px$^2$，$10^6$ 次试验得 $\hat{e}_{\text{mean}}(\hat{x}_{\text{map}})\approx-33.0$cm（小但非零的系统偏差）、$\hat{e}_{\text{sq}}(\hat{x}_{\text{map}})\approx4.41$m$^2$\cite{barfoot2024state}（真值即抽自估计器所用先验仍有偏；实务中先验未知、偏差常更大）。\textbf{方法论提醒}：评判一个非线性滤波器，先问它\emph{捕获完整后验的哪个方面}（均值？众数？），目标不同时谈“谁更准”不公平。

\begin{figure}[ht]
\centering
\begin{tikzpicture}[scale=1.0]
  \draw[->,thick,gray!60!black] (-0.2,0)--(9.0,0) node[right,font=\scriptsize]{$\hat{x}_{\text{map}}$};
  \draw[->,thick,gray!60!black] (0,-0.1)--(0,3.5) node[above,font=\scriptsize]{频数};
  \foreach \x/\h in {0.5/0.2, 1.1/0.5, 1.7/1.0, 2.3/1.7, 2.9/2.6, 3.5/3.1, 4.1/2.7, 4.7/1.8, 5.3/1.1, 5.9/0.6, 6.5/0.35, 7.1/0.18, 7.7/0.08} {
    \fill[main!22, draw=main!55!black] (\x,0) rectangle (\x+0.55,\h);
  }
  \draw[main!70!black, very thick] (3.55,0)--(3.55,3.0);
  \node[main!70!black, font=\scriptsize, left] at (3.55,3.25) {$\mathbb{E}[\hat{x}_{\text{map}}]$};
  \draw[second, thick, dashed] (4.3,0)--(4.3,3.0);
  \node[second, font=\scriptsize, right] at (4.3,3.25) {$\check{x}$};
  \draw[<->, thick, gray!50!black] (3.55,0.55)--(4.3,0.55);
  \node[font=\scriptsize, fill=white, inner sep=1pt] at (3.93,0.55) {$\hat{e}_{\text{mean}}\!\approx\!-33$cm};
\end{tikzpicture}
\caption{立体相机 MAP 估计值在 $10^6$ 次试验下的直方图（每次从先验抽新 $x_{\text{true}}$、从观测模型抽新 $y$）：虚线为先验均值 $\check{x}$、实线为估计期望 $\mathbb{E}[\hat{x}_{\text{map}}]$，两者之差即偏差 $\hat{e}_{\text{mean}}\approx-33$cm（仿 Barfoot Fig.~4.4 重绘）。趋均值的 ISPKF 在同设置下偏差仅 $\approx-3.84$cm（约 $1/10$，\cref{algo:ispkf} 旁），直方图整体更贴近 $\check{x}$。}
\label{fig:map-bias-hist}
\end{figure}

\begin{insight}
\textbf{本质洞察}：从 KF 到所有非线性滤波，分水岭是“一个高斯过非线性后非高斯，均值 $\ne$ 众数”。于是出现两条路：\emph{找众数}（MAP / 迭代线性化 / IEKF）与\emph{找均值}（蒙特卡洛 / sigma 点 / 迭代 SPKF）。EKF 是\emph{不迭代}地线性化一次，离这两者都不确定——这就是它“与完整后验无明确关系”的根源（\cref{tab:filter-taxo}）。
\end{insight}

\paragraph{EKF 的三步构造与逐式推导。}
扩展卡尔曼滤波（EKF）= (i) 约束置信度为高斯 + (ii) 假设噪声高斯 + (iii) \emph{在当前最佳估计处线性化}模型，从而闭式完成贝叶斯滤波的积分\cite{barfoot2024state,gaoxiang2019slam14}。把运动/观测在工作点一阶泰勒展开（噪声可在非线性内部，故其雅可比折进等效协方差，记带撇号者 $\boldsymbol{\Sigma}_{\mathbf{w}}'$、$\boldsymbol{\Sigma}_{\mathbf{v}}'$；加性噪声时撇号即原协方差）：
\begin{align}
  \mathbf{x}_k&\approx\check{\mathbf{x}}_k+\mathbf{F}(\mathbf{x}_{k-1}-\hat{\mathbf{x}}_{k-1})+\mathbf{w}_k',\quad
  \check{\mathbf{x}}_k=\mathbf{f}(\hat{\mathbf{x}}_{k-1},\mathbf{u}_k,\mathbf{0}),\quad
  \mathbf{F}=\left.\frac{\partial\mathbf{f}}{\partial\mathbf{x}_{k-1}}\right|_{\hat{\mathbf{x}}_{k-1},\mathbf{u}_k,\mathbf{0}},\label{eq:ekf-lin-f}\\
  \mathbf{z}_k&\approx\check{\mathbf{z}}_k+\mathbf{H}(\mathbf{x}_k-\check{\mathbf{x}}_k)+\mathbf{v}_k',\quad
  \check{\mathbf{z}}_k=\mathbf{h}(\check{\mathbf{x}}_k,\mathbf{0}),\quad
  \mathbf{H}=\left.\frac{\partial\mathbf{h}}{\partial\mathbf{x}_k}\right|_{\check{\mathbf{x}}_k,\mathbf{0}}.\label{eq:ekf-lin-h}
\end{align}
\cref{eq:ekf-lin-f} 给出 $\mathbb{E}[\mathbf{x}_k]\approx\check{\mathbf{x}}_k+\mathbf{F}(\mathbf{x}_{k-1}-\hat{\mathbf{x}}_{k-1})$、$\mathrm{Cov}\approx\boldsymbol{\Sigma}_{\mathbf{w}}'$，于是 $p(\mathbf{x}_k\mid\mathbf{x}_{k-1},\mathbf{u}_k)\approx\mathcal{N}(\check{\mathbf{x}}_k+\mathbf{F}(\mathbf{x}_{k-1}-\hat{\mathbf{x}}_{k-1}),\boldsymbol{\Sigma}_{\mathbf{w}}')$；同理 $p(\mathbf{z}_k\mid\mathbf{x}_k)\approx\mathcal{N}(\check{\mathbf{z}}_k+\mathbf{H}(\mathbf{x}_k-\check{\mathbf{x}}_k),\boldsymbol{\Sigma}_{\mathbf{v}}')$。把这两个高斯代入贝叶斯滤波 \cref{eq:bayes-filter}：积分项（用“高斯过线性映射仍高斯”）给出预测先验 $\mathcal{N}(\check{\mathbf{x}}_k,\mathbf{F}\hat{\mathbf{P}}_{k-1}\mathbf{F}^\top+\boldsymbol{\Sigma}_{\mathbf{w}}')$；再乘观测似然（两高斯归一化乘积）即得后验。整理出：

\begin{proposition}[EKF 递归]\label{prop:ekf}
\begin{align}
  &\check{\mathbf{x}}_k=\mathbf{f}(\hat{\mathbf{x}}_{k-1},\mathbf{u}_k), &
  &\check{\mathbf{P}}_k=\mathbf{F}\hat{\mathbf{P}}_{k-1}\mathbf{F}^\top+\boldsymbol{\Sigma}_{\mathbf{w}}', \label{eq:ekf-predict}\\
  &\mathbf{K}_k=\check{\mathbf{P}}_k\mathbf{H}^\top\big(\mathbf{H}\check{\mathbf{P}}_k\mathbf{H}^\top+\boldsymbol{\Sigma}_{\mathbf{v}}'\big)^{-1},\\
  &\hat{\mathbf{x}}_k=\check{\mathbf{x}}_k+\mathbf{K}_k\big(\mathbf{z}_k-\mathbf{h}(\check{\mathbf{x}}_k)\big), &
  &\hat{\mathbf{P}}_k=(\mathbf{I}-\mathbf{K}_k\mathbf{H})\check{\mathbf{P}}_k. \label{eq:ekf-correct}
\end{align}
与 KF（\cref{thm:kf}）相比仅两处不同：\emph{均值}用非线性 $\mathbf{f},\mathbf{h}$ 传播、\emph{协方差}用雅可比 $\mathbf{F},\mathbf{H}$ 传播。即“EKF $=$ 把 KF 里的 $\mathbf{A},\mathbf{C}$ 换成线性化雅可比 $\mathbf{F},\mathbf{H}$，并用原非线性模型传均值”\cite{gaoxiang2019slam14}。
\end{proposition}

\begin{proposition}[广义高斯滤波：EKF 的更干净推导]\label{prop:ggf}
不必先线性化，只需\emph{一开始就假设}估计为高斯，写状态与最新观测的联合高斯
$\mathcal{N}\big([\boldsymbol{\mu}_{x,k};\boldsymbol{\mu}_{y,k}],[\boldsymbol{\Sigma}_{xx,k},\boldsymbol{\Sigma}_{xy,k};\boldsymbol{\Sigma}_{yx,k},\boldsymbol{\Sigma}_{yy,k}]\big)$，按舒尔补条件化即得通用校正步\cite{barfoot2024state}：
\begin{equation}
  \mathbf{K}_k=\boldsymbol{\Sigma}_{xy,k}\boldsymbol{\Sigma}_{yy,k}^{-1},\qquad
  \hat{\mathbf{P}}_k=\check{\mathbf{P}}_k-\mathbf{K}_k\boldsymbol{\Sigma}_{xy,k}^\top,\qquad
  \hat{\mathbf{x}}_k=\check{\mathbf{x}}_k+\mathbf{K}_k(\mathbf{z}_k-\boldsymbol{\mu}_{y,k}),
  \label{eq:ggf}
\end{equation}
其中 $\boldsymbol{\mu}_{x,k}=\check{\mathbf{x}}_k$、$\boldsymbol{\Sigma}_{xx,k}=\check{\mathbf{P}}_k$。除非模型线性，$\boldsymbol{\mu}_{y,k},\boldsymbol{\Sigma}_{yy,k},\boldsymbol{\Sigma}_{xy,k}$（观测的均值、协方差与互协方差）不能精确算——\emph{用什么方法近似算这三个矩}，就决定了你得到 EKF（线性化）还是 UKF（sigma 点）。这是统一所有高斯滤波器的“总接口”。
\end{proposition}

\begin{pitfall}[思维陷阱：EKF 的三宗罪——工作点偏差、单次线性化、不一致]
EKF 不保证对一般非线性系统够用，三个根本问题\cite{barfoot2024state,gaoxiang2019slam14}：
\begin{enumerate}
  \item \textbf{工作点是估计均值而非真值}：线性化点取在 $\hat{\mathbf{x}}$ 而非未知的 $\mathbf{x}_{\text{true}}$，这“看似微小”的差别在强非线性下可让 EKF \emph{剧烈发散}，常见后果是\emph{有偏 $+$ 不一致}（报告的协方差比真实误差乐观，\cref{sec:consistency} 量化）。
  \item \textbf{只线性化一次}：EKF 在 $\hat{\mathbf{x}}_{k-1}$ 处展开一次就把后验算出，默认该线性近似在整个后验范围有效；离工作点远时一阶泰勒不成立。
  \item \textbf{破坏马尔可夫性}：近似一旦引入，贝叶斯滤波赖以成立的马尔可夫因子化就不再严格——而\textbf{一旦把马尔可夫假设内建进估计器就无法去除}，这是“批量优化是更优起点”的理论依据（\cref{sec:filter-vs-opt}）。
\end{enumerate}
选型决策语（缩进 if-then 三问）：\emph{若}非线性弱 $\to$ EKF 够用；\emph{若}强非线性但可重线性化 $\to$ 迭代（IEKF）或改批量；\emph{若}雅可比难求或非光滑 $\to$ sigma 点（UKF）。一句话记忆：\textbf{“EKF 大致只是优化里的\emph{一次}迭代”}——优化每步都在新估计点重新线性化，故适用范围更广\cite{gaoxiang2019slam14}。
\end{pitfall}

\begin{practice}
\begin{enumerate}
  \item 对观测 $h(x)=\tfrac{fb}{x}$（立体相机视差–深度），写出 $H=\partial h/\partial x$ 并讨论 $x$ 较小时线性化误差为何变大。答案要点：$H=-fb/x^2$，$x$ 小则曲率大、一阶泰勒失效。
  \item 由 \cref{prop:ggf} 出发，把观测模型设为线性 $\mathbf{z}=\mathbf{C}\mathbf{x}+\mathbf{v}$，验证 $\boldsymbol{\mu}_y=\mathbf{C}\check{\mathbf{x}}$、$\boldsymbol{\Sigma}_{yy}=\mathbf{C}\check{\mathbf{P}}\mathbf{C}^\top+\boldsymbol{\Sigma}_{\mathbf{v}}$、$\boldsymbol{\Sigma}_{xy}=\check{\mathbf{P}}\mathbf{C}^\top$，从而 \cref{eq:ggf} 退回 \cref{thm:kf}。
  \item （概念）为什么说 MAP 对 \cref{eq:stereo-bayes} 是有偏的，而把先验也算进去仍无法消除该偏差？答案要点：偏差源于观测模型非线性使众数偏离均值，与先验是否参与无关。答错回看本节动机例。
\end{enumerate}
\end{practice}

% ════════════════════════════════════════════════════════════════
\section{把高斯过非线性的三种办法与滤波器谱系}
\label{sec:nonlin-pdf}

\cref{prop:ggf} 把所有高斯滤波器归结为“怎样近似 $\boldsymbol{\mu}_y,\boldsymbol{\Sigma}_{yy},\boldsymbol{\Sigma}_{xy}$”，也即“怎样让一个高斯过非线性 $\mathbf{y}=\mathbf{g}(\mathbf{x})$”。Barfoot\cite{barfoot2024state} 给出三种办法：\emph{蒙特卡洛}、\emph{线性化}（EKF 用）、\emph{sigma 点 / 无迹变换}（UKF 用）。本节脉络层为\textbf{次优}（理解谱系用），可在主线读完后回看。

\paragraph{办法一：蒙特卡洛（暴力，作基准）。}
从输入密度抽大量样本 $\mathbf{x}_i$，每个\emph{精确}过非线性 $\mathbf{y}_i=\mathbf{g}(\mathbf{x}_i)$，再由变换后样本算统计矩。大数定律保证样本数 $\to\infty$ 时收敛到真值。优点：适用任意 PDF、任意非线性（不要求可微甚至连续）、不需 $\mathbf{g}$ 的解析式（黑箱也行）、是“anytime”算法（样本换精度）。缺点：低效，高维尤甚。\textbf{关键事实}：输出密度均值 $\ne$ 输入均值过非线性。反例：$x\sim\mathrm{Uniform}[0,1]$、$y=x^2$，输入均值 $1/2$ 过非线性给 $1/4$，但真实输出均值 $\int_0^1 x^2\mathrm{d}x=1/3$。蒙特卡洛能逼近这个 $1/3$，而线性化不能。

\paragraph{办法二：线性化（EKF）。}
均值\emph{精确}过非线性、协方差过\emph{线性化版本}，工作点取均值。不准的根源：(i) 高斯过非线性非高斯，只保前两阶矩丢弃高阶；(ii) 用线性化近似真实输出协方差；(iii) 工作点常是估计均值而非真先验均值；(iv) 用先验均值过非线性近似输出均值，不等于真输出均值。还需算雅可比（解析或数值）。优点：轻非线性 + 高斯输入时简单快速，且若非线性局部可逆则\emph{过程可逆}（可由输出 PDF 过逆非线性精确恢复输入 PDF；sigma 点变换则不可逆）。

\paragraph{办法三：sigma 点 / 无迹变换（UKF）。}
不线性化、不求雅可比，而是用一组\emph{确定性}采样点（sigma 点）表示输入高斯，各点精确过非线性再重组矩\cite{barfoot2024state}。由 $\mathcal{N}(\boldsymbol{\mu}_x,\boldsymbol{\Sigma}_{xx})$（维度 $L$）生成 $2L+1$ 个点（Cholesky $\mathbf{L}\mathbf{L}^\top=\boldsymbol{\Sigma}_{xx}$，$\kappa$ 为用户参数）：
\begin{equation}
  \mathbf{x}_0=\boldsymbol{\mu}_x,\quad
  \mathbf{x}_i=\boldsymbol{\mu}_x+\sqrt{L+\kappa}\,\mathrm{col}_i\mathbf{L},\quad
  \mathbf{x}_{i+L}=\boldsymbol{\mu}_x-\sqrt{L+\kappa}\,\mathrm{col}_i\mathbf{L},\quad i=1\dots L,
  \label{eq:sigma-points}
\end{equation}
权 $\alpha_0=\tfrac{\kappa}{L+\kappa}$、$\alpha_i=\tfrac{1}{2(L+\kappa)}\,(i\ne0)$（和为 1，可验证 $\boldsymbol{\mu}_x=\sum_i\alpha_i\mathbf{x}_i$、$\boldsymbol{\Sigma}_{xx}=\sum_i\alpha_i(\mathbf{x}_i-\boldsymbol{\mu}_x)(\mathbf{x}_i-\boldsymbol{\mu}_x)^\top$）。各点过非线性 $\mathbf{y}_i=\mathbf{g}(\mathbf{x}_i)$，重组 $\boldsymbol{\mu}_y=\sum_i\alpha_i\mathbf{y}_i$、$\boldsymbol{\Sigma}_{yy}=\sum_i\alpha_i(\mathbf{y}_i-\boldsymbol{\mu}_y)(\mathbf{y}_i-\boldsymbol{\mu}_y)^\top$。优于线性化：免雅可比、只用线性代数（Cholesky/外积/求和）、成本与“数值雅可比的线性化”相当、不要求 $\mathbf{g}$ 光滑。（“unscented”一名源自 Julier 的玩笑，与“无香除臭剂”同词\cite{barfoot2024state}。）

\begin{example}[$f(x)=x^2$：三法对比，看清线性化为何过自信]\label{ex:three-methods}
输入 $\mathcal{N}(\mu_x,\sigma_x^2)$。
\emph{蒙特卡洛/精确}：$y=(\mu_x+\delta x)^2=\mu_x^2+2\mu_x\delta x+\delta x^2$，用 $\mathbb{E}[\delta x]=0,\mathbb{E}[\delta x^2]=\sigma_x^2,\mathbb{E}[\delta x^3]=0,\mathbb{E}[\delta x^4]=3\sigma_x^4$：
\[
  \mu_y=\mu_x^2+\sigma_x^2,\qquad \sigma_y^2=4\mu_x^2\sigma_x^2+2\sigma_x^4.
\]
\emph{线性化}：$y\approx\mu_x^2+2\mu_x\delta x$，得 $\mu_y=\mu_x^2$、$\sigma_y^2=4\mu_x^2\sigma_x^2$。对比可见\textbf{线性化均值有偏（少了 $\sigma_x^2$）、方差偏小（过自信，少了 $2\sigma_x^4$）}。
\emph{sigma 点}（$L=1$，三点 $x_0=\mu_x$、$x_{1,2}=\mu_x\pm\sqrt{1+\kappa}\,\sigma_x$）：重组得 $\mu_y=\mu_x^2+\sigma_x^2$（\emph{与 $\kappa$ 无关，恰等真值}）、$\sigma_y^2=4\mu_x^2\sigma_x^2+\kappa\sigma_x^4$。\textbf{取 $\kappa=2$ 时 $\sigma_y^2=4\mu_x^2\sigma_x^2+2\sigma_x^4$，与真值完全相同}。为什么 $\kappa=2$？因 $\kappa$ 缩放点到均值的距离，不影响前三阶矩（均值、方差、零偏度），只调四阶矩；令 sigma 点四阶矩 $=$ 高斯真实峰度 $3\sigma_x^4$ 解出 $\kappa=2$（一维高斯）。结论：\emph{捕获真均值时，sigma 点优于线性化；变换后 PDF 非高斯，均值 $\ne$ 众数}。
\end{example}

\paragraph{SPKF / UKF：把 sigma 点塞进广义高斯滤波。}
弃线性化、改用 sigma 点过非线性运动/观测模型，即 \emph{sigma 点卡尔曼滤波}（SPKF）/ 无迹卡尔曼滤波（UKF）\cite{barfoot2024state}。它把预测与校正各\emph{当作一次独立的无迹变换}（\cref{eq:sigma-points} 的“生成—过非线性—重组”三步），逐式如下，满足“每步不删减”。

\begin{algo}[SPKF/UKF 预测步]\label{algo:spkf-predict}
由先验 $\{\hat{\mathbf{x}}_{k-1},\hat{\mathbf{P}}_{k-1}\}$ 得预测 $\{\check{\mathbf{x}}_k,\check{\mathbf{P}}_k\}$\cite{barfoot2024state}：
\begin{enumerate}
  \item \textbf{堆叠先验与过程噪声}（因二者都有不确定度），$L=\dim\boldsymbol{\mu}_z$：
  \begin{equation}
    \boldsymbol{\mu}_z=\begin{bmatrix}\hat{\mathbf{x}}_{k-1}\\\mathbf{0}\end{bmatrix},\qquad
    \boldsymbol{\Sigma}_{zz}=\begin{bmatrix}\hat{\mathbf{P}}_{k-1}&\mathbf{0}\\\mathbf{0}&\boldsymbol{\Sigma}_{\mathbf{w}}\end{bmatrix}.
    \label{eq:spkf-pred-stack}
  \end{equation}
  \item \textbf{生成 $2L+1$ 个 sigma 点}（Cholesky $\mathbf{L}\mathbf{L}^\top=\boldsymbol{\Sigma}_{zz}$）：$\mathbf{z}_0=\boldsymbol{\mu}_z$，$\mathbf{z}_i=\boldsymbol{\mu}_z+\sqrt{L+\kappa}\,\mathrm{col}_i\mathbf{L}$，$\mathbf{z}_{i+L}=\boldsymbol{\mu}_z-\sqrt{L+\kappa}\,\mathrm{col}_i\mathbf{L}$（$i=1\dots L$），权 $\alpha_i$ 同 \cref{eq:sigma-points}。
  \item \textbf{拆分并过运动模型}：把 $\mathbf{z}_i=[\hat{\mathbf{x}}_{k-1,i};\mathbf{w}_{k,i}]$ 拆为状态与噪声分量，各\emph{精确}过非线性
  \begin{equation}
    \check{\mathbf{x}}_{k,i}=\mathbf{f}(\hat{\mathbf{x}}_{k-1,i},\mathbf{u}_k,\mathbf{w}_{k,i}),\qquad i=0\dots 2L.
    \label{eq:spkf-pred-pass}
  \end{equation}
  \item \textbf{重组预测矩}：
  \begin{equation}
    \check{\mathbf{x}}_k=\sum_{i=0}^{2L}\alpha_i\,\check{\mathbf{x}}_{k,i},\qquad
    \check{\mathbf{P}}_k=\sum_{i=0}^{2L}\alpha_i(\check{\mathbf{x}}_{k,i}-\check{\mathbf{x}}_k)(\check{\mathbf{x}}_{k,i}-\check{\mathbf{x}}_k)^\top.
    \label{eq:spkf-pred-recombine}
  \end{equation}
\end{enumerate}
\end{algo}

\begin{algo}[SPKF/UKF 校正步]\label{algo:spkf-correct}
第二次无迹变换，得 \cref{eq:ggf} 所需的 $\boldsymbol{\mu}_{y,k},\boldsymbol{\Sigma}_{yy,k},\boldsymbol{\Sigma}_{xy,k}$\cite{barfoot2024state}：
\begin{enumerate}
  \item \textbf{堆叠预测态与观测噪声}：$\boldsymbol{\mu}_z=[\check{\mathbf{x}}_k;\mathbf{0}]$、$\boldsymbol{\Sigma}_{zz}=\mathrm{diag}(\check{\mathbf{P}}_k,\boldsymbol{\Sigma}_{\mathbf{v}})$，重取 $L=\dim\boldsymbol{\mu}_z$、按上式重生成 sigma 点 $\mathbf{z}_i=[\check{\mathbf{x}}_{k,i};\mathbf{v}_{k,i}]$。
  \item \textbf{过观测模型}：$\check{\mathbf{z}}_{k,i}=\mathbf{h}(\check{\mathbf{x}}_{k,i},\mathbf{v}_{k,i})$，$i=0\dots2L$。
  \item \textbf{重组观测矩与互协方差}：
  \begin{align}
    \boldsymbol{\mu}_{y,k}&=\sum_{i=0}^{2L}\alpha_i\,\check{\mathbf{z}}_{k,i},\qquad
    \boldsymbol{\Sigma}_{yy,k}=\sum_{i=0}^{2L}\alpha_i(\check{\mathbf{z}}_{k,i}-\boldsymbol{\mu}_{y,k})(\check{\mathbf{z}}_{k,i}-\boldsymbol{\mu}_{y,k})^\top,\label{eq:spkf-corr-moments}\\
    \boldsymbol{\Sigma}_{xy,k}&=\sum_{i=0}^{2L}\alpha_i(\check{\mathbf{x}}_{k,i}-\check{\mathbf{x}}_k)(\check{\mathbf{z}}_{k,i}-\boldsymbol{\mu}_{y,k})^\top.\label{eq:spkf-corr-cross}
  \end{align}
  \item \textbf{代入广义高斯校正步}\cref{eq:ggf}：$\mathbf{K}_k=\boldsymbol{\Sigma}_{xy,k}\boldsymbol{\Sigma}_{yy,k}^{-1}$、$\hat{\mathbf{x}}_k=\check{\mathbf{x}}_k+\mathbf{K}_k(\mathbf{z}_k-\boldsymbol{\mu}_{y,k})$、$\hat{\mathbf{P}}_k=\check{\mathbf{P}}_k-\mathbf{K}_k\boldsymbol{\Sigma}_{xy,k}^\top$。
\end{enumerate}
\end{algo}

\begin{proposition}[轻非线性极限下 $\boldsymbol{\Sigma}_{yy}$ 与 EKF 增益等价]\label{prop:spkf-ekf}
把观测模型在预测态\emph{线性化}（仅为分析等价性，算法本身不线性化）$\check{\mathbf{z}}_{k,i}\approx\mathbf{h}(\check{\mathbf{x}}_k,\mathbf{0})+\mathbf{H}(\check{\mathbf{x}}_{k,i}-\check{\mathbf{x}}_k)+\mathbf{v}_{k,i}'$，代入 \cref{eq:spkf-corr-moments,eq:spkf-corr-cross}，则 SPKF 校正退回 EKF 校正。
\end{proposition}
\begin{proof}
线性化给 $\check{\mathbf{z}}_{k,i}-\boldsymbol{\mu}_{y,k}\approx\mathbf{H}(\check{\mathbf{x}}_{k,i}-\check{\mathbf{x}}_k)+\mathbf{v}_{k,i}'$。代入 \cref{eq:spkf-corr-moments} 的 $\boldsymbol{\Sigma}_{yy,k}$ 并展开外积，因 $\boldsymbol{\Sigma}_{zz}$ \emph{块对角}（状态与噪声 sigma 点不相关），交叉求和 $\sum_i\alpha_i(\check{\mathbf{x}}_{k,i}-\check{\mathbf{x}}_k)\mathbf{v}_{k,i}'^\top=\mathbf{0}$，余下
\[
  \boldsymbol{\Sigma}_{yy,k}\approx\mathbf{H}\underbrace{\Big(\sum_i\alpha_i(\check{\mathbf{x}}_{k,i}-\check{\mathbf{x}}_k)(\cdot)^\top\Big)}_{=\check{\mathbf{P}}_k}\mathbf{H}^\top+\underbrace{\sum_i\alpha_i\mathbf{v}_{k,i}'\mathbf{v}_{k,i}'^\top}_{=\boldsymbol{\Sigma}_{\mathbf{v}}'}=\mathbf{H}\check{\mathbf{P}}_k\mathbf{H}^\top+\boldsymbol{\Sigma}_{\mathbf{v}}'.
\]
同理 \cref{eq:spkf-corr-cross} 给 $\boldsymbol{\Sigma}_{xy,k}\approx\check{\mathbf{P}}_k\mathbf{H}^\top$。于是 $\mathbf{K}_k=\boldsymbol{\Sigma}_{xy,k}\boldsymbol{\Sigma}_{yy,k}^{-1}\approx\check{\mathbf{P}}_k\mathbf{H}^\top(\mathbf{H}\check{\mathbf{P}}_k\mathbf{H}^\top+\boldsymbol{\Sigma}_{\mathbf{v}}')^{-1}$——正是 EKF 增益（\cref{prop:ekf}）。故\textbf{SPKF 在轻非线性极限与 EKF 一致，强非线性时因保留了高阶矩而更准}\cite{barfoot2024state}。
\end{proof}

\noindent\textbf{两点工程便利}\cite{barfoot2024state}：(i) SPKF 全程不需任何解析导数（免雅可比）；(ii) 只用基本线性代数（Cholesky、外积、求和），$\mathbf{f},\mathbf{h}$ 可为黑箱函数。\emph{加性观测噪声特例}（$\mathbf{z}=\mathbf{h}(\mathbf{x})+\mathbf{v}$）下还可只用 $2N+1$ 个状态 sigma 点（$N=\dim\mathbf{x}$），观测噪声以 $+\boldsymbol{\Sigma}_{\mathbf{v}}$ 解析补到 $\boldsymbol{\Sigma}_{yy,k}$，并用 SMW 把 $\boldsymbol{\Sigma}_{yy,k}^{-1}$ 化为对一个 $(2N+1)$ 维矩阵求逆——测量维高时显著省算\cite{barfoot2024state}（完整提速推导见 \cref{der:spkf-additive}）。

\begin{derivation}[加性观测噪声下 SPKF 只需 $2N+1$ 点（$\beta$ 重加权 $+$ SMW 降维）]\label{der:spkf-additive}
设 $\mathbf{z}_k=\mathbf{h}(\mathbf{x}_k)+\mathbf{v}_k$（噪声加性），堆叠维 $L=N+M$（$N=\dim\mathbf{x}$、$M=\dim\mathbf{v}$）。因校正步 $\boldsymbol{\Sigma}_{zz}=\mathrm{diag}(\check{\mathbf{P}}_k,\boldsymbol{\Sigma}_{\mathbf{v}})$ \emph{块对角}，$2L+1$ 个 sigma 点可按块分成 $2N+1$ 个“状态点”（噪声分量取 $\mathbf{0}$）与 $2M$ 个“噪声点”（状态分量取 $\check{\mathbf{x}}_k$）。过观测模型后
\[
  \check{\mathbf{z}}_{k,j}=\begin{cases}\mathbf{h}(\check{\mathbf{x}}_{k,j}),& j=0\dots2N\ (\text{状态点}),\\[2pt] \mathbf{h}(\check{\mathbf{x}}_k)+\mathbf{v}_{k,j},& j=2N+1\dots2L\ (\text{噪声点}).\end{cases}
\]
\textbf{第一步：把噪声点折进权重（$\beta$ 重加权）。} 噪声点的确定部分 $\mathbf{h}(\check{\mathbf{x}}_k)$ 相同、随机部分 $\mathbf{v}_{k,j}$ 关于 $\mathbf{0}$ 对称（成对相消），代入 $\boldsymbol{\mu}_{y,k}=\sum_j\alpha_j\check{\mathbf{z}}_{k,j}$ 后噪声点之和退化为 $\mathbf{h}(\check{\mathbf{x}}_k)\sum_{j>2N}\alpha_j$，于是只剩 $2N+1$ 个状态点、权重重组为
\[
  \beta_0=\alpha_0+\!\!\sum_{j=2N+1}^{2L}\!\!\alpha_j,\qquad \beta_i=\alpha_i\ (i=1\dots2N),
\]
（仍和为 $1$、与原权同形）使 $\boldsymbol{\mu}_{y,k}=\sum_{j=0}^{2N}\beta_j\check{\mathbf{z}}_{k,j}$，且\emph{无近似}地
\[
  \boldsymbol{\Sigma}_{yy,k}=\sum_{j=0}^{2N}\beta_j(\check{\mathbf{z}}_{k,j}-\boldsymbol{\mu}_{y,k})(\check{\mathbf{z}}_{k,j}-\boldsymbol{\mu}_{y,k})^\top+\boldsymbol{\Sigma}_{\mathbf{v}}.
\]
代价上只调用 $\mathbf{h}$ 共 $2N+1$ 次（不再随测量维 $M$ 增长）。\textbf{第二步：SMW 把 $M\times M$ 求逆降到 $(2N+1)\times(2N+1)$。} 记 $\mathrm{col}_j\mathbf{Z}_k=\sqrt{\beta_j}\,(\check{\mathbf{z}}_{k,j}-\boldsymbol{\mu}_{y,k})$（$\mathbf{Z}_k$ 为 $M\times(2N+1)$），则 $\boldsymbol{\Sigma}_{yy,k}=\mathbf{Z}_k\mathbf{Z}_k^\top+\boldsymbol{\Sigma}_{\mathbf{v}}$，对其用矩阵求逆引理（SMW，\cref{ch:slam_est}）
\[
  \boldsymbol{\Sigma}_{yy,k}^{-1}=\boldsymbol{\Sigma}_{\mathbf{v}}^{-1}-\boldsymbol{\Sigma}_{\mathbf{v}}^{-1}\mathbf{Z}_k\big(\underbrace{\mathbf{Z}_k^\top\boldsymbol{\Sigma}_{\mathbf{v}}^{-1}\mathbf{Z}_k+\mathbf{I}}_{(2N+1)\times(2N+1)}\big)^{-1}\mathbf{Z}_k^\top\boldsymbol{\Sigma}_{\mathbf{v}}^{-1},
\]
括号内只是 $(2N+1)$ 维求逆。当 $\boldsymbol{\Sigma}_{\mathbf{v}}^{-1}$ 易得（如 $\boldsymbol{\Sigma}_{\mathbf{v}}=\sigma^2\mathbf{I}\Rightarrow\boldsymbol{\Sigma}_{\mathbf{v}}^{-1}=\sigma^{-2}\mathbf{I}$）、测量维 $M\gg N$ 时省算显著（对应 Barfoot 式 4.91--4.97）。\hfill$\square$
\end{derivation}

\paragraph{迭代：IEKF 趋众数、ISPKF 趋均值。}
把观测在\emph{任意工作点} $\mathbf{x}_{\text{op}}$ 线性化，并反复令 $\mathbf{x}_{\text{op}}\leftarrow\hat{\mathbf{x}}_k$ 重新线性化，即\emph{迭代 EKF}（IEKF）\cite{barfoot2024state}：
\begin{equation}
  \hat{\mathbf{x}}_k=\check{\mathbf{x}}_k+\mathbf{K}_k\big(\mathbf{z}_k-\mathbf{h}(\mathbf{x}_{\text{op}})-\mathbf{H}(\check{\mathbf{x}}_k-\mathbf{x}_{\text{op}})\big),\qquad
  \mathbf{x}_{\text{op}}\leftarrow\hat{\mathbf{x}}_k,
  \label{eq:iekf}
\end{equation}
首次取 $\mathbf{x}_{\text{op}}=\check{\mathbf{x}}_k$ 时 IEKF 退化为 EKF（\textbf{EKF $=$ IEKF 的第一次迭代}）；协方差只在均值收敛后算一次。\textbf{IEKF 收敛到后验的众数，即一个 MAP 估计}（与高斯牛顿同源——单时刻的 GN，\cref{ch:nlopt} 证其等价）。把 sigma 点也迭代，即\emph{迭代 sigma 点卡尔曼滤波}（ISPKF\cite{barfoot2024state}），则\textbf{趋近后验均值}。

\begin{algo}[ISPKF 校正步迭代（Sibley 等）]\label{algo:ispkf}
与 SPKF 校正的唯一差别是：sigma 点围绕\emph{可更新的工作点} $\mathbf{x}_{\text{op}}$（而非固定的 $\check{\mathbf{x}}_k$）生成，并反复重置工作点\cite{barfoot2024state}：
\begin{enumerate}
  \item 初始化 $\mathbf{x}_{\text{op}}\leftarrow\check{\mathbf{x}}_k$。
  \item 堆叠 $\boldsymbol{\mu}_z=[\mathbf{x}_{\text{op}};\mathbf{0}]$、$\boldsymbol{\Sigma}_{zz}=\mathrm{diag}(\check{\mathbf{P}}_k,\boldsymbol{\Sigma}_{\mathbf{v}})$，生成 sigma 点，过观测模型 $\mathbf{z}_{\text{op},k,i}=\mathbf{h}(\mathbf{x}_{\text{op},k,i},\mathbf{v}_{k,i})$，按 \cref{eq:spkf-corr-moments,eq:spkf-corr-cross} 重组 $\boldsymbol{\mu}_{y,k},\boldsymbol{\Sigma}_{yy,k},\boldsymbol{\Sigma}_{xy,k}$，并另算 $\boldsymbol{\Sigma}_{xx,k}=\sum_i\alpha_i(\mathbf{x}_{\text{op},k,i}-\mathbf{x}_{\text{op}})(\cdot)^\top$。
  \item 用\emph{统计量}（而非解析雅可比）做一步校正，注意新息里要补\emph{工作点不在预测均值}的修正项 $\boldsymbol{\Sigma}_{yx,k}\boldsymbol{\Sigma}_{xx,k}^{-1}(\check{\mathbf{x}}_k-\mathbf{x}_{\text{op}})$：
  \begin{equation}
    \mathbf{K}_k=\boldsymbol{\Sigma}_{xy,k}\boldsymbol{\Sigma}_{yy,k}^{-1},\quad
    \hat{\mathbf{x}}_k=\check{\mathbf{x}}_k+\mathbf{K}_k\big(\mathbf{z}_k-\boldsymbol{\mu}_{y,k}-\boldsymbol{\Sigma}_{yx,k}\boldsymbol{\Sigma}_{xx,k}^{-1}(\check{\mathbf{x}}_k-\mathbf{x}_{\text{op}})\big),\quad
    \hat{\mathbf{P}}_k=\boldsymbol{\Sigma}_{xx,k}-\mathbf{K}_k\boldsymbol{\Sigma}_{yx,k}.
    \label{eq:ispkf-update}
  \end{equation}
  \item 令 $\mathbf{x}_{\text{op}}\leftarrow\hat{\mathbf{x}}_k$，回步 2，直至 $\hat{\mathbf{x}}_k$ 变化足够小。
\end{enumerate}
首次迭代 $\mathbf{x}_{\text{op}}=\check{\mathbf{x}}_k$ 时修正项为零，\cref{eq:ispkf-update} 退回 SPKF 校正 \cref{eq:ggf}——即 \textbf{SPKF $=$ ISPKF 的第一次迭代}（与“EKF $=$ IEKF 首次迭代”平行）。
\end{algo}

\noindent\textbf{众数 vs 均值的数值证据}：对前述立体相机例（夸大差异取 $x_{\text{true}}=26$m）数值显示 $\hat{x}_{\text{iekf}}=\hat{x}_{\text{map}}=24.5694$、后验均值 $\bar{x}=24.7770$、$\hat{x}_{\text{ispkf}}=24.7414$——\textbf{IEKF 落在众数、ISPKF 落在均值附近}\cite{barfoot2024state}。更有力的是\emph{统计}证据：在与 \cref{sec:ekf} 的 MAP 偏差实验\emph{同一}设置（$10^6$ 次试验，每次从先验抽新 $x_{\text{true}}$、从观测模型抽新 $y$）下，ISPKF 的偏差 $\hat{e}_{\text{mean}}(\hat{x}_{\text{ispkf}})\approx-3.84$cm、均方误差 $\hat{e}_{\text{sq}}\approx4.32$m$^2$\cite{barfoot2024state}；对照 \cref{eq:e-sq-def} 旁 MAP 的 $-33.0$cm$/4.41$m$^2$——\textbf{均方误差几乎持平，但偏差小了约一个数量级（$\approx 10\times$）}。这把“求均值（ISPKF）优于求众数（MAP）”从定性说法坐实成可复算的量化证据：MAP 找的是偏离均值的众数、系统性偏大，而趋均值的 ISPKF 把这份偏差压掉了约九成。这也把 \cref{tab:filter-taxo} 的“IEKF$\to$众数、ISPKF$\to$均值”坐实为可复算的数字，也再次印证 \cref{sec:ekf} 的本质洞察：非线性后验偏斜，找众数（迭代线性化）与找均值（迭代 sigma 点）是两条分岔的路。

\paragraph{粒子滤波：非高斯、非线性的样本表示。}
当噪声/后验明显非高斯（多峰、重尾、被强非线性掰歪），高斯滤波那套“只存前两阶矩”就不够了。\emph{粒子滤波}（particle filter，PF）回到 \cref{thm:bayes-filter} 的另一支对策——\textbf{用一组带权随机样本}（粒子）$\{\check{\mathbf{x}}_{k,m},w_{k,m}\}_{m=1}^M$ \textbf{表示置信度本身}：经验分布 $p(\mathbf{x}_k\mid\cdot)\approx\sum_m w_{k,m}\,\delta(\mathbf{x}_k-\check{\mathbf{x}}_{k,m})$。粒子数 $M\to\infty$ 时它依分布收敛到真后验，故 PF 是“近似 PDF”一路里最不挑分布形状的方法\cite{barfoot2024state,thrun2005probabilistic}。

\begin{algo}[序贯重要性重采样粒子滤波（SIR / bootstrap）]\label{algo:pf}
由 $\{\hat{\mathbf{x}}_{k-1,m}\}_{m=1}^M$ 推到 $\{\hat{\mathbf{x}}_{k,m}\}_{m=1}^M$\cite{barfoot2024state,thrun2005probabilistic}：
\begin{enumerate}
  \item \textbf{采样（从建议分布抽预测粒子）}：从先验与过程噪声联合抽 $\big[\hat{\mathbf{x}}_{k-1,m};\mathbf{w}_{k,m}\big]\sim p(\mathbf{x}_{k-1}\mid\cdot)\,p(\mathbf{w}_k)$（实践中两因子分别抽），各过运动模型得\emph{预测粒子}
  \begin{equation}
    \check{\mathbf{x}}_{k,m}=\mathbf{f}(\hat{\mathbf{x}}_{k-1,m},\mathbf{u}_k,\mathbf{w}_{k,m}),\qquad m=1\dots M.
    \label{eq:pf-predict}
  \end{equation}
  \item \textbf{加权（用观测似然校正）}：给每个预测粒子配未归一化权重，再归一化（$\sum_m w_{k,m}=1$）
  \begin{equation}
    \tilde{w}_{k,m}=p(\mathbf{z}_k\mid\check{\mathbf{x}}_{k,m}),\qquad
    w_{k,m}=\frac{\tilde{w}_{k,m}}{\sum_{m'}\tilde{w}_{k,m'}}.
    \label{eq:pf-weight}
  \end{equation}
  似然常由模拟观测 $\check{\mathbf{z}}_{k,m}=\mathbf{h}(\check{\mathbf{x}}_{k,m},\mathbf{0})$ 代入已知噪声密度（如 $\mathcal{N}(\check{\mathbf{z}}_{k,m},\boldsymbol{\Sigma}_{\mathbf{v}})$）算出。
  \item \textbf{重采样}：按权重 $\{w_{k,m}\}$ 对粒子集\emph{有放回}抽 $M$ 个，得等权后验粒子 $\hat{\mathbf{x}}_{k,m}$（高权粒子被复制、低权被淘汰），权重重置为 $1/M$。
\end{enumerate}
\end{algo}

\paragraph{建议分布（proposal）：粒子滤波的“关节”。}
更一般地，PF 是\emph{重要性采样}（importance sampling）在贝叶斯递归上的展开：我们想从目标后验抽样却抽不动，于是从一个易抽的\emph{建议分布}（proposal）$q(\mathbf{x}_k\mid\cdot)$ 抽，再用权重修正这一“抽错了分布”的偏差，$w_{k,m}\propto w_{k-1,m}\,\dfrac{p(\mathbf{z}_k\mid\check{\mathbf{x}}_{k,m})\,p(\check{\mathbf{x}}_{k,m}\mid\hat{\mathbf{x}}_{k-1,m})}{q(\check{\mathbf{x}}_{k,m}\mid\cdot)}$。\textbf{bootstrap 版即取建议分布 $=$ 运动模型先验} $q=p(\mathbf{x}_k\mid\mathbf{x}_{k-1},\mathbf{u}_k)$，于是分子的转移项被约掉、权重退化为纯似然 \cref{eq:pf-weight}——实现最简，但因“没看观测就盲目按运动撒点”，观测很准时大多数粒子落在似然尾部、权重接近零，效率低\cite{thrun2005probabilistic}。好的建议分布应\emph{把当前观测也纳入}（如以 EKF/UKF 在每个粒子处算一个局部高斯当 proposal），让粒子撒向高似然区。

\begin{pitfall}[粒子退化与有效样本数 $N_{\text{eff}}$：别每步都重采样，也别从不重采样]
\emph{权退化}（weight degeneracy）：迭代若干步后，几乎所有权重集中到\emph{极少数}粒子上，其余粒子权重趋零却仍占着计算量——经验分布实际只由几个点支撑，方差爆炸。量化退化程度用\textbf{有效样本数}\cite{sarkka2013bayesian,thrun2005probabilistic}
\begin{equation}
  N_{\text{eff}}=\frac{1}{\sum_{m=1}^{M}w_{k,m}^2}\in[1,\,M],
  \label{eq:pf-neff}
\end{equation}
权重全相等（$w_{k,m}=1/M$）时 $N_{\text{eff}}=M$（最健康）；退化到单粒子（某 $w=1$）时 $N_{\text{eff}}=1$（最坏）。\textbf{正确策略是“按需重采样”}：仅当 $N_{\text{eff}}<N_{\text{thresh}}$（常取 $M/2$）才重采样，否则只更新并\emph{携带}权重到下一步。两个极端都错：\emph{每步必重采样}会因重复抽样不断丢弃多样性、加剧\emph{样本贫化}（sample impoverishment，粒子collapse到少数重复值）；\emph{从不重采样}则任由退化、最终所有权重压到一个粒子。重采样本身不改变期望、只重置权重方差，故只在确有必要时做。
\end{pitfall}

\paragraph{重采样的实现：Madow 系统重采样。}
最简单的“按权多项式抽样”方差偏大；Madow（1949）的\emph{系统重采样}方差更小、$O(M)$ 一遍扫完\cite{barfoot2024state}：由归一化权建累积分布的 bin 边界 $\beta_m=\sum_{n=1}^{m}w_{k,n}$（$\beta_M=1$），它把 $[0,1)$ 切成 $M$ 个宽度等于各粒子权重的小区间；取单个随机起点 $\rho\sim\mathrm{Uniform}[0,\tfrac1M)$，令指针 $\rho,\rho+\tfrac1M,\rho+\tfrac2M,\dots,\rho+\tfrac{M-1}{M}$ 等步长扫过 $[0,1)$，每个指针落入的 bin 对应的粒子被选入新集合。等步长保证“凡权重大于 $1/M$ 的粒子至少被选一次”，避免了多项式抽样可能整批漏掉中等权粒子的高方差。

\begin{insight}
\textbf{粒子滤波回指 FastSLAM/RBPF}：\cref{ch:intro} 埋过伏笔（“若把运动模型描述为更一般非高斯分布的样本，则导出 Rao--Blackwell 化的粒子滤波 FastSLAM”，\cref{eq:intro-dw-update} 那段）。其核心 \emph{Rao--Blackwell 化}（RBPF）是：当状态可拆成“给定一部分后另一部分条件高斯”时，\textbf{只对难的那部分撒粒子、对条件高斯部分用解析 KF}，从而大幅降低所需粒子数。FastSLAM 正是这一思想在 SLAM 上的落地——机器人\emph{轨迹}用粒子表示（每个粒子是一条候选轨迹），\emph{每个粒子条件下的各路标}用独立小 EKF 解析维护（给定轨迹后路标相互独立，这正是 \cref{ch:slam_est} “给定位姿则地图可分”的稀疏结构）。这把“维度随路标数指数爆炸”降为“路标数 $\times$ 粒子数”的线性代价，是古典 SLAM 时代与 EKF-SLAM 并立的两大主线之一（\cref{ch:vio} 再述其与现代优化后端的取舍）。
\end{insight}

\noindent\textbf{优缺点小结}：PF 处理非高斯 $+$ 非线性，甚至不需 $\mathbf{f},\mathbf{h}$ 的解析式或可微性，且是 anytime 算法（粒子越多越准）。代价是：达到给定精度所需粒子数随状态维\emph{指数}增长，高维（如直接对 $6$ 维位姿 $+$ 大量路标撒点）易样本贫化，故实践多走 RBPF 把可解析的维度剥离；且与所有滤波器一样，再多粒子也\emph{不能优于}观测信息设定的 CRLB（\cref{ch:slam_est}）\cite{barfoot2024state}。

\begin{table}[H]\centering\small
\caption{高斯/贝叶斯滤波器谱系（顶为精确但不可实现的贝叶斯滤波）}
\label{tab:filter-taxo}
\begin{tabular}{llll}
\toprule
滤波器 & 置信度表示 / 过非线性的方法 & 迭代 & 逼近完整后验的 \\
\midrule
贝叶斯滤波 & 精确 PDF（不可实现） & — & 完整后验 \\
直方图 / 离散贝叶斯 & 离散网格质量（非参数） & — & 任意（含多峰，受网格分辨率限） \\
EKF & 高斯 · 线性化 & 否 & 无明确关系 \\
IEKF & 高斯 · 线性化（重线性化） & 是 & \textbf{众数（MAP）} \\
UKF/SPKF & 高斯 · sigma 点 & 否 & 无明确关系 \\
迭代 SPKF（ISPKF） & 高斯 · sigma 点 & 是 & \textbf{均值} \\
粒子滤波 & 加权样本 · 蒙特卡洛 & — & 任意（非高斯） \\
\midrule
\multicolumn{4}{l}{\emph{以下三者不在“如何表示 PDF/过非线性”这条主轴上，列此以备查（详见章末定理速查表）：}} \\
信息滤波 IF & 高斯 · 信息形式（$\boldsymbol{\Lambda},\boldsymbol{\eta}$） & 否 & 同 KF（KF 的对偶\emph{形式}，非新近似） \\
互补滤波 & 无协方差 · 频域高/低通互补 & 否 & 无（确定性增益，不表 PDF） \\
InEKF & 高斯 · 流形不变误差线性化 & 否 & 同 IEKF，但传播\emph{与状态无关}（\cref{paper:inekf}） \\
\bottomrule
\end{tabular}
\end{table}

\noindent 表末三行与上半部不属同一分类轴：IF 是 KF 的\emph{信息形式对偶}（无独立的 PDF 近似策略，\cref{prop:info-filter}）；互补滤波\emph{不传播协方差}（\cref{eq:complementary-filter}）；InEKF 是高斯滤波的\emph{流形不变}变体（\cref{paper:inekf}）——三者均在本章\emph{章末定理速查表}另有出处，故并入此表仅为读者按“滤波族”检索时不致遗漏。

\paragraph{（选读）滑动窗口滤波（SWF）：递归与批量的折中}\label{par:swf}
\textbf{从优化视角看，EKF 是劣的}：它不迭代至收敛、雅可比只算一次，且把马尔可夫性\emph{永久}内建。批量 GN（\cref{ch:nlopt}）对整条轨迹迭代、可重线性化，但要离线、非常数时间。\emph{滑动窗口滤波}（SWF\cite{barfoot2024state}）取两者之长：在一个固定大小的时间窗口上做小批量 GN、迭代至收敛，然后窗口右扩一步纳入新时刻、左缩一步\emph{边缘化}最旧状态以保持窗口大小。关键是\textbf{收缩用边缘化而非朴素删除}：把窗口批量视为联合高斯，对最旧状态 $\mathbf{x}_0$ 用舒尔补消去（信息矩阵 $\mathbf{H}^\top\mathbf{W}^{-1}\mathbf{H}$ 块三对角，消元后仍块三对角），其信息以“先验块”形式留在保留变量上。这样：在\emph{多于一个}时步上迭代（故收敛到比 IEKF 更接近完整批量解的结果），又保持常数时间在线。\textbf{IEKF $=$ 窗口大小为 1、且只迭代校正步的 SWF}；SWF 对任意窗口大小迭代预测 $+$ 校正两步。这是 VIO 中“边缘化老关键帧”做法（\cref{ch:vio}）的理论原型。

\begin{practice}
\begin{enumerate}
  \item 把 $\mathcal{N}(\mu_x,\sigma_x^2)$ 过 $f(x)=x^3$，分别用线性化与 sigma 点（含均值点）求输出均值，说明线性化为何漏掉了与 $\sigma_x^2$ 有关的项。提示：用高斯高阶矩 $\mathbb{E}[\delta x^4]=3\sigma_x^4$。
  \item 说明取 $\mathbf{x}_{\text{op}}=\check{\mathbf{x}}_k$ 时 IEKF 的首次迭代恰为 EKF（把它代入 \cref{eq:iekf}）。
  \item （概念）结合 \cref{tab:filter-taxo} 解释“EKF 与完整后验无明确关系，迭代后才靠拢众数或均值”，并说明这与 \cref{ex:three-methods} 的“线性化过自信”如何呼应。答错回看 \cref{ex:three-methods} 与本节迭代段。
\end{enumerate}
\end{practice}

% ════════════════════════════════════════════════════════════════
\section{估计器的健康度：无偏、一致与偏置}
\label{sec:consistency}

EKF 的“三宗罪”之一是\emph{不一致}（报告的协方差比真实误差乐观）。本节给出诊断一致性的标准工具 NEES/NIS，并用一个完整的误差动力学推导说明\textbf{未补偿的偏置如何让 KF 系统性过自信、且偏置效应随时间无界增长}；再用状态增广把偏置纳入估计，并以可观测性判断它\emph{能否}估出。本节脉络层为\textbf{次优}（诊断与调参用）。

\paragraph{无偏与一致：定义与统计检验。}
设估计 $\mathcal{N}(\hat{\mathbf{x}}_k,\hat{\mathbf{P}}_k)$、有真值 $\mathbf{x}_{\text{true},k}$，单步误差 $\hat{\mathbf{e}}_k=\hat{\mathbf{x}}_k-\mathbf{x}_{\text{true},k}$。\textbf{无偏}：$\mathbb{E}[\hat{\mathbf{e}}_k]=\mathbf{0}$（误差在零附近徘徊）。\textbf{一致}：误差大体落在 $\pm3\sigma_k$ 包络内（$99.7\%$），标量化为 $\mathbb{E}[\hat{e}_k^2/\hat{P}_k]=1$\cite{barfoot2024state}。实践用\emph{遍历假设}（沿时间平均代替沿试验平均）+ 统计检验：样本均值 $\hat{e}_{\text{mean}}\sim\mathcal{N}(\mu,\sigma^2/K)$ 落入置信区间则判无偏；一致性用卡方检验 $\sum_k\hat{e}_k^2/\hat{P}_k\in[Q_{\chi^2(K)}(\ell),Q_{\chi^2(K)}(u)]$（$\ell=0.025,u=0.975$ 为 $95\%$ 区间）。

\begin{definition}[NEES 与 NIS]\label{def:nees-nis}
$N$ 维估计的一致性用\emph{归一化估计误差平方}（NEES）：
\begin{equation}
  \epsilon_{\text{nees},k}=\hat{\mathbf{e}}_k^\top\hat{\mathbf{P}}_k^{-1}\hat{\mathbf{e}}_k,\qquad
  \mathbb{E}[\epsilon_{\text{nees},k}]=N\ \text{（一致）},
  \label{eq:nees}
\end{equation}
这是平方马氏距离；$\mathbb{E}[\epsilon_{\text{nees}}]<N$ 称\emph{欠自信}，$>N$ 称\emph{过自信}。无真值时用\emph{归一化新息平方}（NIS）：
\begin{equation}
  \epsilon_{\text{nis},k}=\mathbf{e}_{y,k}^\top\mathbf{S}_k^{-1}\mathbf{e}_{y,k},\quad
  \mathbf{e}_{y,k}=\mathbf{z}_k-\mathbf{h}(\check{\mathbf{x}}_k),\quad
  \mathbf{S}_k=\mathbf{H}\check{\mathbf{P}}_k\mathbf{H}^\top+\boldsymbol{\Sigma}_{\mathbf{v}},\quad
  \mathbb{E}[\epsilon_{\text{nis},k}]=M.
  \label{eq:nis}
\end{equation}
\end{definition}
NIS 用\emph{预测态} $\check{\mathbf{x}}_k$ 算新息（$\mathbf{z}_k$ 尚未并入），故可\emph{在线}评估、无需真值；统计检验 $\sum_k\epsilon_{\text{nis},k}\in[Q_{\chi^2(M(K+1))}(\ell),Q_{\chi^2(M(K+1))}(u)]$。同理可定义\emph{运动 NIS}（用预测误差 $\mathbf{e}_{v,k}=\mathbf{f}(\hat{\mathbf{x}}_{k-1},\mathbf{u}_k)-\hat{\mathbf{x}}_k$、$\mathbf{S}_{v,k}=\mathbf{F}\hat{\mathbf{P}}_{k-1}\mathbf{F}^\top+\boldsymbol{\Sigma}_{\mathbf{w}}$、判据 $\mathbb{E}=N$）。NEES/NIS 是 KF/EKF 关键参数调参的标准依据\cite{barfoot2024state}。

\textbf{大 $K$ 的实用近似.}\;长轨迹（$K>100$）时直接查卡方分位数 $Q_{\chi^2(K)}$ 不便，可用标准正态分位数闭式逼近（Bar-Shalom 等\cite{barfoot2024state}）
\begin{equation}
  Q_{\chi^2(K)}(y)\approx\tfrac12\big(Q_{\mathcal{N}(0,1)}(y)+\sqrt{2K-1}\big)^2,
  \label{eq:chi2-largeK}
\end{equation}
便于在长轨迹一致性检验中手算上下界。\textbf{相关轨迹的聚合 NEES.}\;上面沿时间平均的“遍历版”默认各时刻误差\emph{独立}，但同一条轨迹的估计在不同时刻通常\emph{高度相关}，独立性假设失效、逐时刻卡方边界不再可信。此时改用\emph{整条轨迹的单样本聚合 NEES}：把全轨迹误差堆成 $\hat{\mathbf{e}}=[\hat{\mathbf{e}}_0;\dots;\hat{\mathbf{e}}_K]$、配\emph{完整}（含时刻间相关块的）协方差 $\hat{\mathbf{P}}$，
\begin{equation}
  \epsilon_{\text{nees}}=\hat{\mathbf{e}}^\top\hat{\mathbf{P}}^{-1}\hat{\mathbf{e}},\qquad \mathbb{E}[\epsilon_{\text{nees}}]=NK,
  \label{eq:nees-aggregate}
\end{equation}
判据为 $NK$（因只有一条全轨迹样本）\cite{barfoot2024state}。

\begin{theorem}[未补偿偏置使 KF 过自信，且效应无界增长]\label{thm:bias-overconfident}
设带偏置的 LTI 系统 $\mathbf{x}_k=\mathbf{A}\mathbf{x}_{k-1}+\mathbf{B}(\mathbf{u}_k+\bar{\mathbf{u}})+\mathbf{w}_k$、$\mathbf{z}_k=\mathbf{C}\mathbf{x}_k+\bar{\mathbf{z}}+\mathbf{v}_k$，其中 $\bar{\mathbf{u}},\bar{\mathbf{z}}$ 为未知输入/观测偏置，但 KF 按\emph{零偏置}模型运行。则 KF 估计\emph{有偏}（$\mathbb{E}[\hat{\mathbf{e}}_k]\ne\mathbf{0}$）、且报告的协方差比真实误差协方差\emph{偏小}（过自信、不一致），无论偏置正负皆如此，效应随 $k$ 无界增长。
\end{theorem}
\begin{proof}[误差动力学（逐步代数）]
定义预测/校正误差 $\check{\mathbf{e}}_k=\check{\mathbf{x}}_k-\mathbf{x}_k$、$\hat{\mathbf{e}}_k=\hat{\mathbf{x}}_k-\mathbf{x}_k$。由 KF 递推与真实（带偏置）动力学相减，得误差动力学
\[
  \check{\mathbf{e}}_k=\mathbf{A}\hat{\mathbf{e}}_{k-1}-(\mathbf{B}\bar{\mathbf{u}}+\mathbf{w}_k),\qquad
  \hat{\mathbf{e}}_k=(\mathbf{I}-\mathbf{K}_k\mathbf{C})\check{\mathbf{e}}_k+\mathbf{K}_k(\bar{\mathbf{z}}+\mathbf{v}_k).
\]
取期望（$\mathbb{E}[\mathbf{w}_k]=\mathbb{E}[\mathbf{v}_k]=\mathbf{0}$，$\mathbb{E}[\hat{\mathbf{e}}_0]=\mathbf{0}$）：$\mathbb{E}[\check{\mathbf{e}}_1]=-\mathbf{B}\bar{\mathbf{u}}$、$\mathbb{E}[\hat{\mathbf{e}}_1]=-(\mathbf{I}-\mathbf{K}_1\mathbf{C})\mathbf{B}\bar{\mathbf{u}}+\mathbf{K}_1\bar{\mathbf{z}}$——只要 $\bar{\mathbf{u}}\ne\mathbf{0}$ 或 $\bar{\mathbf{z}}\ne\mathbf{0}$ 即\emph{有偏}。再算二阶矩：把 $\check{\mathbf{e}}_1$ 写成“零均值部分 $+$ 偏置常量 $-\mathbf{B}\bar{\mathbf{u}}$”，外积取期望，交叉项因零均值部分期望为零而消失，得
\[
  \mathbb{E}[\check{\mathbf{e}}_1\check{\mathbf{e}}_1^\top]=\check{\mathbf{P}}_1+(\mathbf{B}\bar{\mathbf{u}})(\mathbf{B}\bar{\mathbf{u}})^\top
  \;\Longrightarrow\;
  \check{\mathbf{P}}_1=\mathbb{E}[\check{\mathbf{e}}_1\check{\mathbf{e}}_1^\top]-\mathbb{E}[\check{\mathbf{e}}_1]\mathbb{E}[\check{\mathbf{e}}_1]^\top,
\]
即 KF 报告的 $\check{\mathbf{P}}_1$ \emph{比真实误差不确定度小}了一个偏置外积（恒半正定）。校正同理 $\hat{\mathbf{P}}_1=\mathbb{E}[\hat{\mathbf{e}}_1\hat{\mathbf{e}}_1^\top]-\mathbb{E}[\hat{\mathbf{e}}_1]\mathbb{E}[\hat{\mathbf{e}}_1]^\top$。因偏置外积\emph{非负且与偏置正负无关}，KF 恒\textbf{过自信}；递推可见偏置项随 $k$ 累积、无界增长\cite{barfoot2024state}。
\end{proof}

\paragraph{已知偏置：直接补偿；未知偏置：状态增广。}
若偏置已知，把它塞进 KF 五式即恢复无偏一致：预测均值加 $\mathbf{B}\bar{\mathbf{u}}$、校正新息减 $\bar{\mathbf{z}}$。但通常偏置未知（且随时间漂移）。此时把偏置建成\emph{随机游走} $\bar{\mathbf{u}}_k=\bar{\mathbf{u}}_{k-1}+\mathbf{s}_k$（$\mathbf{s}_k\sim\mathcal{N}(\mathbf{0},\mathbf{W})$）纳入\emph{增广状态} $\mathbf{x}_k'=[\mathbf{x}_k;\bar{\mathbf{u}}_k]$，得增广系统 $\mathbf{A}'=\big[\begin{smallmatrix}\mathbf{A}&\mathbf{B}\\\mathbf{0}&\mathbf{I}\end{smallmatrix}\big]$、$\boldsymbol{\Sigma}_{\mathbf{w}}'=\mathrm{diag}(\boldsymbol{\Sigma}_{\mathbf{w}},\mathbf{W})$、$\mathbf{C}'=[\mathbf{C}\ \mathbf{0}]$，回到“无偏系统”再跑 KF。\textbf{能否估出，取决于增广可观测性}：原系统 $\mathrm{rank}\,\mathcal{O}=N$ 不够，须 $\mathrm{rank}\,\mathcal{O}'=N+U$（$U=\dim\bar{\mathbf{u}}$），$\mathcal{O}'=[\mathbf{C}';\mathbf{C}'\mathbf{A}';\dots;\mathbf{C}'\mathbf{A}'^{(N+U-1)}]$。

\begin{example}[偏置可观测性三例：加速度偏置可观、双偏置不可观、测量偏置（极简 SLAM）不可观]\label{ex:bias-obs}
取小车 $\mathbf{A}=\big[\begin{smallmatrix}1&1\\0&1\end{smallmatrix}\big]$、$\mathbf{C}=[1\ 0]$（位置可测），原系统 $\mathrm{rank}\,\mathcal{O}=2=N$ 可观\cite{barfoot2024state}。
\begin{itemize}
  \item \textbf{加速度上的输入偏置}（$\mathbf{B}=[0;1]$，$U=1$）：$\mathcal{O}'=\big[\begin{smallmatrix}1&0&0\\1&1&0\\1&2&1\end{smallmatrix}\big]$，$\mathrm{rank}=3=N+U$，\emph{可观}——偏置能被估出。
  \item \textbf{速度与加速度上都有输入偏置}（$\mathbf{B}=\mathbf{I}_2$，$U=2$）：$\mathcal{O}'$ 第 2、3 列相同，$\mathrm{rank}=3<4=N+U$，\emph{不可观}——估不出。
  \item \textbf{测量偏置（极简 SLAM）}（$\bar{\mathbf{z}}\ne\mathbf{0}$，$\mathbf{C}'=[\mathbf{C}\ 1]$，$U=1$）：$\mathcal{O}'=\big[\begin{smallmatrix}1&0&1\\1&1&1\\1&2&1\end{smallmatrix}\big]$ 第 1、3 列相同，$\mathrm{rank}=2<3$，\emph{不可观}，零空间 $\mathrm{span}\{[1,0,-1]^\top\}$ 表示“把小车与地标一起平移、测量不变”。此时批量解有\emph{无穷多解}、KF 把偏置\emph{停在初值}。这正是 SLAM 全局平移不可观的极简原型（“定位”=车状态、“地图”=地标=偏置负值）。
\end{itemize}
\end{example}

\paragraph{协方差从哪来：监督式 vs 自适应估计。}
KF/EKF 的 $\boldsymbol{\Sigma}_{\mathbf{w}},\boldsymbol{\Sigma}_{\mathbf{v}}$ 常靠数据手册或试错，不准就直接导致不一致。两条从数据估出协方差的系统办法\cite{barfoot2024state}：\emph{监督式}（有 groundtruth）——在与目标场景相近、且有高质量真值 $\mathbf{x}_{\text{true}}$ 的训练集上，先算误差样本均值 $\bar{\mathbf{e}}_{v}=\tfrac{1}{K}\sum_{k}\mathbf{e}_{v,k}(\mathbf{x}_{\text{true}})$、$\bar{\mathbf{e}}_{y}=\tfrac{1}{K+1}\sum_{k}\mathbf{e}_{y,k}(\mathbf{x}_{\text{true}})$（去偏），再用 Bessel 修正的样本协方差估
\begin{equation}
  \boldsymbol{\Sigma}_{\mathbf{w}}=\frac{1}{K-1}\sum_{k=1}^{K}(\mathbf{e}_{v,k}-\bar{\mathbf{e}}_v)(\mathbf{e}_{v,k}-\bar{\mathbf{e}}_v)^\top,\quad
  \boldsymbol{\Sigma}_{\mathbf{v}}=\frac{1}{K}\sum_{k=0}^{K}(\mathbf{e}_{y,k}-\bar{\mathbf{e}}_y)(\mathbf{e}_{y,k}-\bar{\mathbf{e}}_y)^\top,
  \label{eq:eskf-supcov}
\end{equation}
估出后用于无真值的实际运行（用 \cref{def:nees-nis} 的一致性检验在\emph{独立验证集}上核验质量）。\emph{自适应}（无 groundtruth）——用拖尾窗 $L$ 个新息 $\mathbf{e}_{y,\ell}=\mathbf{z}_\ell-\mathbf{h}(\check{\mathbf{x}}_\ell)$ 的样本协方差 $\mathbf{S}_{y,k}$，令其等于理论新息协方差 $\mathbf{H}\check{\mathbf{P}}\mathbf{H}^\top+\boldsymbol{\Sigma}_{\mathbf{v}}$ 反解
\begin{equation}
  \boldsymbol{\Sigma}_{\mathbf{v},k}=\mathbf{S}_{y,k}-\frac{1}{L}\sum_{\ell=k-1}^{k-L}\mathbf{H}_\ell\check{\mathbf{P}}_\ell\mathbf{H}_\ell^\top,
  \label{eq:eskf-adapcov}
\end{equation}
（过程噪声同理用预测误差 $\mathbf{e}_{v,k}=\mathbf{f}(\hat{\mathbf{x}}_{k-1},\mathbf{u}_k)-\hat{\mathbf{x}}_k$ 与 $\mathbf{F}\hat{\mathbf{P}}\mathbf{F}^\top+\boldsymbol{\Sigma}_{\mathbf{w}}$ 反解 $\boldsymbol{\Sigma}_{\mathbf{w},k}$）。它假设滤波器在估协方差前已无偏一致、噪声加性、协方差变化不太快——\textbf{本质就是让滤波器按 NIS 检验变得一致}（与 \cref{thm:robust-iw} 的 MAP 协方差估计是同一思想的两个落点：一个估常值噪声、一个为每条边膨胀协方差）。

\begin{pitfall}[思维陷阱：把 NEES/NIS 当“越小越好”]
新手以为 $\epsilon_{\text{nees}}$ 越小代表估计越准。\textbf{实际上}判据是 $\mathbb{E}[\epsilon_{\text{nees}}]=N$：太大（过自信，协方差设小或偏置未补偿，\cref{thm:bias-overconfident}）和太小（欠自信，协方差设大）\emph{都是病}。过自信尤其危险——滤波“自信地走向错误”，新息被错误地小权处理，最终发散。\textbf{正确做法}：跑一致性检验，过自信就增大 $\boldsymbol{\Sigma}$ 或排查偏置/线性化，欠自信就减小 $\boldsymbol{\Sigma}$；自适应协方差估计（用拖尾窗样本协方差反解 $\boldsymbol{\Sigma}_{\mathbf{v}},\boldsymbol{\Sigma}_{\mathbf{w}}$）本质就是\emph{让滤波器按 NIS 检验变得一致}\cite{barfoot2024state}。
\end{pitfall}

\begin{practice}
\begin{enumerate}
  \item 对 \cref{ex:bias-obs} 的第一例，写出 $\mathbf{A}',\mathbf{C}'$ 并自行算 $\mathcal{O}'$ 验证 $\mathrm{rank}=3$。
  \item （证明）由 \cref{thm:bias-overconfident} 的 $\hat{\mathbf{P}}_k=\mathbb{E}[\hat{\mathbf{e}}_k\hat{\mathbf{e}}_k^\top]-\mathbb{E}[\hat{\mathbf{e}}_k]\mathbb{E}[\hat{\mathbf{e}}_k]^\top$ 说明：为什么不管 $\bar{\mathbf{u}}$ 正负，KF 都过自信？答案要点：减去的是外积（半正定），与符号无关。
  \item （开放）某 EKF 跑出 $\sum_k\epsilon_{\text{nis},k}\approx 3MK$（远超 $MK$），列出三种可能原因与对应排查。答案要点：协方差设太小 / 偏置未补偿 / 线性化失效；分别对应增大 $\boldsymbol{\Sigma}$、状态增广、改 IEKF 或重线性化。答错回看本节 pitfall 与 \cref{ex:bias-obs}。
\end{enumerate}
\end{practice}

% ════════════════════════════════════════════════════════════════
\section{误差状态卡尔曼滤波（ESKF）}
\label{sec:eskf}

估计含姿态的状态时，直接把四元数/旋转放进 KF 状态会出问题。ESKF 是 VIO、惯导融合的事实标准\cite{sola2017quaternion}，且与本书\emph{右扰动}约定天然契合。本节按 Solà\cite{sola2017quaternion} 的三态分解，把每条误差动力学与每个雅可比都\emph{独立推完}，读者无须翻原文。

\paragraph{动机：姿态不能当普通向量塞进状态。}
四元数有 $4$ 参数却只有 $3$ 自由度（过参数化）$\Rightarrow$ 其 $4\times4$ 协方差奇异、需强加单位范数约束；欧拉角有万向锁。直接 EKF 还把线性化点放在远离原点的姿态上，有效性存疑。ESKF 绕开这些：把状态分成\emph{标称态}（大信号，非线性确定性积分）与\emph{误差态}（小信号，最小 $3$ 参数，线性高斯滤波），姿态误差用 $\delta\boldsymbol{\theta}\in\mathbb{R}^3$ 表示。

\begin{definition}[三种状态]\label{def:eskf-states}
真值 $=$ 标称 $\oplus$ 误差：向量分量用加法，姿态分量用\emph{右乘}流形合成（与 \cref{ch:lie} 右扰动一致）
\begin{equation}
  \mathbf{x}_t=\mathbf{x}\oplus\delta\mathbf{x},\qquad
  \mathbf{p}_t=\mathbf{p}+\delta\mathbf{p},\quad
  \mathbf{v}_t=\mathbf{v}+\delta\mathbf{v},\quad
  \mathbf{R}_t=\mathbf{R}\,\mathrm{Exp}(\delta\boldsymbol{\theta}),
  \label{eq:eskf-states}
\end{equation}
（四元数写法 $\mathbf{q}_t=\mathbf{q}\otimes\delta\mathbf{q}$，$\delta\mathbf{q}=\mathrm{Exp}(\delta\boldsymbol{\theta})=e^{\delta\boldsymbol{\theta}/2}\approx[1;\tfrac12\delta\boldsymbol{\theta}]$；$\delta\mathbf{R}=\mathrm{Exp}([\delta\boldsymbol{\theta}]_\times)$）。bias、重力等也用加法：$\mathbf{a}_{bt}=\mathbf{a}_b+\delta\mathbf{a}_b$、$\boldsymbol{\omega}_{bt}=\boldsymbol{\omega}_b+\delta\boldsymbol{\omega}_b$、$\mathbf{g}_t=\mathbf{g}+\delta\mathbf{g}$。完整误差态 $\delta\mathbf{x}=[\delta\mathbf{p};\delta\mathbf{v};\delta\boldsymbol{\theta};\delta\mathbf{a}_b;\delta\boldsymbol{\omega}_b;\delta\mathbf{g}]\in\mathbb{R}^{18}$。
\end{definition}

\begin{insight}
\textbf{为何用误差状态（四条优势）}\cite{sola2017quaternion}：(1) 姿态误差\emph{最小参数化}（$3$ 参数 $=3$ 自由度），\textbf{无过参数化、协方差不奇异}；(2) 误差态始终在\emph{原点附近}，远离万向锁/奇异，\textbf{线性化有效性恒成立}（直接呼应 EKF 三宗罪之“工作点偏差”被消除）；(3) 误差小 $\Rightarrow$ 二阶项可忽略，\textbf{雅可比简单}（多为单位块）、计算快；(4) 大信号动力学都进了标称态，误差动力学\emph{慢}，\textbf{校正可低于预测频率}施加（高频 IMU 预测 $+$ 低频外部校正）。
\end{insight}

\paragraph{真态与名义态运动学。}
以 IMU 融合为例（角速度、角误差均在\emph{机体系局部}定义，本书右扰动主线）。\emph{真态}连续运动学（$\mathbf{a}_m,\boldsymbol{\omega}_m$ 为 IMU 读数，$\mathbf{a}_n,\boldsymbol{\omega}_n$ 测量噪声，$\mathbf{a}_w,\boldsymbol{\omega}_w$ 偏置随机游走噪声）\cite{sola2017quaternion}：
\begin{equation}
  \dot{\mathbf{v}}_t=\mathbf{R}_t(\mathbf{a}_m-\mathbf{a}_{bt}-\mathbf{a}_n)+\mathbf{g}_t,\quad
  \dot{\mathbf{q}}_t=\tfrac12\mathbf{q}_t\otimes(\boldsymbol{\omega}_m-\boldsymbol{\omega}_{bt}-\boldsymbol{\omega}_n),\quad
  \dot{\mathbf{a}}_{bt}=\mathbf{a}_w,\ \dot{\boldsymbol{\omega}}_{bt}=\boldsymbol{\omega}_w,\ \dot{\mathbf{g}}_t=\mathbf{0},
  \label{eq:eskf-true-kin}
\end{equation}
$\dot{\mathbf{p}}_t=\mathbf{v}_t$。\emph{名义态}是其无噪声版（去掉所有 $\mathbf{a}_n,\boldsymbol{\omega}_n,\mathbf{a}_w,\boldsymbol{\omega}_w$）：
\begin{equation}
  \dot{\mathbf{p}}=\mathbf{v},\quad
  \dot{\mathbf{v}}=\mathbf{R}(\mathbf{a}_m-\mathbf{a}_b)+\mathbf{g},\quad
  \dot{\mathbf{q}}=\tfrac12\mathbf{q}\otimes(\boldsymbol{\omega}_m-\boldsymbol{\omega}_b),\quad
  \dot{\mathbf{a}}_b=\dot{\boldsymbol{\omega}}_b=\dot{\mathbf{g}}=\mathbf{0}.
  \label{eq:eskf-nominal-kin}
\end{equation}
名义态吸收“大信号”，会因丢弃噪声而累积误差——这些误差正是误差态要追踪的。

\begin{theorem}[误差态连续运动学（右/局部扰动）]\label{thm:eskf-error-kin}
线性化后的误差态动力学为\cite{sola2017quaternion}
\begin{align}
  \dot{\delta\mathbf{p}}&=\delta\mathbf{v},\label{eq:eskf-err-p}\\
  \dot{\delta\mathbf{v}}&=-\mathbf{R}[\mathbf{a}_m-\mathbf{a}_b]_\times\,\delta\boldsymbol{\theta}-\mathbf{R}\,\delta\mathbf{a}_b+\delta\mathbf{g}-\mathbf{R}\mathbf{a}_n,\label{eq:eskf-err-v}\\
  \dot{\delta\boldsymbol{\theta}}&=-[\boldsymbol{\omega}_m-\boldsymbol{\omega}_b]_\times\,\delta\boldsymbol{\theta}-\delta\boldsymbol{\omega}_b-\boldsymbol{\omega}_n,\label{eq:eskf-err-theta}\\
  \dot{\delta\mathbf{a}}_b&=\mathbf{a}_w,\quad \dot{\delta\boldsymbol{\omega}}_b=\boldsymbol{\omega}_w,\quad \dot{\delta\mathbf{g}}=\mathbf{0}.\label{eq:eskf-err-bias}
\end{align}
\end{theorem}
\begin{proof}[线速度误差 \cref{eq:eskf-err-v} 的逐步证明]
记体坐标大/小信号加速度 $\mathbf{a}_B\triangleq\mathbf{a}_m-\mathbf{a}_b$、$\delta\mathbf{a}_B\triangleq-\delta\mathbf{a}_b-\mathbf{a}_n$，则真态 $\dot{\mathbf{v}}_t=\mathbf{R}_t(\mathbf{a}_B+\delta\mathbf{a}_B)+\mathbf{g}+\delta\mathbf{g}$、名义 $\dot{\mathbf{v}}=\mathbf{R}\mathbf{a}_B+\mathbf{g}$。用右扰动一阶近似 $\mathbf{R}_t=\mathbf{R}(\mathbf{I}+[\delta\boldsymbol{\theta}]_\times)+O(\|\delta\boldsymbol{\theta}\|^2)$ 展开真态（$\dot{\mathbf{v}}_t=\dot{\mathbf{v}}+\dot{\delta\mathbf{v}}$）：
\[
  \mathbf{R}\mathbf{a}_B+\mathbf{g}+\dot{\delta\mathbf{v}}=\mathbf{R}\mathbf{a}_B+\mathbf{R}\delta\mathbf{a}_B+\mathbf{R}[\delta\boldsymbol{\theta}]_\times\mathbf{a}_B+\underbrace{\mathbf{R}[\delta\boldsymbol{\theta}]_\times\delta\mathbf{a}_B}_{O(2)}+\mathbf{g}+\delta\mathbf{g}.
\]
消去 $\mathbf{R}\mathbf{a}_B+\mathbf{g}$、弃二阶项，并用叉积反交换 $[\delta\boldsymbol{\theta}]_\times\mathbf{a}_B=-[\mathbf{a}_B]_\times\delta\boldsymbol{\theta}$：
\[
  \dot{\delta\mathbf{v}}=\mathbf{R}\big(\delta\mathbf{a}_B-[\mathbf{a}_B]_\times\delta\boldsymbol{\theta}\big)+\delta\mathbf{g}
  =-\mathbf{R}[\mathbf{a}_m-\mathbf{a}_b]_\times\delta\boldsymbol{\theta}-\mathbf{R}\delta\mathbf{a}_b+\delta\mathbf{g}-\mathbf{R}\mathbf{a}_n.
\]
（若加速度计噪声白、各向同性 $\mathbb{E}[\mathbf{a}_n\mathbf{a}_n^\top]=\sigma_a^2\mathbf{I}$，则 $\mathbf{R}\mathbf{a}_n$ 与 $\mathbf{a}_n$ 同分布，可吸收记号写成 $-\mathbf{a}_n$。）
\end{proof}
\begin{proof}[角误差 \cref{eq:eskf-err-theta} 的逐步证明]
记体坐标大/小信号角速度 $\boldsymbol{\omega}\triangleq\boldsymbol{\omega}_m-\boldsymbol{\omega}_b$、$\delta\boldsymbol{\omega}\triangleq-\delta\boldsymbol{\omega}_b-\boldsymbol{\omega}_n$，$\boldsymbol{\omega}_t=\boldsymbol{\omega}+\delta\boldsymbol{\omega}$。对 $\mathbf{q}_t=\mathbf{q}\otimes\delta\mathbf{q}$ 求导并代入 $\dot{\mathbf{q}}_t=\tfrac12\mathbf{q}_t\otimes\boldsymbol{\omega}_t$、$\dot{\mathbf{q}}=\tfrac12\mathbf{q}\otimes\boldsymbol{\omega}$：
\[
  \tfrac12\mathbf{q}\otimes\boldsymbol{\omega}\otimes\delta\mathbf{q}+\mathbf{q}\otimes\dot{\delta\mathbf{q}}=\tfrac12\mathbf{q}\otimes\delta\mathbf{q}\otimes\boldsymbol{\omega}_t.
\]
左消 $\mathbf{q}$、解 $\dot{\delta\mathbf{q}}$，并用 $2\dot{\delta\mathbf{q}}=[0;\dot{\delta\boldsymbol{\theta}}]$、左右乘矩阵：
\[
  \begin{bmatrix}0\\\dot{\delta\boldsymbol{\theta}}\end{bmatrix}=\delta\mathbf{q}\otimes\boldsymbol{\omega}_t-\boldsymbol{\omega}\otimes\delta\mathbf{q}
  =\begin{bmatrix}-\delta\boldsymbol{\omega}^\top\delta\boldsymbol{\theta}/2\,+O(2)\\ \delta\boldsymbol{\omega}-[\boldsymbol{\omega}]_\times\delta\boldsymbol{\theta}-\tfrac12[\delta\boldsymbol{\omega}]_\times\delta\boldsymbol{\theta}\end{bmatrix}+O(\|\delta\boldsymbol{\theta}\|^2).
\]
标量行是二阶无穷小（无用）；向量行去二阶项得 $\dot{\delta\boldsymbol{\theta}}=-[\boldsymbol{\omega}]_\times\delta\boldsymbol{\theta}+\delta\boldsymbol{\omega}$，代回 $\boldsymbol{\omega},\delta\boldsymbol{\omega}$ 即 \cref{eq:eskf-err-theta}。其余 \cref{eq:eskf-err-p,eq:eskf-err-bias} 由对应真态/名义态相减平凡得到。
\end{proof}

\paragraph{离散预测：标称态非线性积分，误差态线性传协方差。}
标称态按 \cref{eq:eskf-nominal-kin} 离散积分（$\mathbf{p}\!\leftarrow\!\mathbf{p}+\mathbf{v}\Delta t+\tfrac12(\mathbf{R}(\mathbf{a}_m-\mathbf{a}_b)+\mathbf{g})\Delta t^2$，$\mathbf{v}\!\leftarrow\!\mathbf{v}+(\mathbf{R}(\mathbf{a}_m-\mathbf{a}_b)+\mathbf{g})\Delta t$，$\mathbf{q}\!\leftarrow\!\mathbf{q}\otimes\mathrm{Exp}((\boldsymbol{\omega}_m-\boldsymbol{\omega}_b)\Delta t)$，细节属 \cref{ch:imu}）。误差态满足线性化差分 $\delta\mathbf{x}\leftarrow\mathbf{F}_{\mathbf{x}}\,\delta\mathbf{x}+\mathbf{F}_{\mathbf{i}}\,\mathbf{i}$，其均值与协方差递推为\textbf{ESKF 预测方程}：
\begin{equation}
  \widehat{\delta\mathbf{x}}\leftarrow\mathbf{F}_{\mathbf{x}}\,\widehat{\delta\mathbf{x}},\qquad
  \mathbf{P}\leftarrow\mathbf{F}_{\mathbf{x}}\,\mathbf{P}\,\mathbf{F}_{\mathbf{x}}^\top+\mathbf{F}_{\mathbf{i}}\,\mathbf{Q}_{\mathbf{i}}\,\mathbf{F}_{\mathbf{i}}^\top,
  \label{eq:eskf-predict}
\end{equation}
其中误差态转移矩阵（按 $[\delta\mathbf{p};\delta\mathbf{v};\delta\boldsymbol{\theta};\delta\mathbf{a}_b;\delta\boldsymbol{\omega}_b;\delta\mathbf{g}]$ 排列；姿态块用\emph{闭式} $\mathrm{Exp}((\boldsymbol{\omega}_m-\boldsymbol{\omega}_b)\Delta t)^\top$ 而非 Euler 的 $\mathbf{I}-[\boldsymbol{\omega}]_\times\Delta t$）\cite{sola2017quaternion}：
\begin{equation}
  \mathbf{F}_{\mathbf{x}}=\begin{bmatrix}
  \mathbf{I}&\mathbf{I}\Delta t&\mathbf{0}&\mathbf{0}&\mathbf{0}&\mathbf{0}\\
  \mathbf{0}&\mathbf{I}&-\mathbf{R}[\mathbf{a}_m-\mathbf{a}_b]_\times\Delta t&-\mathbf{R}\Delta t&\mathbf{0}&\mathbf{I}\Delta t\\
  \mathbf{0}&\mathbf{0}&\mathrm{Exp}((\boldsymbol{\omega}_m-\boldsymbol{\omega}_b)\Delta t)^\top&\mathbf{0}&-\mathbf{I}\Delta t&\mathbf{0}\\
  \mathbf{0}&\mathbf{0}&\mathbf{0}&\mathbf{I}&\mathbf{0}&\mathbf{0}\\
  \mathbf{0}&\mathbf{0}&\mathbf{0}&\mathbf{0}&\mathbf{I}&\mathbf{0}\\
  \mathbf{0}&\mathbf{0}&\mathbf{0}&\mathbf{0}&\mathbf{0}&\mathbf{I}
  \end{bmatrix},
  \label{eq:eskf-Fx}
\end{equation}
扰动雅可比 $\mathbf{F}_{\mathbf{i}}$ 把噪声脉冲 $\mathbf{i}=[\mathbf{v}_i;\boldsymbol{\theta}_i;\mathbf{a}_i;\boldsymbol{\omega}_i]$ 映到 $\delta\mathbf{v},\delta\boldsymbol{\theta},\delta\mathbf{a}_b,\delta\boldsymbol{\omega}_b$ 四行（其余行为零），脉冲协方差 $\mathbf{Q}_{\mathbf{i}}=\mathrm{diag}(\mathbf{V}_i,\boldsymbol{\Theta}_i,\mathbf{A}_i,\boldsymbol{\Omega}_i)$，其中 $\mathbf{V}_i=\sigma_{a_n}^2\Delta t^2\mathbf{I}$、$\boldsymbol{\Theta}_i=\sigma_{\omega_n}^2\Delta t^2\mathbf{I}$、$\mathbf{A}_i=\sigma_{a_w}^2\Delta t\,\mathbf{I}$、$\boldsymbol{\Omega}_i=\sigma_{\omega_w}^2\Delta t\,\mathbf{I}$（由连续噪声密度在 $\Delta t$ 上积分而得）。
\textbf{注意}：误差均值初始化为零，\cref{eq:eskf-predict} 左式恒返回零、可在代码中跳过（仍保留以防遗漏）；但\textbf{协方差预测绝不可省}——$\mathbf{F}_{\mathbf{i}}\mathbf{Q}_{\mathbf{i}}\mathbf{F}_{\mathbf{i}}^\top$ 非零，使不确定度在预测步持续增长（应当如此）\cite{sola2017quaternion}。

\paragraph{校正：观测误差态。}
外部传感器（GPS/视觉）观测 $\mathbf{z}=\mathbf{h}(\mathbf{x}_t)+\mathbf{v}$，$\mathbf{v}\sim\mathcal{N}(\mathbf{0},\boldsymbol{\Sigma}_{\mathbf{v}})$。滤波器估计的是\emph{误差态}，用标准 KF 校正（此阶段误差均值为零，故评估点取名义态 $\mathbf{x}$）：
\begin{equation}
  \mathbf{K}=\mathbf{P}\mathbf{H}^\top(\mathbf{H}\mathbf{P}\mathbf{H}^\top+\boldsymbol{\Sigma}_{\mathbf{v}})^{-1},\quad
  \widehat{\delta\mathbf{x}}\leftarrow\mathbf{K}\big(\mathbf{z}-\mathbf{h}(\mathbf{x})\big),\quad
  \mathbf{P}\leftarrow(\mathbf{I}-\mathbf{K}\mathbf{H})\mathbf{P},
  \label{eq:eskf-correct}
\end{equation}
（协方差更新建议用 Joseph 形式，见 \cref{sec:kf} 数值陷阱）。关键是\emph{测量雅可比}须对误差态求，用链式法则
\begin{equation}
  \mathbf{H}=\left.\frac{\partial\mathbf{h}}{\partial\delta\mathbf{x}}\right|_{\mathbf{x}}
  =\underbrace{\left.\frac{\partial\mathbf{h}}{\partial\mathbf{x}_t}\right|_{\mathbf{x}}}_{\mathbf{H}_{\mathbf{x}}\ \text{标准 EKF 雅可比}}
  \underbrace{\frac{\partial\mathbf{x}_t}{\partial\delta\mathbf{x}}}_{\mathbf{X}_{\delta\mathbf{x}}\ \text{合成雅可比}}.
  \label{eq:eskf-H}
\end{equation}
$\mathbf{H}_{\mathbf{x}}$ 是 $\mathbf{h}$ 对真态的常规雅可比（依具体传感器，与普通 EKF 一致）；$\mathbf{X}_{\delta\mathbf{x}}=\partial\mathbf{x}_t/\partial\delta\mathbf{x}$ 只依赖状态合成方式：向量块全是单位阵，唯独四元数块（右扰动版）为
\begin{equation}
  \mathbf{Q}_{\delta\boldsymbol{\theta}}=\left.\frac{\partial(\mathbf{q}\otimes\delta\mathbf{q})}{\partial\delta\boldsymbol{\theta}}\right|_{\mathbf{q}}
  =[\mathbf{q}]_L\,\tfrac12\begin{bmatrix}\mathbf{0}^\top\\\mathbf{I}_3\end{bmatrix}
  =\tfrac12\begin{bmatrix}-q_x&-q_y&-q_z\\ q_w&-q_z&q_y\\ q_z&q_w&-q_x\\ -q_y&q_x&q_w\end{bmatrix},
  \label{eq:eskf-Qdtheta}
\end{equation}
其中用了 $\delta\mathbf{q}\xrightarrow{\delta\boldsymbol{\theta}\to\mathbf{0}}[1;\tfrac12\delta\boldsymbol{\theta}]$ 与左乘矩阵 $[\mathbf{q}]_L$（\cref{ch:lie}/四元数）。于是 $\mathbf{X}_{\delta\mathbf{x}}=\mathrm{diag}(\mathbf{I}_6,\mathbf{Q}_{\delta\boldsymbol{\theta}},\mathbf{I}_9)$。

\paragraph{注入与复位（ESKF 特有）。}
解出误差态后，\emph{注入}（inject）回标称态，再把误差态\emph{复位}（reset）：
\begin{align}
  \text{注入：}\quad
  &\mathbf{x}\leftarrow\mathbf{x}\oplus\widehat{\delta\mathbf{x}},\qquad
  \mathbf{R}\leftarrow\mathbf{R}\,\mathrm{Exp}(\widehat{\delta\boldsymbol{\theta}})\ \ (\text{姿态右乘},\ \mathbf{q}\leftarrow\mathbf{q}\otimes\mathrm{Exp}(\widehat{\delta\boldsymbol{\theta}})); \label{eq:eskf-inject}\\
  \text{复位：}\quad
  &\widehat{\delta\mathbf{x}}\leftarrow\mathbf{0},\qquad
  \mathbf{P}\leftarrow\mathbf{G}\,\mathbf{P}\,\mathbf{G}^\top, \label{eq:eskf-reset}
\end{align}
向量分量直接相加。复位雅可比 $\mathbf{G}=\mathrm{diag}(\mathbf{I}_6,\mathbf{I}-\tfrac12[\widehat{\delta\boldsymbol{\theta}}]_\times,\mathbf{I}_9)$，多数实现取 $\widehat{\delta\boldsymbol{\theta}}\approx\mathbf{0}\Rightarrow\mathbf{G}\approx\mathbf{I}$（平凡复位），保留姿态修正项可降低长期漂移\cite{sola2017quaternion}。注入用\emph{右乘} $\mathrm{Exp}(\widehat{\delta\boldsymbol{\theta}})$ 正是本书右扰动约定的落地。

\begin{algo}[ESKF 一个周期]\label{algo:eskf}
\textbf{预测}（每个高频惯性输入）：(1) 标称态按 \cref{eq:eskf-nominal-kin} 非线性积分；(2) 误差协方差 \cref{eq:eskf-predict}。\\
\textbf{校正}（每个低频外部观测）：(3) 算 $\mathbf{H}=\mathbf{H}_{\mathbf{x}}\mathbf{X}_{\delta\mathbf{x}}$（\cref{eq:eskf-H,eq:eskf-Qdtheta}）与增益 $\mathbf{K}$；(4) 由新息 $\mathbf{z}-\mathbf{h}(\mathbf{x})$ 解 $\widehat{\delta\mathbf{x}}$、更新 $\mathbf{P}$（\cref{eq:eskf-correct}，Joseph 形式）。\\
\textbf{注入 $+$ 复位}：(5) $\mathbf{x}\leftarrow\mathbf{x}\oplus\widehat{\delta\mathbf{x}}$（\cref{eq:eskf-inject}）；⑥ $\widehat{\delta\mathbf{x}}\leftarrow\mathbf{0}$，$\mathbf{P}\leftarrow\mathbf{G}\mathbf{P}\mathbf{G}^\top$（\cref{eq:eskf-reset}）。
\end{algo}

\paragraph{右扰动（局部）vs 左扰动（全局）。}
姿态误差可定义在\emph{机体系}（右乘，$\mathbf{R}_t=\mathbf{R}\,\mathrm{Exp}(\delta\boldsymbol{\theta})$，本书主线）或\emph{世界系}（左乘，$\mathbf{R}_t=\mathrm{Exp}(\delta\boldsymbol{\theta})\,\mathbf{R}$）。Solà\cite{sola2017quaternion} §7 指出全局角误差（对接 Li–Mourikis VIO）某些性质更好。两者只改动几处雅可比与注入/复位的乘法侧，列于 \cref{tab:eskf-pert}（角速度始终在机体系，与误差定义无关）。全局版的关键改变：角误差动力学变为 $\dot{\delta\boldsymbol{\theta}}=-\mathbf{R}\delta\boldsymbol{\omega}_b-\mathbf{R}\boldsymbol{\omega}_n$（无 $-[\boldsymbol{\omega}]_\times\delta\boldsymbol{\theta}$ 项）。

\begin{table}[H]\centering\small
\caption{ESKF 角误差：右扰动（局部/机体系）vs 左扰动（全局/世界系）三处差异（其余完全相同）}
\label{tab:eskf-pert}
\begin{tabular}{lll}
\toprule
环节 & 右扰动（本书主线，局部） & 左扰动（全局） \\
\midrule
误差合成 & $\mathbf{q}_t=\mathbf{q}\otimes\delta\mathbf{q}$ & $\mathbf{q}_t=\delta\mathbf{q}\otimes\mathbf{q}$ \\
$\partial\delta\mathbf{v}^+/\partial\delta\boldsymbol{\theta}$ & $-\mathbf{R}[\mathbf{a}_m-\mathbf{a}_b]_\times\Delta t$ & $-[\mathbf{R}(\mathbf{a}_m-\mathbf{a}_b)]_\times\Delta t$ \\
姿态转移块 $\partial\delta\boldsymbol{\theta}^+/\partial\delta\boldsymbol{\theta}$ & $\mathrm{Exp}((\boldsymbol{\omega}_m-\boldsymbol{\omega}_b)\Delta t)^\top$ & $\mathbf{I}$ \\
$\partial\delta\boldsymbol{\theta}^+/\partial\delta\boldsymbol{\omega}_b$ & $-\mathbf{I}\Delta t$ & $-\mathbf{R}\Delta t$ \\
注入 & $\mathbf{q}\leftarrow\mathbf{q}\otimes\mathrm{Exp}(\widehat{\delta\boldsymbol{\theta}})$ & $\mathbf{q}\leftarrow\mathrm{Exp}(\widehat{\delta\boldsymbol{\theta}})\otimes\mathbf{q}$ \\
复位姿态块 $\mathbf{G}_{\theta}$ & $\mathbf{I}-\tfrac12[\widehat{\delta\boldsymbol{\theta}}]_\times$ & $\mathbf{I}+\tfrac12[\widehat{\delta\boldsymbol{\theta}}]_\times$ \\
\bottomrule
\end{tabular}
\end{table}

\begin{pitfall}[把姿态当普通向量做 EKF]
直接把四元数 $4$ 维或欧拉角放进 KF 状态向量并线性化，会遇到：四元数协方差奇异、需在每步强行归一化（破坏一致性）；欧拉角万向锁。ESKF 用最小 $3$ 参数误差 $\delta\boldsymbol{\theta}$ 滤波、再注入回流形，从根上回避。这与 \cref{ch:lie} “优化位姿要在李代数/扰动上做、而非直接对矩阵元素”是同一原则。
\end{pitfall}

\paragraph{ESKF vs 直接 EKF vs UKF。}
直接（流形）EKF 把姿态放进状态，有上述过参数化/奇异问题；ESKF 多一步注入 $+$ 复位、测量雅可比多乘一个 $\mathbf{X}_{\delta\mathbf{x}}$，换来协方差规整、线性化恒有效（误差恒在原点附近）、雅可比近常数，在“高频预测 $+$ 低频校正”的惯性融合中最划算。UKF 免雅可比、对强非线性更准，但 sigma 点数随维度增长、流形上需特殊处理，维度可控且难求雅可比时才更合适。

\begin{insight}
\textbf{ESKF 一致性优于直接 EKF——回收 \cref{sec:consistency} 的语言}。EKF 不一致的头号根源是“工作点取在估计均值而非真值，强非线性下线性化失效”（\cref{sec:ekf} 三宗罪之一）。ESKF 的\emph{误差态恒被复位到原点附近}（\cref{eq:eskf-reset}），故每次线性化都发生在 $\delta\mathbf{x}\approx\mathbf{0}$——线性化误差被系统性压到最小，雅可比 \cref{eq:eskf-Fx} 也近常数。用 \cref{def:nees-nis} 的话说：ESKF 的 NEES 通常比直接 EKF 更贴近 $N$（不那么过自信）。但这\emph{并未}根除不一致——姿态/速度/位置间的可观测性耦合仍可能让线性化误差与状态相关，从而虚假地“增加信息”、压乐观协方差。彻底从对称性上解决这一点，是下面\emph{不变卡尔曼滤波}的主题。
\end{insight}

\paragraph{轻量替代：互补滤波（Mahony / Madgwick）。}
ESKF 之外，姿态估计还有一支\emph{确定性、无协方差}的工程方案——\emph{互补滤波}（complementary filter），常作“为什么需要概率滤波”的反面对照。它的出发点是频域互补：陀螺积分给出\emph{短期准、长期漂移}的姿态（高频可信），加速度计/磁力计给出\emph{长期准（重力/地磁方向）、短期受运动加速度污染}的姿态（低频可信）。互补滤波用一个\emph{高通}滤陀螺、一个\emph{低通}滤加速度计，两者相加（传递函数互补为 $1$），取长补短\cite{mahony2008nonlinear}。Mahony 等的\emph{非线性互补滤波}把它\emph{直接定义在} $\mathrm{SO}(3)$ 上（避开欧拉角奇异），用观测向量与其估计方向的\emph{叉积}构造姿态误差、反馈修正陀螺积分，并在线估计陀螺偏置——形式上是一个 PI 控制律：
\begin{equation}
  \dot{\hat{\mathbf{R}}}=\hat{\mathbf{R}}\big(\boldsymbol{\omega}_m-\hat{\mathbf{b}}+k_P\,\boldsymbol{\omega}_{\text{corr}}\big)^\wedge,\qquad
  \dot{\hat{\mathbf{b}}}=-k_I\,\boldsymbol{\omega}_{\text{corr}},\qquad
  \boldsymbol{\omega}_{\text{corr}}=\textstyle\sum_i k_i\,\hat{\mathbf{v}}_i\times(\hat{\mathbf{R}}^\top\mathbf{v}_i^{\text{ref}}),
  \label{eq:complementary-filter}
\end{equation}
其中 $\mathbf{v}_i^{\text{ref}}$ 是世界系参考向量（重力、地磁），$\hat{\mathbf{v}}_i$ 是机体系测量方向，叉积给出把估计姿态\emph{转向}观测的修正角速度；$k_P,k_I$ 为比例/积分增益。Madgwick 滤波\cite{madgwick2011estimation} 是同一思想的梯度下降变体（用四元数表示、以加速度计/磁力计残差的梯度方向修正），以低算力、易调参著称，广泛用于开源 IMU 姿态估计。

\begin{insight}
\textbf{互补滤波 vs ESKF：何时用哪个}。互补滤波\emph{不传播协方差}、增益 $k_P,k_I$ 靠经验整定（而非由噪声统计导出），故\textbf{无法报告不确定度、无法做一致性检验、无法最优融合异质传感器}——但它\emph{极轻量}（几次叉积与积分）、无矩阵求逆、数值稳健，适合超低算力的 IMU-only 姿态估计（小型无人机、可穿戴）。ESKF 则给出完整协方差、可与 GPS/视觉等任意观测按精度最优融合、可联动一致性诊断，代价是要维护并传播 $\mathbf{P}$。一句话：\emph{只要姿态、且算力极紧} $\to$ 互补滤波；\emph{要协方差或要融合多源} $\to$ ESKF。\cref{ch:imu} 会在 IMU 标定语境再用互补滤波做初始对准。
\end{insight}

\paragraph{现代去向：不变卡尔曼滤波（InEKF）与一致性。}
ESKF 把误差态压在原点附近、改善了一致性，但\emph{并未}从根上消除“线性化误差与状态相关”这一不一致之源。过去十年滤波侧最重要的进展之一——\emph{不变（扩展）卡尔曼滤波}（Invariant EKF, InEKF；Barrau--Bonnabel\cite{barrau2017invariant}，Barfoot 第 2 版新增材料）——用李群\emph{对称性}给出了更彻底的方案，且与本书\emph{右扰动}主线天然契合。

\begin{paper}[不变卡尔曼滤波：群仿射系统与误差自治]\label{paper:inekf}
\textbf{核心思想}\cite{barrau2017invariant}：当状态活在李群 $G$ 上、且其确定性动力学满足\emph{群仿射}（group-affine）结构 $f(\mathbf{X}_1\mathbf{X}_2)=f(\mathbf{X}_1)\mathbf{X}_2+\mathbf{X}_1 f(\mathbf{X}_2)-\mathbf{X}_1 f(\mathbf{I})\mathbf{X}_2$ 时，定义\emph{右不变误差} $\boldsymbol{\eta}=\hat{\mathbf{X}}\mathbf{X}^{-1}$（或左不变 $\mathbf{X}^{-1}\hat{\mathbf{X}}$），其误差动力学有一条惊人性质：\emph{精确}（未线性化）的误差动力学就已是\textbf{轨迹自治}的（autonomous error，其右端只含群元素本身、不显含状态轨迹）；由此推论，该精确误差经 $\mathrm{Log}$ 后的\emph{线性化}传播矩阵也不依赖当前估计 $\hat{\mathbf{X}}$（log-linear property）。是“精确自治”这一更强的结构，保证了线性化传播与状态轨迹无关。

\textbf{为什么这治好了不一致}：EKF/ESKF 不一致的根源正是“线性化雅可比依赖估计状态”，使得估计误差与真值的耦合方式随轨迹漂移、虚假地“制造可观测性”、把协方差压得过乐观（\cref{thm:bias-overconfident} 是其偏置版，可观测性退化是其几何版）。群仿射 $+$ 不变误差让传播\emph{与状态解耦}，于是线性化在更大范围内精确、可观测性结构不被线性化扭曲，\emph{一致性显著改善}——这是 FEJ（首次估计雅可比）等启发式之上更具原理性的解法。

\textbf{需要的状态空间——扩展位姿群 $\mathrm{SE}_2(3)$}：把惯导的\emph{旋转、位置、速度}打包成一个 $5\times5$ 矩阵李群元素
\[
  \mathbf{X}=\begin{bmatrix}\mathbf{R}&\mathbf{r}&\mathbf{v}\\\mathbf{0}^\top&1&0\\\mathbf{0}^\top&0&1\end{bmatrix}\in\mathrm{SE}_2(3),\qquad
  \mathbf{R}\in\mathrm{SO}(3),\ \mathbf{r},\mathbf{v}\in\mathbb{R}^3,
\]
（\textbf{记号}：列序与位置/速度字母\emph{完全遵循}权威定义 \cref{def:se23}——第 2 列位置 $\mathbf{r}$、第 3 列速度 $\mathbf{v}$；此处的位置 $\mathbf{r}$ 即本章 ESKF 误差态 \cref{eq:eskf-Fx} 里的位置量 $\mathbf{p}$，仅因群上沿用 \cref{ch:lie} 的 $\mathbf{r}$ 记法而换字母，二者同指机体在世界系的位置。）正是在这个“双平移”扩展特殊欧氏群上，IMU 运动学才呈现群仿射结构（普通 $\mathrm{SE}(3)$ 只含旋转 $+$ 位置，装不下速度，不满足群仿射）。$\mathrm{SE}_2(3)$ 的\emph{正式}李群定义（指数/对数映射、伴随、左右雅可比）见 \cref{sec:se23}；本章只用到它“$5\times5$、含 $\mathbf{R}/\mathbf{r}/\mathbf{v}$、右扰动 $\mathbf{X}\,\mathrm{Exp}(\delta\boldsymbol{\xi})$”这一结构要点。\textbf{一处实务限制}：把陀螺/加计\emph{零偏}增广进状态会\emph{破坏}群仿射结构（偏置项不满足上面的 $f(\mathbf{X}_1\mathbf{X}_2)$ 恒等式），故工程上用\emph{不完美不变滤波}（imperfect InEKF）——把偏置当作群仿射动力学的扰动项近似处理；这也是为何 InEKF 并非“无损地把偏置吞进群”。

\textbf{去向}：InEKF 已成当代 VIO/惯导一致性研究的主流框架之一（\cref{ch:vio} 的“右不变卡尔曼滤波”去向、足式机器人的接触辅助惯导 InEKF 等）。它与本章 ESKF 的关系是：\emph{同样在流形上滤波、同样右扰动}，但 InEKF 进一步选对了\emph{群结构与误差定义}，使线性化误差与状态无关，从而把“误差恒在原点附近”升级为“误差传播也与轨迹无关”。
\end{paper}

\begin{practice}
\begin{enumerate}
  \item 用 \cref{ch:lie} 的右扰动求 $\partial(\mathbf{R}\mathbf{p})/\partial\delta\boldsymbol{\theta}=-\mathbf{R}[\mathbf{p}]_\times$（提示 $\mathbf{R}\,\mathrm{Exp}(\delta\boldsymbol{\theta})\mathbf{p}\approx\mathbf{R}\mathbf{p}+\mathbf{R}[\delta\boldsymbol{\theta}]_\times\mathbf{p}$），把它作为某视觉观测 $\mathbf{H}_{\mathbf{x}}$ 的姿态块。
  \item 独立复现 \cref{thm:eskf-error-kin} 的角误差证明（\cref{eq:eskf-err-theta}），并指出哪一步丢了二阶项 $\tfrac12[\delta\boldsymbol{\omega}]_\times\delta\boldsymbol{\theta}$、为什么可以丢。
  \item 依 \cref{tab:eskf-pert} 把一个右扰动 ESKF 改写为左扰动版，指出 $\mathbf{F}_{\mathbf{x}}$ 中哪三块变、其余为何不变。答案要点：仅姿态相关三块变（$\partial\delta\mathbf{v}/\partial\delta\boldsymbol{\theta}$、姿态转移块、$\partial\delta\boldsymbol{\theta}/\partial\delta\boldsymbol{\omega}_b$），位置/速度/偏置/重力块与误差定义无关。答错回看 \cref{tab:eskf-pert}。
\end{enumerate}
\end{practice}

% ════════════════════════════════════════════════════════════════
\section{实践：滤波还是优化？兼论鲁棒化}
\label{sec:filter-vs-opt}

ESKF 等滤波器在嵌入式、高频惯性融合上仍是主力；但现代视觉 SLAM 后端的主流是 BA/因子图\emph{优化}（平滑）。本节给工程抉择依据，并补上让两类方法都不可或缺的\emph{鲁棒化}（离群点处理）。

\paragraph{EKF-SLAM 的四条局限（转向优化的论据）。}\cite{gaoxiang2019slam14}
\begin{enumerate}
  \item \textbf{马尔可夫假设难处理回环}：滤波只用邻近时刻，当前帧与很久前数据相关（回环）时难以利用；优化（Full-SLAM）用全部历史，回环天然纳入。
  \item \textbf{只线性化一次 $\Rightarrow$ 不一致}：EKF 在固定点线性化一次就定后验；优化每次迭代在新估计点\emph{重新}线性化，状态变化大时更稳。“EKF 约等于优化的一次迭代”（\cref{prop:ekf} 与 \cref{eq:iekf} 已坐实）。
  \item \textbf{协方差随状态平方增长}：把路标纳入状态时，EKF 要维护稠密协方差，存储/计算 $O(n^2)$，且边缘化历史会让信息矩阵\emph{稠密化}（fill-in），不适合大场景；优化（平滑）的信息矩阵\emph{保持稀疏}（块三对角 / 箭头形），可扩展到百万级\cite{carlone2026handbook}。
  \item \textbf{无异常检测易发散}：滤波无内建外点剔除，而视觉匹配/光流常有误匹配；优化可加鲁棒核函数（见下）。
\end{enumerate}

\paragraph{为什么优化保稀疏、滤波却稠密——共视的代价。}
回顾 \cref{ch:nlopt}/\cref{ch:slam_est}：BA 的信息矩阵 $\mathbf{H}=\mathbf{J}^\top\mathbf{J}$ 因“每条观测只连一个相机一个路标”而\emph{块稀疏}，按相机/路标分块为 $\big[\begin{smallmatrix}\mathbf{B}&\mathbf{E}\\\mathbf{E}^\top&\mathbf{C}\end{smallmatrix}\big]$（$\mathbf{B},\mathbf{C}$ 对角块），用舒尔补 $\mathbf{S}=\mathbf{B}-\mathbf{E}\mathbf{C}^{-1}\mathbf{E}^\top$ 边缘化路标得位姿子问题——$\mathbf{S}$ 的非零块恰对应两相机有\emph{共视}路标。而滤波\emph{每步都把过去状态边缘化掉}，边缘化会在被消去变量的邻居间引入新边（fill-in），信息矩阵逐渐稠密化。\textbf{对偶要点}：批量/平滑保留全部变量、只在求解时利用稀疏；滤波强制边缘化、以稠密化换常数维度。这就是“同等算力下优化更优”的结构性原因。\cref{par:swf} 的 SWF 正是用“窗口内不消、只消出窗者”在两者间取平衡。

\paragraph{滤波 vs 平滑：一张对照表。}
\begin{table}[H]\centering\small
\caption{滤波（filtering）与优化/平滑（smoothing）的工程对照}
\label{tab:filter-vs-opt}
\begin{tabular}{lll}
\toprule
维度 & 滤波（KF/EKF/ESKF） & 优化 / 平滑（BA/因子图） \\
\midrule
时间范围 & 马尔可夫，邻近时刻 & Full-SLAM，全部历史 \\
线性化 & 固定点一次 & 每次迭代重线性化 \\
回环 & 难 & 天然纳入 \\
信息矩阵 & 边缘化历史 $\Rightarrow$ 稠密（fill-in） & 保持稀疏、可扩展 \\
外点 & 无内建机制 & 鲁棒核 / RANSAC \\
计算 & 高频、轻量、在线、常数时间 & 较重、可后台慢跑、非常数时间 \\
典型场合 & 惯性融合、嵌入式、高频 & 视觉 SLAM 后端、离线建图 \\
\bottomrule
\end{tabular}
\end{table}

\noindent 理论与实践比较都表明：\emph{同等计算量}下，非线性优化在精度与鲁棒性上通常优于滤波\cite{gaoxiang2019slam14,carlone2026handbook}。这正是上一章把 SLAM 钉成“因子图上的 MAP”、并留待后续 \cref{ch:nlopt} 专门求解的原因。但这\emph{不}意味着滤波过时——

\begin{insight}
\textbf{何时仍选滤波}：算力/内存受限的嵌入式平台；状态维度低、模型简单；需要\emph{高频}输出（如 IMU 200--1000\,Hz 的 ESKF 预测）。常见的现代架构是\emph{混合}：用 ESKF/滤波做高频前端传播，用因子图优化做低频后端精修——自 PTAM 起，后端优化就被认为“无须实时响应每帧”，故与前端分线程运行\cite{gaoxiang2019slam14}。VIO（\cref{ch:vio}）正是这种分工的典型。选型决策语：\emph{若}需高频且算力紧 $\to$ ESKF；\emph{若}要回环 + 大地图 + 高精度 $\to$ 因子图优化；\emph{若}两者都要 $\to$ ESKF 前端 + 优化后端混合。
\end{insight}

\paragraph{先有正确对应：数据关联（外部 vs 内部）。}
离群点处理的上游是\emph{数据关联}（data association）——弄清哪条测量对应模型/地图里的哪个元素\cite{barfoot2024state}。它分两类：\emph{外部关联}（external，又称“已知关联”）用估计问题之外的专门信息建立对应，如 GPS 卫星在时戳消息里嵌唯一码、视觉条形码、给目标涂唯一颜色——把对应在工程上事先做掉，估计问题随之简化；\emph{内部关联}（internal，又称“未知关联”）只能用测量与模型本身做对应，常基于“给定模型时该测量的似然”，最简单是取\emph{最可能}的关联，复杂者把多关联假设一并带入估计。\emph{两个对照范例把这条分界讲透}\cite{barfoot2024state}：GPS 定位中，卫星轨道作为时间的函数已知、测距靠时戳消息里嵌的\emph{唯一码}识别来自哪颗卫星，对应被工程\emph{事先做掉}（外部关联，随后 trilateration 解出位置即可），是“易关联”的样板；而\emph{星敏定姿}（star tracker）以预先编绘的\emph{星图}（用自然界作地图、事先测绘）为参照、传感器只看到一堆亮点，“这个亮点对应星图里哪颗星”必须\emph{在估计内部}用观测本身判定（内部关联），远比 GPS 难——它正是“内部 / 难关联”的经典动机范例。对这类三维地标/星表，常用地标\emph{星座}（constellation）辅助：点对间距离可作唯一标识，如 DARCES（数据对齐刚性约束穷举搜索，Chen 1999\cite{chen1999ransac}）。\textbf{一条工程铁律}：估计器一旦发散，头号嫌疑往往是坏的数据关联，而非滤波公式本身——因此必须承认误关联会发生，并用下面的离群点机制把它挡在估计之外。

\paragraph{鲁棒化：离群点为什么必须专门处理。}
即便关联正确，仍有测量本身很差（GPS 多径反射给出过长伪距、视线被遮挡时接收机无从分辨）。\emph{离群点}（outlier）即按测量模型判断极不可能的测量（一维常以偏离均值 $>3\sigma$ 为界）。无论滤波还是优化，二范数平方代价对离群点极敏感：单个大误差项主导二次代价、梯度巨大，会把解“拽”过去以迎合一条本不该存在的边\cite{gaoxiang2019slam14}。两条经典对策（可单用或合用）\cite{barfoot2024state}：\emph{RANSAC}（预处理，随机抽最小子集拟合、按内点数投票）与\emph{M-估计/鲁棒核}（改代价形状）。RANSAC 单次迭代五步（以直线拟合为例）：(1) 随机取最小子集作\emph{假设内点}（拟合直线取两点）；(2) 对假设内点拟合模型；(3) 用其余数据按残差判内点/离群，内点太少则本次作废；(4) 用\emph{假设内点 $+$ 已判内点}重拟合；(5) 以全部内点残差评估。多次迭代取残差最小者。所需迭代次数由“至少一次抽到全内点子集”的概率 $1-p=(1-w^n)^k$ 反解：
\begin{equation}
  k=\frac{\ln(1-p)}{\ln(1-w^n)},
  \label{eq:eskf-ransac}
\end{equation}
$p$ 目标置信、$w$ 单点内点率、$n$ 子集大小；现实中数据多为顺序选取（非独立），\cref{eq:eskf-ransac} 是\emph{上界}。

\paragraph{鲁棒核与 IRLS。}
把二范数 $\rho(u)=\tfrac12u^2$（$u_i=\sqrt{\mathbf{e}_i^\top\boldsymbol{\Sigma}_i^{-1}\mathbf{e}_i}$）换成增长更慢、仍光滑的 $\rho(u)$，即\emph{鲁棒核}\cite{gaoxiang2019slam14,barfoot2024state}（此处只列滤波侧最常用的三款，与 \cref{ch:nlopt} 的鲁棒核\emph{是同一套}——\cref{sec:nlopt-robust} 的 \cref{tab:robust} 给出含 Tukey 在内的全家福与完整 M-估计理论，本节字母 $\boldsymbol{\Sigma}_i,\mathbf{Y}_i$ 亦与之对齐）：
\begin{equation}
  \rho_{\text{Huber}}(u)=\begin{cases}\tfrac12u^2,&|u|\le\delta\\ \delta(|u|-\tfrac12\delta),&|u|>\delta\end{cases},\quad
  \rho_{\text{Cauchy}}(u)=\tfrac{c^2}{2}\ln\!\Big(1+\tfrac{u^2}{c^2}\Big),\quad
  \rho_{\text{GM}}(u)=\tfrac12\frac{u^2}{1+u^2}.
  \label{eq:robust-kernels}
\end{equation}
（Cauchy 写成带尺度参数 $c$ 的形式，与 \cref{tab:robust} 一致；取 $c=1$ 即简化版 $\tfrac12\ln(1+u^2)$。）Huber 在 $|u|>\delta$ 后由二次转一次（限梯度上界、保光滑）；Cauchy、Geman–McClure（GM）增长更慢。对鲁棒代价 $J'=\sum_i\alpha_i\rho(u_i(\mathbf{x}))$ 求梯度并整理，可写回“加权最小二乘”的形式，只是协方差被\emph{膨胀}：
\begin{equation}
  \frac{\partial J'}{\partial\mathbf{x}}=\sum_i\mathbf{e}_i^\top\mathbf{Y}_i^{-1}\frac{\partial\mathbf{e}_i}{\partial\mathbf{x}},\qquad
  \mathbf{Y}_i^{-1}=\frac{\alpha_i}{u_i}\frac{\partial\rho}{\partial u_i}\boldsymbol{\Sigma}_i^{-1}.
  \label{eq:irls-weight}
\end{equation}
在上一迭代点 $\mathbf{x}_{\text{op}}$ 处冻结 $\mathbf{Y}_i(\mathbf{x}_{\text{op}})$、每步解一个普通加权最小二乘，即\emph{迭代重加权最小二乘}（IRLS）：收敛时梯度与原鲁棒问题相同、故同极小，只是路径不同。以 Cauchy（取 $c=1$）为例 $\mathbf{Y}_i(\mathbf{x}_{\text{op}})=\tfrac1{\alpha_i}(1+\mathbf{e}_i^\top\boldsymbol{\Sigma}_i^{-1}\mathbf{e}_i)\boldsymbol{\Sigma}_i$——正是“误差越大、信任越低”的协方差膨胀（与 \cref{ch:nlopt} 的 \cref{eq:nlopt-cauchy-Y} 同一结果，$\mathbf{Y}_i$ 记法一致）。

\begin{theorem}[（选读）鲁棒核 $=$ 协方差逆 Wishart 先验下的 MAP]\label{thm:robust-iw}
把每项协方差 $\mathbf{M}_i$ 也当未知量、加\emph{逆 Wishart} 先验 $\mathbf{M}_i\sim\mathcal{W}^{-1}(\boldsymbol{\Psi}_i,\nu_i)$ 做联合 MAP，对 $\mathbf{M}_i^{-1}$ 求导置零得最优协方差
\begin{equation}
  \mathbf{M}_i(\mathbf{x})=\tfrac{1}{\alpha_i}\boldsymbol{\Psi}_i+\tfrac{1}{\alpha_i}\mathbf{e}_i(\mathbf{x})\mathbf{e}_i(\mathbf{x})^\top\qquad(\alpha_i=\nu_i+M_i+2),
  \label{eq:iw-cov}
\end{equation}
代回消去 $\mathbf{M}$，目标函数化为 $J'(\mathbf{x})=\tfrac12\sum_i\alpha_i\ln(1+\mathbf{e}_i^\top\boldsymbol{\Psi}_i^{-1}\mathbf{e}_i)$——当 $\boldsymbol{\Psi}_i=\boldsymbol{\Sigma}_i$ 时\textbf{恰是加权 Cauchy 鲁棒代价}。（这是与 \cref{ch:nlopt} 共享的结论；含 Sherman--Morrison 的\emph{完整}逐步证明见 \cref{sec:nlopt-robust} 的 \cref{eq:nlopt-robust-cauchy}，本处只给要点。）
\end{theorem}
\begin{proof}[要点]
逆 Wishart 负对数似然含 $\tfrac12(\mathbf{e}_i^\top\mathbf{M}_i^{-1}\mathbf{e}_i-\alpha_i\ln\det\mathbf{M}_i^{-1}+\operatorname{tr}(\boldsymbol{\Psi}_i\mathbf{M}_i^{-1}))$，对 $\mathbf{M}_i^{-1}$ 用 $\partial\ln\det\mathbf{M}^{-1}/\partial\mathbf{M}^{-1}=\mathbf{M}$、$\partial\operatorname{tr}(\boldsymbol{\Psi}\mathbf{M}^{-1})/\partial\mathbf{M}^{-1}=\boldsymbol{\Psi}$ 求导置零即得 \cref{eq:iw-cov}（含一个“常数 $+$ 残差膨胀”结构，与 IRLS 的 $\mathbf{Y}_i$ 同形）；代回并丢弃与 $\mathbf{x}$ 无关项得 Cauchy 型代价\cite{barfoot2024state}。
\end{proof}

\begin{insight}
\textbf{本质洞察}：鲁棒核不是“临时补丁”，而是对“噪声协方差本身有不确定性”这一事实的 MAP 解——\textbf{所有鲁棒代价都可看成对协方差矩阵的某种先验分布的特定选择}\cite{barfoot2024state}。这把“离群点处理”从工程技巧提升为概率推断，与本章“一切估计都是 MAP/贝叶斯”的主线一以贯之。
\end{insight}

\paragraph{选哪个核、怎么避免坏局部极小（GNC）。}
鲁棒代价虽抑制离群点，却把目标变\emph{非凸}，IRLS 易陷入差的局部极小。两条现代缓解\cite{barfoot2024state}：(1) \emph{自适应核族}——Barron（arXiv:1701.03077）证明 Huber/Cauchy/Geman--McClure/Welsch 等可统一为单参数族 $\rho(u;\beta)$，在估计中\emph{自适应}地学出形状参数 $\beta$，省去手调核宽\cite{barron2019general}；(2) \emph{逐步非凸化}（graduated non-convexity, GNC，Yang 2020\cite{yang2020gnc}）——先用接近二次（凸）的代价求一个好初值，再沿一个连续参数把代价\emph{逐步}收紧到目标鲁棒核，使解平滑地从凸问题过渡到非凸问题，显著降低落入坏极小的风险。工程上 GNC 已成为位姿图/点云配准对抗高比例离群点的标配（\cref{ch:lidar_slam}、\cref{ch:vio} 会再用）。

\begin{practice}
\begin{enumerate}
  \item 给定一个算力受限的轮式机器人、低维状态、需 100\,Hz 输出，论证选 ESKF 还是因子图优化。答案要点：高频 + 低维 + 算力紧 $\to$ ESKF（\cref{tab:filter-vs-opt} 末行）。
  \item 解释“边缘化历史使信息矩阵稠密（fill-in）”为何让滤波在大地图上变慢，对照 \cref{ch:slam_est} 的稀疏性讨论与本节共视段。
  \item （证明骨架）由 \cref{thm:robust-iw} 出发，验证 $\boldsymbol{\Psi}_i=\boldsymbol{\Sigma}_i$ 时消去 $\mathbf{M}$ 的目标 $\tfrac12\alpha_i\ln(1+\mathbf{e}_i^\top\boldsymbol{\Sigma}_i^{-1}\mathbf{e}_i)$ 的梯度，与 \cref{eq:irls-weight} 的 Cauchy $\mathbf{Y}_i$ 一致。答错回看 \cref{eq:irls-weight} 与 \cref{thm:robust-iw}。
  \item 由 \cref{eq:eskf-ransac}：内点率 $w=0.5$、子集 $n=2$、目标置信 $p=0.99$，估计 RANSAC 至少需多少次迭代？再讨论为什么内点率从 $0.5$ 降到 $0.3$ 时迭代数激增。答案要点：$k=\ln(0.01)/\ln(1-0.25)\approx16$；$w^n$ 随 $w$ 下降使 $\ln(1-w^n)\to0^-$，$k$ 发散。
  \item （概念）解释“估计器发散时头号嫌疑是坏数据关联”，并区分外部关联与内部关联各自适合什么场景。答案要点：外部关联靠唯一码/颜色等问题外信息（GPS、条形码），内部关联只用测量似然（视觉特征匹配）；答错回看本节数据关联段。
  \item （实现，选做）用 Eigen 实现一维 KF（\cref{thm:kf}），打印每步 $\hat{x}_k,\hat{P}_k$，观察预测–校正下协方差的涨落（对照 \cref{ex:kf-1d}）。
\end{enumerate}
\end{practice}

\begin{codebox}[C++]{一维卡尔曼滤波最小实现（验证 thm:kf 与 ex:kf-1d）}
#include <iostream>
// 标量一维 KF: x_k = x_{k-1} + u + w,  z = x + v
// Sigma_w = Q (过程噪声方差),  Sigma_v = R (观测噪声方差)
int main() {
  double x_hat = 0.0, P_hat = 1.0;        // 后验均值 x_hat, 后验协方差 P_hat
  const double Q = 1.0, R = 0.5;          // 过程 / 观测噪声方差
  double u[2] = {1.0, 1.0};               // 输入序列
  double z[2] = {2.0, 3.1};               // 观测序列
  for (int k = 0; k < 2; ++k) {
    // 预测 (predict): x_check = x_hat + u,  P_check = P_hat + Q
    double x_check = x_hat + u[k];
    double P_check = P_hat + Q;
    // 增益 (gain): K = P_check / (P_check + R)
    double K = P_check / (P_check + R);
    // 校正 (correct): x_hat = x_check + K*(z - x_check)
    x_hat = x_check + K * (z[k] - x_check);
    // 协方差: 用 Joseph 形式更稳 P_hat = (1-K)^2 P_check + K^2 R
    P_hat = (1 - K) * (1 - K) * P_check + K * K * R;
    std::cout << "k=" << k << " x_hat=" << x_hat << " P_hat=" << P_hat << "\n";
  }
  return 0;   // 观察: P_hat 先因预测变大(+Q), 再因观测收缩(<R 与 <P_check)
}
\end{codebox}

% ════════════════════════════════════════════════════════════════
%  深化增节（Kalman 族收口）：吸收
%  Tutorial 20_经典非线性滤波族.md（Scaled UT / CKF / GHKF / Sparse-Grid / SR-UKF）
%  + Tutorial 40_Kalman族全景收口.md（Bell-Cathey IEKF=GN / LM-IEKF / IKFoM 流形 IEKF /
%    非线性 RTS：EKS/UKS/CKS / RTS 三视角与 1D 数值示例 / 平方根·信息形式 / MHE /
%    EnKF·Schmidt·H∞ / 全景对比大表）。
%  与前文重叠处一律 \cref 指过去（UKF/SPKF=algo:spkf-*、IEKF=eq:iekf、ISPKF=algo:ispkf、
%    线性 RTS=eq:rts、信息滤波=prop:info-filter、谱系=tab:filter-taxo），不重写。
%  右扰动为主；源里凡左扰动/JPL/互逆约定处改写为本书约定并在 note 注明。
% ════════════════════════════════════════════════════════════════
\section{Sigma 点族收口：缩放无迹变换、CKF 与 GHKF}
\label{sec:sigma-family}

\cref{sec:nonlin-pdf} 已把\emph{无迹变换}（UT）与无迹/sigma 点卡尔曼滤波（UKF/SPKF）讲透——\cref{eq:sigma-points} 的“生成—过非线性—重组”、\cref{algo:spkf-predict,algo:spkf-correct} 的预测/校正逐式、\cref{prop:spkf-ekf} 的“轻非线性退回 EKF 增益”、\cref{ex:three-methods} 的“sigma 点捕获真均值而线性化过自信”。本节不重复这些，而是\emph{补齐缺口}：(i) 工程实际用的\textbf{缩放无迹变换}（Scaled UT，三参数 $\alpha,\beta,\kappa$），它修正了 \cref{eq:sigma-points} 单参数 $\kappa$ 版在高维下 sigma 点飞太远的弊病；(ii) \textbf{球面径向容积卡尔曼滤波}（CKF）与它是 Scaled UT 特例的证明；(iii) \textbf{高斯--厄米卡尔曼滤波}（GHKF）与稀疏网格，并把整族归一为“同一广义高斯滤波 \cref{prop:ggf}、只换求积规则”的统一观点。本节脉络层为\textbf{次优}（理解 sigma 点族用），可在主线读完后回看。

\paragraph{动机：单参数 $\kappa$ 的 sigma 点在高维会“飞出去”。}
\cref{eq:sigma-points} 的展开半径是 $\sqrt{L+\kappa}$（$L$ 为被采样维度）。维度 $L$ 一大，半径随 $\sqrt{L}$ 增长，sigma 点被推到离均值很远处——落进 $\mathbf{g}$ 的强非线性区，重组出的均值反而偏差更大\cite{vandermerwe2004sigma}。更糟的是 \cref{eq:sigma-points} 的中心权 $\alpha_0=\kappa/(L+\kappa)$：要让半径合理就得取 $\kappa<0$（如 $\kappa=3-L$ 在 $L\ge3$ 时），于是 $\alpha_0<0$，重组协方差可能丢失正定。\textbf{反面}：若无视这一点、在 $15$ 维组合导航里硬套 $\kappa=3-L=-12$，中心权深度为负，Cholesky 在数分钟内抛异常（与 \cref{sec:kf} 数值陷阱同源）。Scaled UT 正是为解耦“展开半径”与“高阶矩补偿”而生。

\begin{definition}[缩放无迹变换（Scaled UT）]\label{def:scaled-ut}
对 $\mathbf{x}\sim\mathcal{N}(\boldsymbol{\mu},\boldsymbol{\Sigma})$、$\mathbf{x}\in\mathbb{R}^n$，引入三参数 $\alpha,\beta,\kappa$，令复合标度
\begin{equation}
  \lambda=\alpha^2(n+\kappa)-n,\qquad \gamma=\sqrt{n+\lambda},
  \label{eq:sut-lambda}
\end{equation}
取 $\boldsymbol{\Sigma}=\mathbf{L}\mathbf{L}^\top$（Cholesky，$\mathrm{col}_i\mathbf{L}$ 为第 $i$ 列），生成 $2n+1$ 个 sigma 点
\begin{equation}
  \boldsymbol{\mathcal{X}}_0=\boldsymbol{\mu},\quad
  \boldsymbol{\mathcal{X}}_i=\boldsymbol{\mu}+\gamma\,\mathrm{col}_i\mathbf{L},\quad
  \boldsymbol{\mathcal{X}}_{i+n}=\boldsymbol{\mu}-\gamma\,\mathrm{col}_i\mathbf{L},\quad i=1\dots n,
  \label{eq:sut-points}
\end{equation}
\textbf{均值权与协方差权分开}（这是与 \cref{eq:sigma-points} 单权版的关键差别）：
\begin{equation}
  W_0^m=\frac{\lambda}{n+\lambda},\quad
  W_0^c=\frac{\lambda}{n+\lambda}+(1-\alpha^2+\beta),\quad
  W_i^m=W_i^c=\frac{1}{2(n+\lambda)}\ (i\ge1).
  \label{eq:sut-weights}
\end{equation}
非线性传播 $\boldsymbol{\mathcal{Y}}_i=\mathbf{g}(\boldsymbol{\mathcal{X}}_i)$ 后用\emph{对应}权重重组：$\boldsymbol{\mu}_y=\sum_i W_i^m\boldsymbol{\mathcal{Y}}_i$、$\boldsymbol{\Sigma}_{yy}=\sum_i W_i^c(\boldsymbol{\mathcal{Y}}_i-\boldsymbol{\mu}_y)(\boldsymbol{\mathcal{Y}}_i-\boldsymbol{\mu}_y)^\top$、$\boldsymbol{\Sigma}_{xy}=\sum_i W_i^c(\boldsymbol{\mathcal{X}}_i-\boldsymbol{\mu})(\boldsymbol{\mathcal{Y}}_i-\boldsymbol{\mu}_y)^\top$。
\end{definition}

\noindent 三参数各司其职\cite{julier2002scaled,vandermerwe2004sigma}：
\begin{itemize}
  \item $\alpha\in[10^{-4},1]$ 控制\emph{展开半径}。$\alpha$ 越小 sigma 点越贴近 $\boldsymbol{\mu}$、抑制高阶非线性误差；它把半径从 $\sqrt{n+\kappa}$ 缩成 $\alpha\sqrt{n+\kappa}$（由 $\gamma=\alpha\sqrt{n+\kappa}$，因 $n+\lambda=\alpha^2(n+\kappa)$），再用 $W_0^c$ 里的 $(1-\alpha^2+\beta)$ 项把协方差“补”回来——这就是“先缩放、再补偿”。
  \item $\beta\ge0$ 修正\emph{先验分布形状}。\textbf{$\beta=2$ 对严格高斯最优}（匹配四阶矩，峰度 $3$）；$\beta=0$ 即不修正。
  \item $\kappa\in\mathbb{R}$ 为次级自由度，常取 $\kappa=0$ 或 $\kappa=3-n$。
\end{itemize}
\textbf{推荐默认}（van der Merwe 博士论文\cite{vandermerwe2004sigma}）：$\alpha=10^{-3},\beta=2,\kappa=0$，适用于多数机器人状态估计。把 \cref{def:scaled-ut} 的点与权代入 \cref{algo:spkf-predict,algo:spkf-correct}（用 $W^m$ 算均值、$W^c$ 算协方差，其余不变）即得工程版 UKF；\cref{eq:sigma-points} 的单参数版是其 $\alpha=1,\beta=0$ 特例（此时 $\lambda=\kappa$、$W_0^m=W_0^c=\kappa/(n+\kappa)$、$W_i=1/2(n+\kappa)$，恰为 \cref{sec:nonlin-pdf} 的权）。

\begin{note}
\textbf{精度阶（与 \cref{ex:three-methods} 衔接）.}\;对称 sigma 点对 $\boldsymbol{\mathcal{X}}_i,\boldsymbol{\mathcal{X}}_{i+n}$ 关于 $\boldsymbol{\mu}$ 成中心对称，故重组时\emph{奇数阶}泰勒项两两抵消：UKF 的\emph{均值估计精确到三阶}泰勒、\emph{协方差精确到二阶}，$\beta=2$ 再令四阶矩系数对齐高斯。相形之下 EKF 均值、协方差均仅一阶。这把 \cref{ex:three-methods} 一维 $x^2$ 的“sigma 点取 $\kappa=2$ 命中真四阶矩”推广到了任意维：在同为 $O(n^3)$ 的代价下，UKF 比 EKF “免费”多两阶精度且免雅可比\cite{barfoot2024state,julier2002scaled}。
\end{note}

\paragraph{容积卡尔曼滤波（CKF）：取消中心点、全正权重。}
UKF 在高维有两处不适\cite{arasaratnam2009cubature}：(a) $\alpha$ 选差则点飞远；(b) 中心协方差权 $W_0^c$ 可能为负（\cref{eq:sut-weights}，当 $\alpha<1$ 时 $1-\alpha^2>0$ 虽缓解，但 $\lambda<0$ 主导时仍可负），破坏正定。Arasaratnam 与 Haykin\cite{arasaratnam2009cubature} 不从“采样高斯”而从\emph{数值积分}重新出发：广义高斯滤波要算的本质是高斯加权积分 $\int_{\mathbb{R}^n}\mathbf{g}(\mathbf{x})\,\mathcal{N}(\mathbf{x};\boldsymbol{0},\mathbf{I})\,\mathrm{d}\mathbf{x}$，对它施\emph{三阶球面径向（spherical-radial）容积规则}，得到一个\textbf{无中心点、$2n$ 个点、权重一律 $1/(2n)$} 的方案。

\begin{derivation}[三阶球面径向容积规则]\label{der:ckf-rule}
目标积分 $I[\mathbf{g}]=\int_{\mathbb{R}^n}\mathbf{g}(\mathbf{x})\,\mathcal{N}(\mathbf{x};\boldsymbol{0},\mathbf{I})\,\mathrm{d}\mathbf{x}$。作球坐标分离 $\mathbf{x}=r\,\mathbf{y}$（$\mathbf{y}^\top\mathbf{y}=1$，$r\in[0,\infty)$），高斯权拆成径向与球面两部分\cite{arasaratnam2009cubature}：
\[
  I[\mathbf{g}]=\frac{1}{(2\pi)^{n/2}}\int_0^\infty\!\!\int_{\mathcal{U}_n}\mathbf{g}(r\mathbf{y})\,r^{n-1}e^{-r^2/2}\,\mathrm{d}\sigma(\mathbf{y})\,\mathrm{d}r,
\]
$\mathcal{U}_n$ 为单位球面。\emph{球面}部分用三阶对称规则——取 $2n$ 个交点 $\{\pm\mathbf{e}_i\}_{i=1}^n$（$\mathbf{e}_i$ 为坐标轴单位向量），对球面上三阶以下多项式精确；\emph{径向}部分用一阶广义高斯--拉盖尔规则，单结点落在 $r^2=n$（即 $r=\sqrt{n}$）处。两者合并：
\begin{equation}
  I[\mathbf{g}]\approx\sum_{i=1}^{2n}\frac{1}{2n}\,\mathbf{g}(\boldsymbol{\xi}_i),\qquad
  \boldsymbol{\xi}_i=\sqrt{n}\,\mathbf{e}_i,\ \ \boldsymbol{\xi}_{n+i}=-\sqrt{n}\,\mathbf{e}_i\ \ (i=1\dots n),
  \label{eq:ckf-points}
\end{equation}
对任意三阶以下多项式 $\mathbf{g}$ 精确。\hfill$\square$
\end{derivation}

\noindent CKF 算法即把 \cref{eq:ckf-points} 的标准点平移缩放到 $\mathbf{x}\sim\mathcal{N}(\boldsymbol{\mu},\boldsymbol{\Sigma})$：取 $\boldsymbol{\Sigma}=\mathbf{L}\mathbf{L}^\top$，容积点 $\boldsymbol{\mathcal{X}}_i=\boldsymbol{\mu}+\mathbf{L}\boldsymbol{\xi}_i$，其余预测/校正完全套用 \cref{algo:spkf-predict,algo:spkf-correct} 的结构——只是\emph{只有 $2n$ 个点、权重一律 $1/(2n)$，没有 $\alpha,\beta,\kappa$、没有负权重}。Arasaratnam--Haykin 原文定理证明该协方差更新保持对称正定。

\begin{proposition}[CKF $=$ Scaled UT 在 $(\alpha,\beta,\kappa)=(1,0,0)$ 的特例]\label{prop:ckf-ukf}
取 $\alpha=1,\beta=0,\kappa=0$ 代入 \cref{def:scaled-ut}，则 Scaled UT 退化为 CKF。
\end{proposition}
\begin{proof}
代 \cref{eq:sut-lambda}：$\lambda=\alpha^2(n+\kappa)-n=1\cdot(n+0)-n=0$。于是
\[
  \gamma=\sqrt{n+\lambda}=\sqrt{n}\ (\text{与 \cref{eq:ckf-points} 半径一致}),\quad
  W_0^m=\frac{\lambda}{n+\lambda}=0,\quad
  W_0^c=0+(1-1+0)=0,\quad
  W_i^m=W_i^c=\frac{1}{2(n+\lambda)}=\frac{1}{2n}.
\]
中心点权为零（等价于取消中心点），其余 $2n$ 点半径 $\sqrt{n}$、权 $1/(2n)$——与 CKF 完全相同。\hfill$\square$
\end{proof}

\begin{insight}
\textbf{CKF 与 UKF 不是“谁更准”，而是“要不要 $\beta$ 自由度”}。二者在三阶精度下等价（\cref{prop:ckf-ukf}）；真正的区别是 CKF 用\emph{全正权重}消除了 UKF 的“中心负权致协方差病态”风险，代价是放弃 $\beta$ 提供的四阶矩调节。工程准则极简：\emph{若}不确定 $\alpha,\beta,\kappa$ 怎么选、或状态维 $>10$ $\to$ 直接用 CKF；\emph{若}强非线性、低维、想用 $\beta=2$ 贴四阶矩 $\to$ 用 UKF\cite{arasaratnam2009cubature}。CKF 在航天器姿态、深空导航与高维状态空间模型中比 UKF 更受欢迎。其平滑版 CKS 见 \cref{sec:rts-family}，五阶 CKF（$2n^2+1$ 点）与连续--离散 CD-CKF 为进阶变体\cite{arasaratnam2011smoother}。
\end{insight}

\paragraph{高斯--厄米卡尔曼滤波（GHKF）与稀疏网格：精度可任意、但维度灾难。}
若要更高精度，可换\emph{高斯--厄米}（Gauss--Hermite）求积\cite{ito2000gaussian}：一维 $m$ 点规则对 $\le 2m-1$ 阶多项式精确，结点 $\xi_i$ 取（概率论约定）厄米多项式 $He_m(x)$ 的根、权由结点处多项式导数定。把它推到 $n$ 维需作\textbf{张量积}：$m^n$ 个点。$m=3,n=10$ 即 $3^{10}=59049$ 点，$n=15$ 已逾千万——\textbf{这正是 GHKF 只能用于 $n\le4$ 的根本原因}，故它的舞台是航天器姿态（$n=3$ 或 $6$）、化工过程（$n=2\text{--}4$）这类低维高精度场景。缓解维度灾难的标准手段是\emph{稀疏网格}（Smolyak\cite{jia2012sparse}）：从张量积里挑“精度预算最有效”的子集，点数从 $m^n$ 降到约 $O(n^{\ell}\log^{n-1}n)$（$\ell$ 为稀疏等级），把精度对维度从指数压成多项式，可覆盖 $n\approx8\text{--}12$ 的中维问题\cite{jia2012sparse}。

\begin{insight}
\textbf{统一观点：所有高斯滤波器只差一条求积规则}\cite{barfoot2024state}。\cref{prop:ggf} 把校正归结为算三个矩 $\boldsymbol{\mu}_y,\boldsymbol{\Sigma}_{yy},\boldsymbol{\Sigma}_{xy}$，而它们都是\emph{高斯加权积分}。于是 EKF、UKF、CKF、GHKF 的全部分歧，就是\textbf{用什么数值积分规则近似这些积分}（\cref{tab:gauss-quadrature}）。这把“一堆看似无关的滤波器”收束为同一模板的不同实例——也是把 \cref{tab:filter-taxo} 的“过非线性方法”一列落到可计算细节的钥匙。
\end{insight}

\begin{table}[H]\centering\small
\caption{高斯滤波器 $=$ 广义高斯滤波（\cref{prop:ggf}）$+$ 不同求积规则}
\label{tab:gauss-quadrature}
\begin{tabular}{lllll}
\toprule
滤波器 & 求积规则 & 点数 & 高斯下精度阶 & 调参 \\
\midrule
EKF & 一阶泰勒（伪积分） & $1$（仅 $\boldsymbol{\mu}$） & $1$ & 无 \\
二阶 EKF & 二阶泰勒（含 Hessian） & $1+$Hessian & $2$ & 无 \\
UKF / Scaled UT & 无迹（球面对称） & $2n+1$ & 均值 $3$、协方差 $2$ & $\alpha,\beta,\kappa$ \\
CKF & 三阶球面径向容积 & $2n$ & $3$ & 无 \\
五阶 CKF & 五阶球面径向容积 & $2n^2+1$ & $5$ & 无 \\
GHKF($m$) & $m$ 阶高斯--厄米张量积 & $m^n$ & $2m-1$ & $m$ \\
稀疏网格 GHKF($\ell$) & Smolyak 稀疏网格 & $\sim O(n^{\ell})$ & 约 $2\ell-1$ & $\ell$ \\
\bottomrule
\end{tabular}
\end{table}

\begin{pitfall}[编程陷阱：用协方差对角线而非 Cholesky 因子撒点]
生成 sigma/容积点时，\textbf{必须}用 $\mathbf{L}=\mathrm{chol}(\boldsymbol{\Sigma})$ 的列向量作展开方向（\cref{eq:sut-points,eq:ckf-points}）。若图省事只沿坐标轴用 $\sqrt{\Sigma_{ii}}$ 撒点，则当 $\boldsymbol{\Sigma}$ \emph{非对角}（状态间有相关）时会\emph{完全丢弃交叉协方差}，传播后的均值与协方差都错。原因：$\mathbf{L}$ 的列向量编码了状态间的完整相关结构，而对角线没有。这与 \cref{sec:nonlin-pdf} 用 Cholesky 是同一要求，但高维、强相关时后果尤烈。
\end{pitfall}

\begin{practice}
\begin{enumerate}
  \item 把 $(\alpha,\beta,\kappa)=(1,0,0)$ 代入 \cref{eq:sut-lambda,eq:sut-weights}，独立复算 \cref{prop:ckf-ukf}：验证 $W_0^m=W_0^c=0$、$W_i=1/(2n)$、$\gamma=\sqrt{n}$。
  \item 取 $n=20$、$(\alpha,\kappa)=(10^{-3},0)$，算 UKF 的 $\lambda$ 与 $W_0^c$，判断中心协方差权是否为负，并说明它为何威胁 Cholesky（联系 \cref{sec:kf} 数值陷阱）。
  \item （概念）为什么 GHKF 精度可任意高却仍少用？答案要点：张量积点数 $m^n$ 随维度指数爆炸，$n>4$ 即不可行（\cref{tab:gauss-quadrature}），稀疏网格只缓解到中维。
\end{enumerate}
\end{practice}

% ════════════════════════════════════════════════════════════════
\section{迭代变体收口：IEKF 是高斯牛顿、流形 IEKF 与 ISPKF}
\label{sec:iekf-gn}

\cref{sec:nonlin-pdf} 已给出迭代 EKF（IEKF，\cref{eq:iekf}）与迭代 sigma 点滤波（ISPKF，\cref{algo:ispkf}），并断言“IEKF 趋众数、ISPKF 趋均值”。本节补齐\emph{为什么}：把 IEKF 的迭代式与高斯牛顿（GN，\cref{ch:nlopt}）逐式对上（Bell--Cathey 定理），给出阻尼/线搜索加强版，并把它抬到\emph{流形}上（IKFoM）——这正是 \cref{ch:vio}、LIO 后端把“滤波”与“优化”焊在一起的代数纽带。本节脉络层为\textbf{骨干}（连接滤波与优化的关键）。

\paragraph{动机：EKF 的“一次线性化”往往不够，但又不想上整条轨迹的 BA。}
\cref{prop:ekf} 的“三宗罪”之首是\emph{只线性化一次}：若先验均值 $\check{\mathbf{x}}_k$ 离真后验众数远，单次雅可比 $\mathbf{H}$ 在错点求值，增益与更新都偏。一个直观例子：观测 $h(x)=\arctan(x)$，先验均值 $\check{x}=5$ 处斜率 $H=1/(1+25)\approx0.038$ 极平，再大的新息也几乎推不动估计；但真后验众数可能在 $x=1$ 附近，那里 $H=1/2$，线性化好得多\cite{barfoot2024state}。IEKF 的迭代正是反复在\emph{新的更好的点}重新线性化、逐步逼近众数——它把“EKF 的一步”补成“一个时间步内的完整 GN”，却\emph{不}付出整条轨迹批量 BA 的离线/非常数时间代价。

\paragraph{IEKF 的迭代更新式（重述并补全搬移项）。}
\cref{eq:iekf} 的紧凑写法在工程上易漏一项。完整的迭代更新（以 $\hat{\mathbf{x}}^{(0)}=\check{\mathbf{x}}_k$ 初始化，$j=0,1,\dots$）为\cite{barfoot2024state,bell1993iterated}
\begin{equation}
  \hat{\mathbf{x}}^{(j+1)}=\check{\mathbf{x}}_k+\mathbf{K}^{(j)}\Big[\mathbf{z}_k-\mathbf{h}(\hat{\mathbf{x}}^{(j)})-\underbrace{\mathbf{H}^{(j)}\big(\check{\mathbf{x}}_k-\hat{\mathbf{x}}^{(j)}\big)}_{\text{先验搬移项}}\Big],\quad
  \mathbf{K}^{(j)}=\check{\mathbf{P}}_k\mathbf{H}^{(j)\top}\big(\mathbf{H}^{(j)}\check{\mathbf{P}}_k\mathbf{H}^{(j)\top}+\boldsymbol{\Sigma}_{\mathbf{v}}\big)^{-1},
  \label{eq:iekf-full}
\end{equation}
$\mathbf{H}^{(j)}=\partial\mathbf{h}/\partial\mathbf{x}|_{\hat{\mathbf{x}}^{(j)}}$。\textbf{先验搬移项 $-\mathbf{H}^{(j)}(\check{\mathbf{x}}_k-\hat{\mathbf{x}}^{(j)})$ 不可省}：它把先验从“以 $\check{\mathbf{x}}_k$ 为均值”搬到当前工作点 $\hat{\mathbf{x}}^{(j)}$ 的线性化坐标系下。$j=0$ 时 $\hat{\mathbf{x}}^{(0)}=\check{\mathbf{x}}_k$ 使该项为零、\cref{eq:iekf-full} 退回 EKF（\cref{prop:ekf}）——故\textbf{EKF $=$ IEKF 第一次迭代}（与 \cref{eq:iekf} 一致）；若漏写此项，每次迭代会“忘记先验”，把后验越拉越偏。收敛后协方差只算一次：$\hat{\mathbf{P}}_k=(\mathbf{I}-\mathbf{K}^{(j^\star)}\mathbf{H}^{(j^\star)})\check{\mathbf{P}}_k$。

\begin{theorem}[Bell--Cathey：IEKF $=$ 对单步 MAP 做高斯牛顿]\label{thm:bell-cathey}
在加性高斯观测 $\mathbf{z}_k=\mathbf{h}(\mathbf{x}_k)+\mathbf{v}_k$、$\mathbf{v}_k\sim\mathcal{N}(\boldsymbol{0},\boldsymbol{\Sigma}_{\mathbf{v}})$ 下，迭代 \cref{eq:iekf-full} 与对单步负对数后验
\begin{equation}
  J(\mathbf{x})=\tfrac12\big\|\mathbf{x}-\check{\mathbf{x}}_k\big\|_{\check{\mathbf{P}}_k^{-1}}^2+\tfrac12\big\|\mathbf{z}_k-\mathbf{h}(\mathbf{x})\big\|_{\boldsymbol{\Sigma}_{\mathbf{v}}^{-1}}^2
  \label{eq:iekf-cost}
\end{equation}
做高斯牛顿\textbf{逐式相同}。故 IEKF 若收敛，其定点是该单步 MAP 目标的一阶驻点（众数），与 \cref{sec:nonlin-pdf} 的“IEKF$\to$众数”一致\cite{bell1993iterated,barfoot2024state}。
\end{theorem}
\begin{proof}
设 $\mathbf{L}_P\mathbf{L}_P^\top=\check{\mathbf{P}}_k$、$\mathbf{L}_R\mathbf{L}_R^\top=\boldsymbol{\Sigma}_{\mathbf{v}}$，把 \cref{eq:iekf-cost} 写成 $J=\tfrac12\|\mathbf{r}(\mathbf{x})\|^2$，残差
\[
  \mathbf{r}(\mathbf{x})=\begin{bmatrix}\mathbf{L}_P^{-1}(\mathbf{x}-\check{\mathbf{x}}_k)\\ \mathbf{L}_R^{-1}\big(\mathbf{z}_k-\mathbf{h}(\mathbf{x})\big)\end{bmatrix},\qquad
  \mathbf{J}_r^{(j)}=\frac{\partial\mathbf{r}}{\partial\mathbf{x}}\Big|_{\hat{\mathbf{x}}^{(j)}}=\begin{bmatrix}\mathbf{L}_P^{-1}\\ -\mathbf{L}_R^{-1}\mathbf{H}^{(j)}\end{bmatrix}.
\]
于是 $\mathbf{J}_r^{(j)\top}\mathbf{J}_r^{(j)}=\check{\mathbf{P}}_k^{-1}+\mathbf{H}^{(j)\top}\boldsymbol{\Sigma}_{\mathbf{v}}^{-1}\mathbf{H}^{(j)}$。GN 步 $\boldsymbol{\Delta}^{(j)}=-(\mathbf{J}_r^\top\mathbf{J}_r)^{-1}\mathbf{J}_r^\top\mathbf{r}(\hat{\mathbf{x}}^{(j)})$ 展开为
\[
  \boldsymbol{\Delta}^{(j)}=-\big(\check{\mathbf{P}}_k^{-1}+\mathbf{H}^{(j)\top}\boldsymbol{\Sigma}_{\mathbf{v}}^{-1}\mathbf{H}^{(j)}\big)^{-1}\Big[\check{\mathbf{P}}_k^{-1}(\hat{\mathbf{x}}^{(j)}-\check{\mathbf{x}}_k)-\mathbf{H}^{(j)\top}\boldsymbol{\Sigma}_{\mathbf{v}}^{-1}\big(\mathbf{z}_k-\mathbf{h}(\hat{\mathbf{x}}^{(j)})\big)\Big].
\]
对中间的信息矩阵逆用矩阵求逆引理（SMW，\cref{ch:slam_est}），代入两条恒等式
\[
  \big(\check{\mathbf{P}}_k^{-1}+\mathbf{H}^\top\boldsymbol{\Sigma}_{\mathbf{v}}^{-1}\mathbf{H}\big)^{-1}\mathbf{H}^\top\boldsymbol{\Sigma}_{\mathbf{v}}^{-1}=\mathbf{K},\qquad
  \big(\check{\mathbf{P}}_k^{-1}+\mathbf{H}^\top\boldsymbol{\Sigma}_{\mathbf{v}}^{-1}\mathbf{H}\big)^{-1}\check{\mathbf{P}}_k^{-1}=\mathbf{I}-\mathbf{K}\mathbf{H},
\]
（$\mathbf{K}$ 同 \cref{eq:iekf-full}）即得 $\hat{\mathbf{x}}^{(j+1)}=\hat{\mathbf{x}}^{(j)}+\boldsymbol{\Delta}^{(j)}=\check{\mathbf{x}}_k+\mathbf{K}^{(j)}[\mathbf{z}_k-\mathbf{h}(\hat{\mathbf{x}}^{(j)})-\mathbf{H}^{(j)}(\check{\mathbf{x}}_k-\hat{\mathbf{x}}^{(j)})]$——正是 \cref{eq:iekf-full}。\hfill$\square$
\end{proof}

\begin{insight}
\textbf{IEKF 的本质不是“多做几次 EKF 更新”，而是“在 MAP 目标 \cref{eq:iekf-cost} 上做 GN”}。每次迭代改善的是\emph{线性化点的质量}，不是信息量——故\textbf{收敛点与最大迭代数 $J_{\max}$ 无关}（凸邻域内始终是同一 MAP 解），$J_{\max}$ 只决定能否到达。最终协方差由收敛点 Hessian 定：$\hat{\mathbf{P}}_k^{-1}=\check{\mathbf{P}}_k^{-1}+\mathbf{H}^{(j^\star)\top}\boldsymbol{\Sigma}_{\mathbf{v}}^{-1}\mathbf{H}^{(j^\star)}$。推论：先验已准（$\check{\mathbf{x}}_k$ 近众数）则一次迭代足矣（退回 EKF）；先验差或观测强非线性则多迭代显著改善——这也解释了 FAST-LIO2 把 \verb@max_iteration@ 默认设 $3$--$4$：点到面残差在配准良好后近线性，再迭代无收益\cite{xu2022fastlio2}。Bell 进一步证明\emph{迭代平滑 $=$ 整条轨迹的 GN}\cite{bell1993iterated}，这是“IEKF $=$ 单步因子图、BA $=$ 全步因子图”谱系（\cref{eq:iekf-bridge}）的理论基石。
\end{insight}

\begin{pitfall}[概念误区：把 IEKF 迭代当成“多次利用同一观测”]
错误想法：“IEKF 对同一个 $\mathbf{z}_k$ 更新了 $J_{\max}$ 次，信息被用了 $J_{\max}$ 倍、协方差会越缩越小。” 实则不然：由上 insight，$J_{\max}$ 次迭代收敛到的定点\emph{与 $J_{\max}$ 无关}（凸条件下始终是同一 MAP 解），迭代次数只影响\emph{能否到达}、不影响到达后的协方差。最终协方差只由收敛点 Hessian 一次决定，绝非按迭代次数累乘——否则就成了“用一条观测人为制造 $J_{\max}$ 倍信息”的过自信病。
\end{pitfall}

\paragraph{阻尼与线搜索：当纯 GN 不收敛。}
纯 GN（\cref{eq:iekf-full}）在强非线性、初值差时可能不下降甚至发散。两种标准加强\cite{barfoot2024state}：
\begin{itemize}
  \item \textbf{LM--IEKF}（Levenberg--Marquardt 化）：在近似 Hessian 上加阻尼 $\mathbf{A}_\lambda^{(j)}=\check{\mathbf{P}}_k^{-1}+\mathbf{H}^{(j)\top}\boldsymbol{\Sigma}_{\mathbf{v}}^{-1}\mathbf{H}^{(j)}+\lambda\mathbf{I}$，$\lambda$ 大时趋最速下降（小步稳）、$\lambda$ 小时趋 GN（大步快）；$\lambda$ 按代价升降自适应调整（接受步缩小 $\lambda$、拒绝步放大 $\lambda$）\footnote{常用的 $\lambda$ 自适应策略（如 Nielsen 系数 $\max(1/3,1-(2\rho-1)^3)$，$\rho$ 为增益比）见 Skoglund 等 FUSION 2015；具体系数以原文为准。}。
  \item \textbf{线搜索 IEKF}：求出 GN 方向 $\boldsymbol{\Delta}^{(j)}$ 后取 $\hat{\mathbf{x}}^{(j+1)}=\hat{\mathbf{x}}^{(j)}+\alpha\boldsymbol{\Delta}^{(j)}$，$\alpha\in(0,1]$ 满足 Armijo 充分下降条件，否则回溯 $\alpha\leftarrow\alpha/2$。Särkkä\cite{sarkka2013bayesian} 给出其完整形式。
\end{itemize}
二者与 \cref{ch:nlopt} 的 LM/线搜索\emph{是同一套}——再次印证“IEKF 就是优化”。收敛判据：状态增量 $\|\hat{\mathbf{x}}^{(j+1)}-\hat{\mathbf{x}}^{(j)}\|<\varepsilon$ 或代价相对下降 $<\delta$，工程常设 $J_{\max}\in\{3,4,5\}$。

\paragraph{流形 IEKF（IKFoM）：把迭代抬到李群上。}
当状态在李群/流形 $\mathcal{M}$（如 $\SO(3)$、$\SE(3)$、$\mathrm{SE}_2(3)$）上时，\cref{eq:iekf-full} 的“减法”与“加法”都不合法，须换成 \cref{ch:lie} 的 $\boxplus/\boxminus$（$\mathbf{x}\boxplus\mathbf{u}=\mathbf{x}\cdot\mathrm{Exp}(\mathbf{u})$、$\boxminus$ 为其逆，右扰动约定）。He、Xu 与 Zhang 的 IKFoM\cite{he2021ikfom} 给出规范形式：误差态 $\delta\mathbf{x}^{(j)}\triangleq\mathbf{x}\boxminus\hat{\mathbf{x}}^{(j)}$，关键多出一个\textbf{流形曲率雅可比} $\mathbf{J}^{(j)}$——它把切空间扰动从 $\hat{\mathbf{x}}^{(j)}$ 搬到先验工作点 $\check{\mathbf{x}}_k$ 所需的线性化矩阵。GN 增量为
\begin{equation}
  \delta^{(j+1)}=-\big(\mathbf{J}^{(j)\top}\check{\mathbf{P}}_k^{-1}\mathbf{J}^{(j)}+\mathbf{H}^{(j)\top}\boldsymbol{\Sigma}_{\mathbf{v}}^{-1}\mathbf{H}^{(j)}\big)^{-1}\Big[\mathbf{J}^{(j)\top}\check{\mathbf{P}}_k^{-1}\big(\hat{\mathbf{x}}^{(j)}\boxminus\check{\mathbf{x}}_k\big)+\mathbf{H}^{(j)\top}\boldsymbol{\Sigma}_{\mathbf{v}}^{-1}\big(\mathbf{h}(\hat{\mathbf{x}}^{(j)})-\mathbf{z}_k\big)\Big],
  \label{eq:ikfom-gn}
\end{equation}
随后 $\hat{\mathbf{x}}^{(j+1)}=\hat{\mathbf{x}}^{(j)}\boxplus\delta^{(j+1)}$，收敛后 $\hat{\mathbf{P}}_k=(\mathbf{I}-\mathbf{K}\mathbf{H})\,\mathbf{J}^{-1}\check{\mathbf{P}}_k\mathbf{J}^{-\top}$。对 $\SO(3)$，$\mathbf{J}^{(j)}$ 由 \cref{ch:lie} 的左雅可比 $\mathbf{J}_l$ 在 $\mathbf{u}=\hat{\mathbf{x}}^{(j)}\boxminus\check{\mathbf{x}}_k$ 处给出（源用 $\mathbf{J}_l(\mathbf{u})^{-\top}$\footnote{曲率雅可比 $\mathbf{J}^{(j)}$ 的具体形式（$\mathbf{J}_l^{-\top}$ 抑或右雅可比 $\mathbf{J}_r$ 的某变体）随扰动约定而异，与本书右扰动约定的对齐以 He--Xu--Zhang 原文（eq.(27)--(28)）为准。}）；在 $\mathbb{R}^n$ 上 $\mathbf{J}=\mathbf{I}$、\cref{eq:ikfom-gn} 退回欧氏 IEKF。\textbf{曲率雅可比 $\mathbf{J}$ 是“在流形上做加法”的代价}：略去它（误当 $\mathbf{J}=\mathbf{I}$）会在大角度更新时引入系统误差。IKFoM 是 FAST-LIO/FAST-LIO2 的核心\cite{xu2021fastlio,xu2022fastlio2}：用户只填 $\mathbf{f},\partial\mathbf{f}/\partial\mathbf{x},\partial\mathbf{f}/\partial\mathbf{w},\mathbf{h},\partial\mathbf{h}/\partial\mathbf{x}$，模板自动处理 $\boxplus/\boxminus/\mathbf{J}$；其高效增益 $\mathbf{K}=(\mathbf{H}^\top\boldsymbol{\Sigma}_{\mathbf{v}}^{-1}\mathbf{H}+\check{\mathbf{P}}_k^{-1})^{-1}\mathbf{H}^\top\boldsymbol{\Sigma}_{\mathbf{v}}^{-1}$ 复杂度只随状态维 $n\approx23$ 而非测量维 $m\sim10^3\text{--}10^4$，是 LIO 实时的关键（这正是 \cref{prop:info-filter} 信息形式“校正即相加、免大逆”在流形迭代上的落地）。

\begin{note}
\textbf{ISPKF 的对偶定理（与 \cref{algo:ispkf} 衔接）.}\;
\cref{algo:ispkf} 的迭代 sigma 点滤波把 \cref{eq:iekf-full} 的\emph{解析}线性化换成\emph{统计线性回归}（用 sigma 点重组的 $\boldsymbol{\Sigma}_{xy},\boldsymbol{\Sigma}_{yy}$ 当“等效雅可比”）。与 \cref{thm:bell-cathey} 对偶的结论是：\textbf{ISPKF 收敛到后验\emph{均值}，而非众数}\cite{barfoot2024state}。根源在 \cref{ex:three-methods} 已点明——UT/sigma 点近似的是\emph{期望积分} $\mathbb{E}[\mathbf{h}(\mathbf{x})]$（矩匹配/KL 意义下的最优高斯，对应均值），泰勒近似的是\emph{点处梯度}（对应众数）。\cref{sec:nonlin-pdf} 末“IEKF$\to$众数、ISPKF$\to$均值”的数值证据（立体相机例 $\hat{x}_{\text{iekf}}=24.5694$ 落众数、$\hat{x}_{\text{ispkf}}=24.7414$ 落均值附近）即此二定理的可复算对照。García-Fernández 等把这一差异系统化为\emph{迭代后验线性化滤波}（IPLF），ISPKF 可视为其特例\footnote{IPLF 出处：A.\ F.\ García-Fernández, L.\ Svensson, M.\ R.\ Morelande, S.\ Särkkä, ``Posterior Linearization Filter: Principles and Implementation Using Sigma Points'', \emph{IEEE Trans.\ Signal Processing} 63(20):5561--5573, 2015；对应的迭代后验线性化平滑器 IPLS 则发表于 \emph{IEEE Trans.\ Automatic Control}（2017）。}。工业界少用 ISPKF：每迭代需重生成 sigma 点并过 $\mathbf{h}$，比 IEKF 贵 $5$--$10$ 倍，且其“均值语义”与因子图后端的“MAP/众数语义”割裂。
\end{note}

\begin{insight}
\textbf{从 IEKF 到因子图：一条“保留多少历史”的连续谱}。把单步 MAP \cref{eq:iekf-cost} 推广到整条轨迹的联合 MAP，并对其做 GN/LM，即 \cref{ch:nlopt} 的束调整（BA）/全图 SLAM。IEKF 是对该全局目标做\emph{顺序舒尔边缘化}：每步把 $\mathbf{x}_{0:k-1}$ 消成等价先验 $(\check{\mathbf{x}}_k,\check{\mathbf{P}}_k)$、仅对 $\mathbf{x}_k$ 做 GN。于是\cite{barfoot2024state,xu2022fastlio2}
\begin{equation}
  \underbrace{\text{EKF}}_{\text{单步、一次线性化}}\ \subset\ \underbrace{\text{IEKF}}_{\text{单步、迭代 GN（\cref{thm:bell-cathey}）}}\ \subset\ \underbrace{\text{SWF/MHE}}_{\text{窗口 GN（\cref{par:swf}）}}\ \subset\ \underbrace{\text{BA / 全图 SLAM}}_{\text{全轨迹 GN}}.
  \label{eq:iekf-bridge}
\end{equation}
四者只差一个参数——\emph{窗口大小} $N$（保留多少历史重线性化）：$N=1$ 只迭代校正即 IEKF；$N=5\text{--}20$ 为滑窗 MHE（VINS-Mono、FAST-LIO2）；$N\to\infty$ 为全局 BA（离线建图）。窗口越大，精度/一致性越好（可重线性化更多历史），但算力/内存/延迟越高。这把 \cref{sec:filter-vs-opt} 的“滤波 vs 优化”从二选一升级为\emph{一条可调谱}：故 MSCKF $\equiv$ 特殊的部分联合 IEKF（相机位姿联合保留、特征零空间投影消去），FAST-LIO2 $\equiv$ $23$ 维流形上的单步 MAP\cite{xu2022fastlio2,mourikis2007msckf}。
\end{insight}

\begin{practice}
\begin{enumerate}
  \item 沿 \cref{thm:bell-cathey} 证明的两条 SMW 恒等式，独立把 GN 步 $\boldsymbol{\Delta}^{(j)}$ 化简到 \cref{eq:iekf-full}，确认先验搬移项的来历。
  \item 用 $h(x)=\arctan(x)$、先验 $\check{x}=5$、$\check{P}=4$、$\boldsymbol{\Sigma}_{\mathbf{v}}=0.1$、观测 $z=\arctan(1)$，手算 EKF 一步与 IEKF 两步，比较 $\hat{x}$ 是否更靠近 $x=1$，体会“重线性化改善线性化点”。
  \item （概念）解释“IEKF 收敛点与 $J_{\max}$ 无关、只与是否到达有关”，并说明这为何反驳“迭代 $J_{\max}$ 次 $=$ 信息用 $J_{\max}$ 倍”。答错回看本节 insight 与 pitfall。
  \item （流形）在 \cref{eq:ikfom-gn} 中令 $\mathcal{M}=\mathbb{R}^n$（$\mathbf{J}=\mathbf{I}$、$\boxminus$ 为减、$\boxplus$ 为加），验证它退回欧氏 IEKF \cref{eq:iekf-full}。
\end{enumerate}
\end{practice}

% ════════════════════════════════════════════════════════════════
\section{平滑器族收口：非线性 RTS（EKS/UKS/CKS）与三视角}
\label{sec:rts-family}

\cref{sec:kf} 已给\emph{线性}RTS 平滑器（前向 KF $+$ 后向修正，\cref{eq:rts,eq:rts-cov}），并指出它精确求解原批量问题、$O(K)$。本节补两件事：(i) 把 RTS 推广到\emph{非线性}（EKS/UKS/CKS），它们与 \cref{sec:nonlin-pdf,sec:sigma-family} 的“过非线性三法”一一对应；(ii) 给出 RTS 的\textbf{三种等价视角}与一个\emph{可复算的 1D 数值示例}，把“平滑利用未来信息”从公式落成数字。本节脉络层为\textbf{次优}（离线精修与理解平滑用）。

\paragraph{动机：滤波只用过去，离线时未来数据本可让过去更准。}
\cref{sec:eskf-why} 已分清滤波（因果、只用 $\le k$）与平滑（非因果、可用未来）。\textbf{反面}：若有一条含\emph{异常观测}的序列，滤波在该时刻只能被异常值拽偏、无从挽回；但事后看，其\emph{后续}观测往往“证明”了系统并未真的偏离——平滑器正是把这份未来信息向后传回，修正过去的估计。线性时这由 \cref{eq:rts} 精确完成；非线性时则需把后向递推里的“矩与交叉协方差”改用近似算法。

\paragraph{非线性 RTS：换的只是“怎么算交叉协方差”。}
Särkkä\cite{sarkka2013bayesian} 把非线性平滑统一为\emph{高斯 RTS 家族}：后向递推的核心是\emph{滤波--预测交叉协方差} $\mathbf{D}_{k}=\mathrm{Cov}(\mathbf{x}_k,\mathbf{x}_{k+1}\mid\mathbf{z}_{0:k})$ 与预测协方差 $\check{\mathbf{P}}_{k+1}$，平滑增益 $\mathbf{G}_k=\mathbf{D}_k\check{\mathbf{P}}_{k+1}^{-1}$，再套 \cref{eq:rts-cov} 形式的修正。\emph{怎么近似 $\mathbf{D}_k,\check{\mathbf{P}}_{k+1}$} 决定了得到哪种平滑器——与 \cref{sec:nonlin-pdf} 的“过非线性三法”严格对应：

\begin{table}[H]\centering\small
\caption{非线性 RTS 平滑器家族（后向递推套 \cref{eq:rts,eq:rts-cov}，仅近似 $\mathbf{D}_k,\check{\mathbf{P}}_{k+1}$ 的方法不同）}
\label{tab:nonlinear-rts}
\begin{tabular}{lllll}
\toprule
平滑器 & 近似方式 & 对应滤波器 & 雅可比 & 权重正定 \\
\midrule
EKS & 一阶泰勒 $\mathbf{F}=\partial\mathbf{f}/\partial\mathbf{x}$ & EKF（\cref{prop:ekf}） & 需 & --- \\
UKS / URTSS & 无迹变换，$2n+1$ 点 & UKF（\cref{algo:spkf-predict}） & 免 & $n\le3$ 常正 \\
CKS & 三阶球面径向容积，$2n$ 点 & CKF（\cref{prop:ckf-ukf}） & 免 & \textbf{始终正} \\
\bottomrule
\end{tabular}
\end{table}

\noindent EKS 在后向 \cref{eq:rts} 中用 $\mathbf{F}=\partial\mathbf{f}/\partial\mathbf{x}|_{\hat{\mathbf{x}}_{k,f}}$ 替线性 $\mathbf{A}_k$；UKS\cite{sarkka2008unscented} 与 CKS\cite{arasaratnam2011smoother} 则用各自 sigma/容积点重组 $\mathbf{D}_k,\check{\mathbf{P}}_{k+1}$（CKS 因 \cref{eq:ckf-points} 权重全正而数值最稳）。\emph{迭代 EKS}（对整条 RTS 再做 GN，IPLS）等价于整条轨迹 MAP 的高斯牛顿，是 \cref{thm:bell-cathey} 的 Bell 推广在平滑上的自然延伸\cite{bell1993iterated}。流形版 RTS（IRTS/ERTS）把 \cref{eq:rts} 的减/加换成 $\boxminus/\boxplus$、增量经 $\boxplus$ 注入；在群仿射系统下，不变 RTS 的雅可比与状态估计无关，其一致性优势与 InEKF（\cref{paper:inekf}）同源\cite{vanderlaan2020invariant}（van der Laan, Cohen, Arsenault, Forbes,《The Invariant Rauch--Tung--Striebel Smoother》, IEEE RA-L 5(4):5067--5074, 2020）。

\begin{theorem}[RTS 的三种等价视角]\label{thm:rts-views}
线性高斯下，下列三者给出\emph{相同}的平滑边缘后验 $\{\hat{\mathbf{x}}_{k\mid K},\hat{\mathbf{P}}_{k\mid K}\}$\cite{barfoot2024state,sarkka2013bayesian}：
\begin{enumerate}
  \item \textbf{前向--后向 RTS}：前向跑 KF，再后向修正 \cref{eq:rts,eq:rts-cov}（单向存储、实现最省）；
  \item \textbf{双滤波器（Fraser--Potter）}：前向 KF 与一个\emph{后向信息滤波器}，二者在每时刻融合；后向初值须用信息形式 $\boldsymbol{\Lambda}\to\boldsymbol{0}$（弥散先验，\cref{prop:info-filter} 的“全局无知”）；
  \item \textbf{批量 MAP}：把 $\mathbf{x}_{0:K}$ 拼成大向量，其信息矩阵块三对角，做 Cholesky——\textbf{前向消元 $=$ KF、后向回代 $=$ RTS}。
\end{enumerate}
\end{theorem}
\begin{proof}[要点]
视角 3 是 \cref{thm:kf-map} 与 \cref{sec:eskf-appendix} 推导的直接推论：块三对角信息矩阵的稀疏 Cholesky，前向消元逐式即 KF（\cref{eq:kf-predict,eq:kf-gain,eq:kf-correct}），后向回代逐式即 RTS（\cref{eq:rts,eq:rts-cov}）。视角 1$\leftrightarrow$3 的等价已在 \cref{sec:eskf-appendix} 的 RTS 后向协方差推导中由矩阵求逆引理坐实。视角 2 由“后验 $=$ 前向因子 $\times$ 后向因子”（\cref{eq:eskf-fwd-bwd}）两因子各自递归、再相乘（信息相加）得到。\hfill$\square$
\end{proof}

\begin{insight}
\textbf{教学上务必掌握视角 3}：它把“滤波 vs 平滑”还原为“消元 vs 回代”，与 \cref{sec:eskf-why} 的“滤波 $=$ 批量 MAP 的因果递归”一脉相承，也为 \cref{ch:nlopt} 的 iSAM2/贝叶斯树（增量重排消元顺序）打地基。一句话：\emph{KF 是从过去往未来消元、RTS 是从未来往过去回代，二者合起来恰好解完整条轨迹的那个块三对角系统}。
\end{insight}

\begin{example}[1D 匀速模型：5 步 KF$+$RTS，看“未来信息”修正过去]\label{ex:rts-1d}
取状态 $\mathbf{x}_k=(p_k,v_k)^\top$（位置、速度），$\mathbf{A}=\big[\begin{smallmatrix}1&1\\0&1\end{smallmatrix}\big]$（$\Delta t=1$）、$\mathbf{C}=[1\ 0]$（只测位置）、$\boldsymbol{\Sigma}_{\mathbf{w}}=0.25\mathbf{I}_2$、$\boldsymbol{\Sigma}_{\mathbf{v}}=[1]$、初值 $\hat{\mathbf{x}}_{0}=(0,1)^\top$、$\hat{\mathbf{P}}_0=\mathbf{I}_2$。真实轨迹匀速 $v=1$（$p_{\text{true}}=0,1,2,3,4$），含噪观测 $\mathbf{z}=(0.2,0.8,2.5,2.7,4.3)$。前向 KF 与后向 RTS（\cref{eq:rts,eq:rts-cov}）的位置分量结果如下（表内数值由本节末代码框复算核对，取两位小数）：

\begin{center}\small
\begin{tabular}{cccccc}
\toprule
时刻 $k$ & 观测 $z_k$ & 滤波 $\hat{p}_{k\mid k}$ & 平滑 $\hat{p}_{k\mid K}$ & 真值 $p_{\text{true}}$ & 滤波 $P_{11}$ / 平滑 $P_{11}$ \\
\midrule
0 & 0.2 & 0.10 & 0.06 & 0 & 0.50 / 0.39 \\
1 & 0.8 & 0.91 & 1.05 & 1 & 0.64 / 0.33 \\
2 & 2.5 & 2.30 & 2.11 & 2 & 0.71 / 0.34 \\
3 & 2.7 & 2.92 & 3.07 & 3 & 0.70 / 0.38 \\
4 & 4.3 & 4.15 & 4.15 & 4 & 0.69 / 0.69 \\
\bottomrule
\end{tabular}
\end{center}

\noindent 三个关键观察：(1) \textbf{平滑方差恒不大于滤波方差}（$P_{k\mid K}\preceq P_{k\mid k}$，因 $\check{\mathbf{P}}_{k+1\mid K}-\check{\mathbf{P}}_{k+1}\preceq\mathbf{0}$，未来数据只增信息），$k=4$ 处二者相等（最后一步无未来可用）；(2) \textbf{$k=3$ 修正显著}：偏低观测 $z_3=2.7$ 把滤波估计拽到 $\hat{p}_{3\mid3}=2.92$，平滑后被后续观测 $z_4=4.3$（暗示速度未降）向上修正到 $\hat{p}_{3\mid K}=3.07$（越过真值 $3.0$ 一点，体现“未来信息”的反向拉动）；(3) \textbf{中间时刻协方差改善最大}（$k=1,2$ 同受前后观测约束，降近 $50\%$）。这与 \cref{thm:rts-views} 视角 3 一致：平滑 $=$ 同时考虑全部观测、解一个块三对角系统。
\end{example}

\begin{pitfall}[陷阱：对含延迟观测的系统直接套 RTS]
某些传感器有固定延迟（如视觉管线 $50$\,ms）。若把延迟观测按\emph{接收}时戳对齐就跑 RTS，会出现“平滑后反而比滤波差、残差统计异常”。根因：RTS 假设 $\mathbf{z}_k$ 是 $\mathbf{x}_k$ 的观测；若 $\mathbf{z}_k$ 实为 $\mathbf{x}_{k-d}$ 的观测（$d$ 为延迟步），直接套用会把“未来观测信息”错误归因到当前状态。正确做法：(a) 把状态增广为 $(\mathbf{x}_k,\dots,\mathbf{x}_{k-d})$ 让延迟观测对应正确分量；(b) 用固定滞后平滑器（fixed-lag smoother）显式处理延迟窗；(c) 在因子图中把延迟观测连到正确时刻节点。
\end{pitfall}

\begin{codebox}[Python]{1D 匀速 KF+RTS 复算（验证 ex:rts-1d 的 P\_smooth <= P\_filt）}
import numpy as np
# system: x=(p,v); A const-velocity, C observes position only
A = np.array([[1.0, 1.0], [0.0, 1.0]])
C = np.array([[1.0, 0.0]])
Sw = 0.25 * np.eye(2)          # process noise cov
Sv = np.array([[1.0]])         # measurement noise cov
x = np.array([0.0, 1.0]); P = np.eye(2)
z = [0.2, 0.8, 2.5, 2.7, 4.3]
xf, Pf, xp, Pp = [], [], [], []          # filter / predicted stores
for k in range(len(z)):
    if k > 0:                            # predict (k=0 starts from prior)
        x = A @ x; P = A @ P @ A.T + Sw
    xp.append(x.copy()); Pp.append(P.copy())
    S = C @ P @ C.T + Sv                 # innovation cov
    K = P @ C.T @ np.linalg.inv(S)
    x = x + K @ (np.array([z[k]]) - C @ x)
    P = (np.eye(2) - K @ C) @ P
    xf.append(x.copy()); Pf.append(P.copy())
xs = [None]*len(z); Ps = [None]*len(z)   # smoothed stores
xs[-1], Ps[-1] = xf[-1].copy(), Pf[-1].copy()
for k in range(len(z)-2, -1, -1):        # RTS backward pass
    G = Pf[k] @ A.T @ np.linalg.inv(Pp[k+1])
    xs[k] = xf[k] + G @ (xs[k+1] - xp[k+1])
    Ps[k] = Pf[k] + G @ (Ps[k+1] - Pp[k+1]) @ G.T
    assert np.all(np.linalg.eigvals(Pf[k] - Ps[k]) >= -1e-9)  # P_smooth <= P_filt
for k in range(len(z)):
    print("k=%d  p_filt=%.2f  p_smooth=%.2f  Pf=%.2f  Ps=%.2f"
          % (k, xf[k][0], xs[k][0], Pf[k][0,0], Ps[k][0,0]))
\end{codebox}

\begin{practice}
\begin{enumerate}
  \item 跑上面的代码框，核对 \cref{ex:rts-1d} 表内数字；观察 $k=3$ 的位置如何被未来观测 $z_4$ 拉回真值。
  \item （概念）把 \cref{tab:nonlinear-rts} 的三种非线性平滑器与 \cref{sec:nonlin-pdf,sec:sigma-family} 的“过非线性三法”一一对应，说明 CKS 为何数值最稳。答案要点：CKS 用 \cref{eq:ckf-points} 全正权重容积规则。
  \item 把 \cref{ex:rts-1d} 扩为 2D（状态 $(p_x,v_x,p_y,v_y)$、$\mathbf{A}$ 块对角、$\mathbf{C}=[\mathbf{I}_2\ \mathbf{0}]$），生成含转弯的 10 步轨迹，验证逐步 $P_{k\mid K}\preceq P_{k\mid k}$，并解释为何转弯处平滑改善大于直线段。
\end{enumerate}
\end{practice}

% ════════════════════════════════════════════════════════════════
\section{平方根与信息形式的 sigma 点版；Kalman 族全景}
\label{sec:srukf-panorama}

本节收两件事：(i) 把 \cref{sec:kf} 提过的\emph{平方根滤波}思想（SRIF/UD）补到\emph{sigma 点}族——\textbf{平方根 UKF}（SR-UKF），它是 UKF/CKF 长时运行不“爆协方差”的工程必备；(ii) 给出贯通全章的\textbf{Kalman 族全景对比表}与“四正交坐标轴”总览，把线性 KF$\to$非线性$\to$流形$\to$迭代/平滑$\to$非高斯这一整族用一张图组织起来。本节脉络层为\textbf{骨干小结}（选型与检索用）。

\paragraph{平方根 UKF：传播 Cholesky 因子而非协方差本身。}
UKF/CKF 每步要 $\mathbf{L}=\mathrm{chol}(\boldsymbol{\Sigma})$。$\boldsymbol{\Sigma}$ 理论上半正定，但因 (a) 浮点累积、(b) \cref{eq:sut-weights} 的 $W_0^c$ 可能为负、(c) 校正减法 $\hat{\boldsymbol{\Sigma}}=\check{\boldsymbol{\Sigma}}-\mathbf{K}\boldsymbol{\Sigma}_{yy}\mathbf{K}^\top$ 破坏正定，实际系统中 Cholesky 常在运行数分钟后失败（与 \cref{sec:kf} 数值陷阱同根）。救援手段由轻到重：Joseph 形式校正、改 $\mathbf{L}\mathbf{D}\mathbf{L}^\top$/特征分解容忍半正定、每步对称化 $\boldsymbol{\Sigma}\leftarrow\tfrac12(\boldsymbol{\Sigma}+\boldsymbol{\Sigma}^\top)$、加对角抖动 $\boldsymbol{\Sigma}\leftarrow\boldsymbol{\Sigma}+\varepsilon\mathbf{I}$；而\emph{从根本上}解决的是 SR-UKF\cite{vandermerwe2004sigma}——\textbf{全程只存并传播平方根因子 $\mathbf{S}$（$\boldsymbol{\Sigma}=\mathbf{S}\mathbf{S}^\top$），不存 $\boldsymbol{\Sigma}$}。sigma 点直接用 $\mathbf{S}$ 生成；协方差更新改用两件数值线性代数原语：

\begin{itemize}
  \item \textbf{QR 分解}把各 sigma 点偏差列拼成的矩阵正交三角化，得预测平方根因子的转置
  \begin{equation}
    \check{\mathbf{S}}_k=\mathrm{qr}\Big\{\big[\sqrt{W_1^c}\,(\check{\boldsymbol{\mathcal{X}}}_{k,1}-\check{\mathbf{x}}_k)\ \cdots\ \sqrt{W_{2n}^c}\,(\check{\boldsymbol{\mathcal{X}}}_{k,2n}-\check{\mathbf{x}}_k)\ \ \sqrt{\boldsymbol{\Sigma}_{\mathbf{w}}}\big]\Big\},
    \label{eq:srukf-qr}
  \end{equation}
  （记号 $\mathrm{qr}\{\cdot\}$ 取 QR 分解的上三角因子 $\mathbf{R}$ 的转置作下三角平方根因子）；
  \item \textbf{Cholesky 秩一更新/降秩}（\verb@cholupdate@）处理中心点贡献——$W_0^c>0$ 时作 update、$W_0^c<0$ 时作 downdate，从而即便中心权为负也\emph{在因子层面}保持数值稳定（这正是 SR-UKF 优于普通 UKF 之处）。
\end{itemize}
SR-UKF 与普通 UKF 数学上等价、复杂度同为 $O(n^3)$，但条件数减半（对因子操作把 $\kappa(\boldsymbol{\Sigma})$ 开方回 $\kappa(\mathbf{S})$，理由同 \cref{sec:kf} 平方根滤波段与 \cref{ch:nlopt} 的 $\sqrt{}$SAM）。CKF 同理有平方根版（SR-CKF）。

\begin{note}
\textbf{平方根与信息形式的两个家族（与 \cref{prop:info-filter} 衔接）.}\;本书的“数值稳健化”分两条线，勿混：(1) \emph{平方根线}传播协方差或信息矩阵的\emph{Cholesky 因子}——线性 KF 侧为 SRIF（信息阵平方根，与 \cref{ch:nlopt} 的 $\sqrt{}$SAM 同对象的递归版）与 UD 分解（Bierman），sigma 点侧为本节 SR-UKF/SR-CKF；目的都是把条件数开方、避免显式求逆与减法相消。(2) \emph{信息形式线}（\cref{prop:info-filter} 的 $\boldsymbol{\Lambda},\boldsymbol{\eta}$）改变的是\emph{代数对偶}（校正变相加、可表全局无知、SLAM 下近似稀疏），不一定配平方根。二者可叠加（SRIF $=$ 信息形式 $+$ 平方根），但解决的问题正交：信息形式省“多观测的大逆”，平方根保“长跑的正定与精度”。\cref{sec:kf} 已把平方根的完整数值论证指向 \cref{ch:nlopt}（优化侧），本节只补 sigma 点族的落地。
\end{note}

\paragraph{Kalman 族全景：四条正交坐标轴。}
把全章（含本章深化）的方法摆在一起，会发现整族由\textbf{四条正交坐标轴}生成\cite{barfoot2024state}：
\begin{enumerate}
  \item \textbf{先验/后验高斯假设}——高斯（KF 族）还是非高斯（粒子滤波、直方图，\cref{tab:filter-taxo}）；
  \item \textbf{线性化策略}——精确（线性 KF）、一阶泰勒（EKF）、sigma/容积/高斯--厄米点（UKF/CKF/GHKF，\cref{tab:gauss-quadrature}）；
  \item \textbf{状态空间几何}——欧氏 $\mathbb{R}^n$（KF/EKF/UKF）还是流形/李群（ESKF/InEKF/IKFoM，\cref{sec:eskf}、\cref{paper:inekf}、\cref{eq:ikfom-gn}）；
  \item \textbf{迭代/平滑维度}——单步、迭代（IEKF/ISPKF，\cref{thm:bell-cathey}）、平滑（RTS/EKS/UKS/CKS，\cref{sec:rts-family}）、窗口/全图（SWF/MHE/BA，\cref{eq:iekf-bridge}）。
\end{enumerate}
任一具体滤波器就是这四轴上各取一个坐标的组合；沿“几何精致度”（欧氏$\to$流形$\to$不变/等变）与“计算策略”（单步$\to$迭代$\to$平滑$\to$联合 MAP）两个方向，\emph{精度与一致性单调上升、实现复杂度与算力也单调上升}——工程选型即在此二维上找帕累托最优。

\begin{center}
\small
\begin{tikzpicture}[
  >=Stealth,
  ax/.style={->, thick, gray!55!black},
  lab/.style={font=\scriptsize, align=center},
  dotk/.style={circle, fill=main!70!black, inner sep=1.3pt},
  dotm/.style={circle, fill=second!70!black, inner sep=1.3pt}]
  % axes: x = geometric refinement, y = computation strategy
  \draw[ax] (0,0) -- (8.6,0) node[right, lab]{几何精致度\\（欧氏$\to$流形$\to$不变/等变）};
  \draw[ax] (0,0) -- (0,5.0) node[above, lab]{计算策略\\（单步$\to$迭代$\to$平滑$\to$联合）};
  % single-step row
  \node[dotk] (kf) at (0.7,0.6) {}; \node[lab, below=1pt of kf]{KF};
  \node[dotk] (ekf) at (2.3,0.6) {}; \node[lab, below=1pt of ekf]{EKF/UKF\\CKF};
  \node[dotm] (eskf) at (4.4,0.6) {}; \node[lab, below=1pt of eskf]{ESKF};
  \node[dotm] (inekf) at (6.6,0.6) {}; \node[lab, below=1pt of inekf]{InEKF/EqF};
  % iterated row
  \node[dotk] (iekf) at (2.3,1.9) {}; \node[lab, right=1pt of iekf]{IEKF/ISPKF};
  \node[dotm] (ikfom) at (5.5,1.9) {}; \node[lab, right=1pt of ikfom]{IKFoM};
  % smoothing row
  \node[dotk] (rts) at (2.3,3.2) {}; \node[lab, right=1pt of rts]{RTS/EKS/UKS/CKS};
  \node[dotm] (irts) at (5.5,3.2) {}; \node[lab, right=1pt of irts]{IRTS/ERTS};
  % joint row
  \node[dotk] (mhe) at (3.2,4.4) {}; \node[lab, right=1pt of mhe]{SWF/MHE/BA};
  % monotone arrow
  \draw[->, thick, third!70!black, dashed] (1.0,0.95) -- (5.0,4.1)
    node[midway, sloped, above, lab]{精度/一致性 $\uparrow$、算力 $\uparrow$};
\end{tikzpicture}
\end{center}

\begin{table}[H]\centering\small
\caption{Kalman 族全景对比（按“线性$\to$非线性欧氏$\to$流形$\to$迭代/平滑$\to$非高斯”排列；本书右扰动为主，左扰动来源已注明转换）}
\label{tab:kalman-panorama}
\begin{tabularx}{\linewidth}{lLLLllL}
\toprule
方法 & 误差/几何 & 系统假设 & 过非线性 & 精度阶 & 迭代/平滑 & 一致性 \\
\midrule
KF & 加性 $\mathbb{R}^n$ & 线性高斯 & 无 & 精确 & 单步 & 天然 \\
信息滤波 IF & 加性（对偶） & 线性高斯 & 无 & 精确 & 单步 & 同 KF \\
SR-KF（Bierman/UD） & 加性 $\mathbb{R}^n$ & 线性高斯 & 无 & 精确 & 单步 & 天然（数值更稳） \\
EKF & 加性 $\mathbb{R}^n$ & 非线性高斯 & 一阶泰勒 & $1$ & 单步 & 伪观测，需 FEJ \\
二阶 EKF & 加性 $\mathbb{R}^n$ & 非线性 & 二阶泰勒 & $2$ & 单步 & 同 EKF \\
IEKF & 加性 $\mathbb{R}^n$ & 非线性 & 一阶泰勒迭代 & $1$（GN） & 迭代 & 同 EKF \\
UKF / Scaled UT & 加性 $\mathbb{R}^n$ & 非线性高斯 & 无迹 $2n+1$ 点 & 均值 $3$/协方差 $2$ & 单步 & 无 \\
SR-UKF & 加性 $\mathbb{R}^n$ & 非线性 & 无迹 $+$ Cholesky 更新 & 同 UKF & 单步 & 数值稳定 \\
CKF & 加性 $\mathbb{R}^n$ & 非线性 & 三阶容积 $2n$ 点 & $3$ & 单步 & 无（权全正） \\
GHKF($m$) & 加性 $\mathbb{R}^n$（低维） & 非线性 & 高斯--厄米 $m^n$ & $2m-1$ & 单步 & 无（维度灾难） \\
ISPKF & 加性 $\mathbb{R}^n$ & 强非线性 & 无迹 $+$ GN & 收敛到\emph{均值} & 迭代 & 无 \\
MEKF & 乘性 $\SO(3)$ & 非线性 $+$ 姿态 & 一阶（切空间） & $1$ & 单步 & 较好 \\
ESKF & 误差态 $\SE(3)$ & 非线性 $+$ 流形 & 一阶 on $\delta\mathbf{x}$ & $1$ & 单步 $+$ 复位 & 一般，需复位雅可比 \\
InEKF & 不变误差，矩阵李群 & 群仿射 & 一阶 on $\boldsymbol{\eta}$ & $1$ & 单步 & \textbf{天然}（传播与状态无关） \\
IKFoM & $\boxminus$ $+$ 曲率 $\mathbf{J}$ & 非线性流形 & 一阶 $+$ 流形 GN & $1$（迭代） & 迭代 & 类 IEKF \\
EqF & 等变误差，齐次空间 & 可积对称群 & 一阶 & $1$ & 单步 & 最强几何一致性 \\
EKS/UKS/CKS & 加性/流形 & 非线性高斯 & 泰勒/无迹/容积 $+$ RTS & $1$/$2$/$3$ & 平滑 & 同对应滤波 \\
RTS on Manifold & $\boxminus$ & 非线性流形 & RTS $+$ $\boxplus$ & $1$ & 平滑 & 群仿射下天然 \\
MHE & 加性/流形 & 非线性 $+$ 约束 & NLP $+$ 到达代价 & 可迭代 GN & 窗口 & FEJ $+$ 约束 \\
EnKF & 加性 $\mathbb{R}^n$（超高维） & 非线性高斯 & $N_e$ 集合隐式协方差 & 蒙特卡洛 & 单步 & 需局部化 $+$ 膨胀 \\
Schmidt-KF & 加性，含“考虑”参数 & 线性/非线性 & 参数不估只携不确定度 & 一阶 & 单步 & 保守（诚实） \\
$H_\infty$ 滤波 & 加性 $\mathbb{R}^n$ & 非线性 & 极小极大（最坏情形） & 一阶 & 单步 & 鲁棒（无概率假设） \\
粒子滤波 & 经验分布 & 任意非高斯 & $N_p$ 粒子 & 蒙特卡洛 & 单步 & 无 \\
RBPF（FastSLAM） & 混合 & 条件高斯 & 位姿粒子 $+$ 地图 EKF & 蒙特卡洛 $+$ 一阶 & 单步 & 粒子退化 \\
\bottomrule
\end{tabularx}
\end{table}

\noindent 表中横跨四轴：上半部沿“几何”从加性欧氏（KF/EKF/UKF/CKF/GHKF）走到乘性/流形（MEKF/ESKF）、再到不变/等变（InEKF/EqF）；纵向沿“计算策略”从单步到迭代（IEKF/ISPKF/IKFoM）、平滑（EKS/UKS/CKS/RTS-on-Manifold）、窗口（MHE）；末四行（EnKF/Schmidt/$H_\infty$/粒子）是离主轴的专用变体——EnKF 占“超高维 $+$ 蒙特卡洛”角（用 $N_e\ll n$ 个集合成员隐式表示协方差，气象/海洋专用，需局部化与协方差膨胀防坍缩），Schmidt-KF 占“保守协方差 $+$ 不估讨厌参数”角（把标定残余等缓变参数“考虑”进协方差但不估其值），$H_\infty$ 占“无概率假设 $+$ 极小极大鲁棒”角（噪声统计未知/对抗时用），粒子滤波见 \cref{algo:pf}。

\begin{note}
\textbf{跨库选型与约定雷区（与本章记号约定呼应）.}\;落地这些方法时，最大的坑不是公式而是\emph{约定}\cite{sola2017quaternion}：\textbf{四元数} Hamilton（GTSAM/Sophus/manif/VINS-Mono/FAST-LIO2）vs JPL（OpenVINS/MSCKF\_VIO），互转须同时核对存储顺序、虚部符号、主动/被动旋转，绝非一次转置或变号；\textbf{扰动方向}右扰动（Sophus/manif/GTSAM/VINS-Mono/IKFoM，本书主用）vs 左扰动（OpenVINS 全局误差），雅可比差一个伴随 $\mathrm{Ad}$（\cref{ch:lie}）。一致性方案三条技术路线：\textbf{FEJ}（OpenVINS/VINS-Mono/MSCKF）、\textbf{不变滤波}（InEKF/UKF-on-Manifold）、\textbf{等变滤波}（EqF）。工程组合：产品线常用 FAST-LIO2 $+$ OpenVINS/VINS-Mono $+$ GTSAM $+$ manif/Sophus；学术对比常用 kalmanif $+$ UKF-M $+$ invariant-ekf。这与本章开头“跨书雷区”记号约定一以贯之——\emph{先对齐约定，再抄公式}。
\end{note}

\begin{practice}
\begin{enumerate}
  \item （概念）说出 SR-UKF 相对普通 UKF 的两点好处，并指出它靠哪两个数值原语（\cref{eq:srukf-qr} 与 \verb@cholupdate@）实现“即便 $W_0^c<0$ 也保稳定”。
  \item 用“四正交坐标轴”给下列方法各标四个坐标：EKF、CKF、ESKF、IKFoM、RTS、粒子滤波（参照 \cref{tab:kalman-panorama}）。
  \item （概念）从 \cref{tab:kalman-panorama} 论证“沿几何与计算两个方向，精度/一致性与算力同向上升”，并据此为“高频 LIO 实时”与“离线高精度建图”各选一个方法、说明理由。答案要点：实时选 IKFoM/FAST-LIO2（单步流形迭代、复杂度随状态维），离线选 BA 或 CKS（平滑/全图、精度优先）。
\end{enumerate}
\end{practice}

% ════════════════════════════════════════════════════════════════
\section*{常见误解汇总}
\addcontentsline{toc}{section}{常见误解汇总}
\begin{table}[H]\centering\small
\begin{tabular}{p{0.46\textwidth} p{0.46\textwidth}}
\toprule
\textbf{常见误解} & \textbf{正确理解} \\
\midrule
KF 和批量最小二乘是两回事 & LG 下 KF $\equiv$ 批量 MAP 的递归，非近似（\cref{thm:kf-map}）\\
EKF 给出最优估计 & 仅 LG 的 KF 最优；EKF 是单次线性化下的近似 MAP，可有偏、不一致（\cref{prop:ekf}）\\
高斯过非线性还是高斯 & 一般不是；后验偏斜、均值 $\ne$ 众数（\cref{sec:nonlin-pdf}，\cref{ex:three-methods}）\\
线性化一次就够 & 离工作点远时一阶泰勒失效；优化/IEKF 每步重线性化更稳 \\
UKF/CKF 是 EKF 的“升级版” & 是三种不同逼近策略（泰勒/sigma 点/数值积分），无线性递进（\cref{tab:filter-taxo}）\\
NEES 越小越好 & 判据是 $\mathbb{E}[\text{NEES}]=N$；过自信($>$)与欠自信($<$)都是病（\cref{def:nees-nis}）\\
偏置不大可以忽略 & 未补偿偏置使 KF 系统性过自信、效应随时间无界增长（\cref{thm:bias-overconfident}）\\
姿态直接进 KF 状态即可 & 四元数过参数化致协方差奇异、欧拉角万向锁；用 ESKF 误差状态（\cref{sec:eskf}）\\
增益越大越好 & 增益按先验/观测精度折中；观测噪声大时应\emph{小}增益、信赖预测 \\
鲁棒核只是工程补丁 & 等价于协方差逆 Wishart 先验下的 MAP（\cref{thm:robust-iw}）\\
滤波过时了 & 高频惯性融合、嵌入式仍首选滤波；主流是滤波 $+$ 优化混合（\cref{sec:filter-vs-opt}）\\
$\hat{\mathbf{P}}=(\mathbf{I}-\mathbf{K}\mathbf{C})\check{\mathbf{P}}$ 总没问题 & 数值上易失对称正定，应用 Joseph 形式或平方根滤波（\cref{sec:kf} 数值陷阱）\\
粒子滤波每步都该重采样 & 否；每步重采样加剧样本贫化，应按 $N_{\text{eff}}<M/2$ 触发（\cref{eq:pf-neff}）\\
InEKF 只是“在李群上跑 EKF” & 关键是群仿射 $+$ 不变误差使线性化传播与状态无关 $\Rightarrow$ 一致（\cref{paper:inekf}）\\
\bottomrule
\end{tabular}
\end{table}

% ════════════════════════════════════════════════════════════════
\section*{本章小结}
\addcontentsline{toc}{section}{本章小结}

\begin{table}[H]\centering\small
\caption{符号表}
\begin{tabular}{lll}
\toprule
符号 & 含义 & 首次出现 \\
\midrule
$\check{\mathbf{x}}_k,\hat{\mathbf{x}}_k$ & 预测 / 校正（先验 / 后验）状态 & \cref{eq:kf-predict} \\
$\check{\mathbf{P}}_k,\hat{\mathbf{P}}_k$ & 预测 / 校正协方差 & \cref{eq:kf-predict} \\
$\mathbf{A}_{k-1},\mathbf{C}_k$ & 线性转移 / 观测矩阵 & \cref{eq:lg-model} \\
$\mathbf{F},\mathbf{H}$ & 非线性运动 / 观测雅可比 & \cref{eq:ekf-lin-f,eq:ekf-lin-h} \\
$\boldsymbol{\Sigma}_{\mathbf{w}},\boldsymbol{\Sigma}_{\mathbf{v}}$ & 过程 / 观测噪声协方差（三书字母互异，本书统一） & \cref{eq:lg-model} \\
$\mathbf{K}_k,\ \mathbf{S}_k$ & 卡尔曼增益 / 新息协方差 & \cref{eq:kf-gain} \\
$\boldsymbol{\Lambda}=\boldsymbol{\Sigma}^{-1},\ \boldsymbol{\eta}=\boldsymbol{\Lambda}\mathbf{x}$ & 信息矩阵 / 信息向量（信息滤波 IF） & \cref{eq:kf-info-P,eq:if-predict} \\
$\mathbf{P}$（稳态）, DARE & 稳态协方差 / 离散代数 Riccati 方程 & \cref{eq:dare} \\
$\mathbf{x}_i,\alpha_i,\kappa$ & sigma 点 / 权 / 调参 & \cref{eq:sigma-points} \\
$w_{k,m},N_{\text{eff}}$ & 粒子权重 / 有效样本数 & \cref{eq:pf-weight,eq:pf-neff} \\
$\epsilon_{\text{nees}},\epsilon_{\text{nis}}$ & 归一化估计/新息误差平方 & \cref{def:nees-nis} \\
$\mathbf{x},\delta\mathbf{x},\mathbf{x}_t$ & 标称 / 误差 / 真值态 & \cref{def:eskf-states} \\
$\delta\boldsymbol{\theta}$ & 姿态误差（$3$ 参数，右扰动 $\mathbf{R}\,\mathrm{Exp}(\delta\boldsymbol{\theta})$） & \cref{eq:eskf-states} \\
$\mathbf{F}_{\mathbf{x}},\mathbf{F}_{\mathbf{i}},\mathbf{Q}_{\mathbf{i}}$ & 误差转移 / 扰动雅可比 / 脉冲协方差 & \cref{eq:eskf-Fx} \\
$\mathbf{X}_{\delta\mathbf{x}},\mathbf{Q}_{\delta\boldsymbol{\theta}},\mathbf{G}$ & 合成 / 四元数合成 / 复位雅可比 & \cref{eq:eskf-H,eq:eskf-Qdtheta,eq:eskf-reset} \\
$\rho(u),\mathbf{Y}_i$ & 鲁棒核 / IRLS 膨胀协方差 & \cref{eq:robust-kernels,eq:irls-weight} \\
$k,w,n$ & RANSAC 迭代数 / 内点率 / 子集大小 & \cref{eq:eskf-ransac} \\
$\mathbf{S}_{y,k}$ & 拖尾窗新息样本协方差（自适应估计） & \cref{eq:eskf-adapcov} \\
\bottomrule
\end{tabular}
\end{table}

\begin{table}[H]\centering\small
\caption{定理速查表}
\begin{tabular}{lll}
\toprule
名称 & 一句话说明（能做什么） & 出处 \\
\midrule
贝叶斯滤波 & 预测–校正递归（精确，一般不可实现，须近似） & \cref{thm:bayes-filter} \\
卡尔曼滤波 & LG 预测/校正 + 增益；MAP 配方与高斯条件化两路同解 & \cref{thm:kf} \\
信息形式 / 信息滤波 IF & 校正 $=$ 信息相加（免大逆），可表全局无知；SEIF 稀疏化 & \cref{eq:kf-info-P,prop:info-filter} \\
KF $\equiv$ MAP & LG 下滤波 $=$ 用至今全部数据的批量 MAP，非近似 & \cref{thm:kf-map} \\
稳态 KF / DARE & 常值增益；唯一半正定解需可镇定 $+$ 可检测四条件 & \cref{eq:dare} \\
Kalman--Bucy & 连续时间 LTV，两条耦合 ODE $+$ 增益 $\mathbf{K}=\hat{\mathbf{P}}\mathbf{C}^\top\mathbf{R}^{-1}$ & \cref{eq:kb-mean,eq:kb-cov} \\
RTS 平滑 & 前向 KF $+$ 后向修正，用未来数据精修过去 & \cref{eq:rts} \\
直方图 / 离散贝叶斯 & 离散网格的预测–校正，最朴素可实现、能表多峰 & \cref{eq:histogram-filter} \\
EKF / 广义高斯滤波 & 线性化 / 高斯条件化得近似后验 & \cref{prop:ekf,prop:ggf} \\
三法过非线性 & 蒙特卡洛 / 线性化（过自信）/ sigma 点（捕真均值） & \cref{ex:three-methods} \\
SPKF/UKF & 无迹变换预测/校正逐式；轻非线性退回 EKF 增益 & \cref{algo:spkf-predict,prop:spkf-ekf} \\
缩放无迹变换 & 三参数 $\alpha,\beta,\kappa$ 解耦展开半径与高阶矩补偿 & \cref{def:scaled-ut} \\
CKF $=$ Scaled UT 特例 & $(\alpha,\beta,\kappa)=(1,0,0)$；$2n$ 点全正权、免调参 & \cref{prop:ckf-ukf,der:ckf-rule} \\
高斯滤波 $=$ 求积规则 & EKF/UKF/CKF/GHKF 只差数值积分规则 & \cref{tab:gauss-quadrature} \\
ISPKF & 绕工作点迭代 sigma 点，趋后验均值 & \cref{algo:ispkf} \\
Bell--Cathey & IEKF $=$ 对单步 MAP 做高斯牛顿（趋众数） & \cref{thm:bell-cathey} \\
流形 IEKF（IKFoM） & $\boxplus/\boxminus$ $+$ 曲率雅可比 $\mathbf{J}$ 的流形 GN & \cref{eq:ikfom-gn} \\
滤波$\to$优化谱 & EKF$\subset$IEKF$\subset$SWF/MHE$\subset$BA，只差窗口大小 & \cref{eq:iekf-bridge} \\
非线性 RTS（EKS/UKS/CKS） & 后向递推套 RTS，仅换交叉协方差近似 & \cref{tab:nonlinear-rts} \\
RTS 三视角 & 前向后向 / 双滤波器 / 批量（消元$=$KF、回代$=$RTS） & \cref{thm:rts-views} \\
平方根 sigma 点（SR-UKF） & 传播 Cholesky 因子，QR $+$ cholupdate 保稳定 & \cref{eq:srukf-qr} \\
Kalman 族全景 & 四正交轴：高斯/线性化/几何/迭代-平滑 & \cref{tab:kalman-panorama} \\
粒子滤波 & SIR 三步 $+$ proposal $+$ $N_{\text{eff}}$ 按需重采样；RBPF/FastSLAM & \cref{algo:pf,eq:pf-neff} \\
滤波器谱系 & IEKF$\to$众数，ISPKF$\to$均值，EKF/UKF 无明确关系 & \cref{tab:filter-taxo} \\
偏置使 KF 过自信 & 未补偿偏置令协方差系统性偏小、随时间无界增长 & \cref{thm:bias-overconfident} \\
NEES/NIS & 诊断一致性（$\mathbb{E}=N$ 或 $M$），NIS 可在线无真值 & \cref{def:nees-nis} \\
ESKF 误差动力学 & 误差态线速度/角误差线性化（右扰动） & \cref{thm:eskf-error-kin} \\
ESKF 周期 & 标称积分 $+$ 误差滤波 $+$ 注入/复位 & \cref{algo:eskf} \\
互补滤波 & 陀螺高通 $+$ 加计低通的确定性姿态估计（无协方差） & \cref{eq:complementary-filter} \\
InEKF / 群仿射 & 不变误差使线性化传播与状态无关 $\Rightarrow$ 一致；需 $\mathrm{SE}_2(3)$ & \cref{paper:inekf} \\
鲁棒 = 逆 Wishart MAP & 鲁棒核源自对协方差的先验，等价 IRLS & \cref{thm:robust-iw} \\
RANSAC 迭代数 & 由内点率与目标置信反解抽样次数 & \cref{eq:eskf-ransac} \\
协方差估计 & 监督式（真值样本协方差）/ 自适应（拖尾窗反解） & \cref{eq:eskf-supcov,eq:eskf-adapcov} \\
\bottomrule
\end{tabular}
\end{table}

\begin{table}[H]\centering\small
\caption{知识点总表}
\begin{tabularx}{\linewidth}{clLcl}
\toprule
\# & 知识点 & 核心要点 & 难度 & 脉络层 \\
\midrule
1 & 贝叶斯滤波 & 预测–校正、精确但不可实现 & $\bigstar\bigstar\bigstar$ & 骨干 \\
2 & 卡尔曼滤波 & 配方/条件化两路、增益折中、信息形式 IF、稳态 DARE & $\bigstar\bigstar\bigstar$ & 骨干 \\
3 & KF $\equiv$ 批量 MAP & 与 \cref{ch:slam_est} 缝合；RTS 平滑、Kalman--Bucy & $\bigstar\bigstar\bigstar\bigstar$ & 骨干 \\
4 & EKF / 广义高斯 & 线性化雅可比、三宗罪 & $\bigstar\bigstar\bigstar$ & 骨干 \\
5 & 贝叶斯滤波谱系 & 直方图/MC/线性化/sigma 点；IEKF/ISPKF/PF（$N_{\text{eff}}$）/SWF & $\bigstar\bigstar\bigstar\bigstar$ & 次优 \\
6 & 一致性与偏置 & NEES/NIS、过自信、状态增广可观测性 & $\bigstar\bigstar\bigstar\bigstar$ & 次优 \\
7 & ESKF / 右扰动 & 三态、误差动力学、注入/复位、左右扰动；互补滤波、InEKF & $\bigstar\bigstar\bigstar\bigstar$ & 骨干 \\
8 & 滤波 vs 优化 + 鲁棒 & 四局限、稀疏对偶、鲁棒核 = 逆 Wishart MAP & $\bigstar\bigstar\bigstar$ & 骨干 \\
9 & Sigma 点族收口 & 缩放 UT（$\alpha,\beta,\kappa$）、CKF=UT 特例、GHKF/稀疏网格、求积统一观 & $\bigstar\bigstar\bigstar$ & 次优 \\
10 & 迭代变体收口 & Bell--Cathey（IEKF=GN）、LM/线搜索、流形 IKFoM、滤波$\to$优化谱 & $\bigstar\bigstar\bigstar\bigstar$ & 骨干 \\
11 & 平滑器族收口 & 非线性 RTS（EKS/UKS/CKS）、三等价视角、1D 数值示例 & $\bigstar\bigstar\bigstar$ & 次优 \\
12 & 平方根 sigma 点 + 全景 & SR-UKF/SR-CKF、四正交轴、Kalman 族全景对比表 & $\bigstar\bigstar\bigstar$ & 骨干 \\
\bottomrule
\end{tabularx}
\end{table}

\paragraph{延伸阅读。}
\begin{itemize}
  \item Barfoot《State Estimation for Robotics》\cite{barfoot2024state} 第 3 章（递归 LG/KF/RTS/信息滤波）、第 4 章（贝叶斯滤波/EKF/广义高斯/IEKF/sigma 点·UKF·ISPKF/粒子滤波/滤波器谱系/批量 MAP·GN/Laplace 协方差/ML 偏差/SWF）、第 5 章（无偏与一致/NEES·NIS/偏置与可观测性/RANSAC/M 估计·IRLS/协方差估计·逆 Wishart）：本章理论骨架与一致性/鲁棒内容的来源。
  \item Solà《Quaternion kinematics for the error-state Kalman filter》\cite{sola2017quaternion}（arXiv:1711.02508）：四元数运动学与 ESKF 三态、预测/校正、注入/复位、局部 vs 全局角误差对照的权威推导（本章 ESKF 一节的完整出处）。
  \item Solà 等《A micro Lie theory for state estimation in robotics》\cite{sola2018micro}：右雅可比与流形上求导，ESKF 雅可比的现代统一记号。
  \item 视觉 SLAM 十四讲\cite{gaoxiang2019slam14} 第 9 章：KF/EKF 的直觉与 MAP 配方法推导、EKF-SLAM 四局限、H 稀疏·Schur·边缘化·鲁棒核、“滤波 vs 优化”叙事。
  \item Thrun, Burgard \& Fox《Probabilistic Robotics》\cite{thrun2005probabilistic}：直方图/网格滤波、粒子滤波（重要性采样、$N_{\text{eff}}$、重采样）、信息滤波与 SEIF、栅格定位的标准教科书出处。
  \item Särkkä《Bayesian Filtering and Smoothing》\cite{sarkka2013bayesian}：贝叶斯滤波/平滑的统一处理，粒子滤波退化与有效样本数判据的清晰推导；其 General Gaussian Filtering 章是“高斯滤波 $=$ 求积规则”统一观（\cref{tab:gauss-quadrature}）与非线性 RTS（EKS/UKS/CKS，\cref{sec:rts-family}）的权威出处。
  \item Bell \& Cathey《The iterated Kalman filter update as a Gauss--Newton method》\cite{bell1993iterated}（IEEE TAC 1993）：IEKF $=$ 对单步 MAP 做高斯牛顿（\cref{thm:bell-cathey}）的奠基论文，Bell 1994 进一步证明迭代平滑 $=$ 整条轨迹 GN。
  \item Julier《The scaled unscented transformation》\cite{julier2002scaled}（ACC 2002）与 van der Merwe《Sigma-Point Kalman Filters for Probabilistic Inference》\cite{vandermerwe2004sigma}（博士论文）：缩放无迹变换三参数（\cref{def:scaled-ut}）、SR-UKF（\cref{sec:srukf-panorama}）、加性/增广 UKF 的权威推导与默认参数来源。
  \item Arasaratnam \& Haykin《Cubature Kalman Filters》\cite{arasaratnam2009cubature}（IEEE TAC 2009）与 Arasaratnam \& Haykin《Cubature Kalman Smoothers》\cite{arasaratnam2011smoother}（Automatica 2011）：三阶球面径向容积规则（\cref{der:ckf-rule}）、CKF 全正权重、CKS 平滑的出处。
  \item Ito \& Xiong《Gaussian filters for nonlinear filtering problems》\cite{ito2000gaussian}（IEEE TAC 2000）与 Jia--Xin--Cheng 稀疏高斯--厄米求积\cite{jia2012sparse}（Automatica 2012）：GHKF 与 Smolyak 稀疏网格（\cref{tab:gauss-quadrature}）的来源。
  \item Särkkä《Unscented Rauch--Tung--Striebel Smoother》\cite{sarkka2008unscented}（IEEE TAC 2008）：UKS/URTSS 的原始推导（\cref{tab:nonlinear-rts}）。
  \item He, Xu \& Zhang《Embedding manifold structures into Kalman filters》（IKFoM）\cite{he2021ikfom}（arXiv:2102.03804）：流形迭代 EKF 的 $\boxplus/\boxminus$ 与曲率雅可比（\cref{eq:ikfom-gn}），FAST-LIO/FAST-LIO2 核心。
  \item Xu 等《FAST-LIO》\cite{xu2021fastlio} 与《FAST-LIO2》\cite{xu2022fastlio2}：流形迭代 EKF 在紧耦合 LiDAR-惯性里程计的工程落地，高效卡尔曼增益（复杂度随状态维而非测量维）的来源。
  \item Barrau \& Bonnabel《The Invariant Extended Kalman Filter as a Stable Observer》\cite{barrau2017invariant}：群仿射系统、(右)不变误差、误差自治（log-linear）性质与一致性的奠基论文，引入扩展位姿群 $\mathrm{SE}_2(3)$。
  \item Mahony, Hamel \& Pflimlin《Nonlinear Complementary Filters on the Special Orthogonal Group》\cite{mahony2008nonlinear}：直接定义在 $\mathrm{SO}(3)$ 上的显式互补滤波（叉积误差 $+$ PI 反馈 $+$ 在线偏置估计）。
  \item Madgwick《An efficient orientation filter for inertial/MARG sensor arrays》\cite{madgwick2011estimation}：低算力四元数梯度下降姿态滤波，开源 IMU 姿态估计的常用实现。
  \item SLAM Handbook\cite{carlone2026handbook} 第 1 章：平滑 vs 滤波（信息矩阵稀疏 vs 稠密）的现代论述与因子图视角。
  \item Barron《A General and Adaptive Robust Loss Function》\cite{barron2019general}（arXiv:1701.03077）——把 Huber/Cauchy/Geman--McClure/Welsch 统一为单参数族并在估计中自适应学出形状，免去手调核宽。
  \item Yang 等《Graduated Non-Convexity for Robust Spatial Perception》\cite{yang2020gnc}（arXiv:1909.08605）——逐步非凸化：从凸代价平滑过渡到目标鲁棒核，对抗高比例离群点、避免坏局部极小，已成位姿图/点云配准标配。
\end{itemize}

\section*{本章与后续章节的关系}
\begin{table}[H]\centering\small
\begin{tabularx}{\linewidth}{lLL}
\toprule
后续章节 & 与本章的关系 & 本章哪个知识点为其铺垫 \\
\midrule
\cref{ch:imu} IMU 模型 & 给 ESKF 填入具体运动学与噪声标定；互补滤波做初始对准 & \cref{algo:eskf}、\cref{eq:eskf-Fx}、\cref{eq:complementary-filter} \\
\cref{ch:vio} VIO & ESKF 前端 + 因子图后端、边缘化老帧；InEKF 一致性去向 & \cref{sec:eskf}、SWF（\cref{par:swf}）、\cref{paper:inekf} \\
\cref{ch:lie} 李群 & 提供 InEKF 所需 $\mathrm{SE}_2(3)$ 的正式定义（本章只给结构要点） & \cref{paper:inekf} \\
\cref{ch:nlopt} 非线性优化 & 批量 MAP 的 GN/LM 解、鲁棒核/IRLS、$\sqrt{}$SAM（平方根滤波同源） & \cref{thm:kf-map}、\cref{eq:irls-weight} \\
\cref{ch:slam_est} 状态估计 & 信息矩阵稀疏、舒尔补、协方差对偶 & \cref{thm:kf}、\cref{sec:filter-vs-opt} \\
\bottomrule
\end{tabularx}
\end{table}

\section*{故障排查手册}
\begin{table}[H]\centering\small
\begin{tabular}{p{0.24\textwidth}p{0.30\textwidth}p{0.36\textwidth}}
\toprule
症状 & 可能原因 & 排查步骤 / 相关节 \\
\midrule
EKF 发散 / 估计跑飞 & 强非线性、工作点偏差、初值差 & 改 IEKF/ESKF；给好初值；减小步长（\cref{sec:ekf}） \\
协方差越来越乐观（不一致） & 单次线性化、$\boldsymbol{\Sigma}$ 设小、偏置未补偿 & 跑 NEES/NIS；Joseph 更新；状态增广估偏置（\cref{sec:consistency}） \\
$\hat{\mathbf{P}}$ 出现负特征值 / NaN & $(\mathbf{I}-\mathbf{K}\mathbf{C})\check{\mathbf{P}}$ 失对称 & 换 Joseph 形式 $\hat{\mathbf{P}}=(\mathbf{I}-\mathbf{K}\mathbf{C})\check{\mathbf{P}}(\cdot)^\top+\mathbf{K}\boldsymbol{\Sigma}_{\mathbf{v}}\mathbf{K}^\top$（\cref{sec:kf}） \\
UKF/CKF 跑几分钟后 Cholesky 失败 & 浮点累积 $+$ 中心权 $W_0^c<0$ $+$ 校正减法 & 对称化/加抖动；根治用 SR-UKF 传播平方根因子（\cref{sec:srukf-panorama}） \\
强非线性观测下 EKF 推不动估计 & 先验点斜率小、只线性化一次 & 改 IEKF 重线性化（\cref{thm:bell-cathey}）；初值差时加 LM/线搜索 \\
平滑后反比滤波差 / 残差异常 & 观测含延迟却按接收时戳套 RTS & 状态增广对齐延迟分量或用固定滞后平滑（\cref{sec:rts-family}） \\
姿态协方差奇异 / 需反复归一化 & 四元数直接进状态 & 改 ESKF 误差状态 $\delta\boldsymbol{\theta}$（\cref{sec:eskf}） \\
大地图越跑越慢 & 滤波边缘化致信息阵稠密（fill-in） & 改因子图优化（保稀疏）或 SWF（\cref{sec:filter-vs-opt}） \\
增益异常、新息被忽略/过信 & $\boldsymbol{\Sigma}_{\mathbf{v}}$ 设定不当 & 按传感器真实噪声标定，或自适应协方差（\cref{sec:consistency}） \\
ESKF 长期漂移 & 复位 $\mathbf{G}\approx\mathbf{I}$ 过简、未补偏置 & 保留 \cref{eq:eskf-reset} 姿态修正块；增广估偏置 \\
少数误匹配毁掉整解 & 二范数对离群点敏感 & 加 RANSAC 预处理 + 鲁棒核/IRLS（\cref{eq:robust-kernels}） \\
\bottomrule
\end{tabular}
\end{table}

\paragraph{研究实践建议。}
给新手：先把 \cref{ex:kf-1d} 的一维 KF 手算 + 代码（\cref{thm:kf} 代码框）跑通，再用 NEES/NIS 检验它是否一致；上 ESKF 前务必先在纯标称态积分（无校正）下看“蒙眼漂移”是否合理。给有经验者：实现 ESKF 时优先 Joseph 协方差更新 + 局部角误差（右扰动），把复位雅可比 $\mathbf{G}$ 的姿态项保留以压长期漂移；做大场景时，用 SWF/因子图替代纯滤波，并把鲁棒核的核宽 $\delta$ 与一致性检验联动调参。

% ════════════════════════════════════════════════════════════════
\section*{附录：本章推导补遗}
\label{sec:eskf-appendix}
\addcontentsline{toc}{section}{附录：本章推导补遗}

\begin{derivation}[KF $\equiv$ 批量 MAP 的递归（\cref{thm:kf-map}）]
用等价先验 $\{\hat{\mathbf{x}}_{k-1},\hat{\mathbf{P}}_{k-1}\}$ 代替全部历史（马尔可夫性，LG 下精确），构造小批量
\[
  \mathbf{z}=\begin{bmatrix}\hat{\mathbf{x}}_{k-1}\\\mathbf{u}_k\\\mathbf{z}_k\end{bmatrix},\quad
  \mathbf{H}=\begin{bmatrix}\mathbf{I}&\\-\mathbf{A}_{k-1}&\mathbf{I}\\&\mathbf{C}_k\end{bmatrix},\quad
  \mathbf{W}=\operatorname{diag}(\hat{\mathbf{P}}_{k-1},\boldsymbol{\Sigma}_{\mathbf{w}},\boldsymbol{\Sigma}_{\mathbf{v}}).
\]
其正规方程 $(\mathbf{H}^\top\mathbf{W}^{-1}\mathbf{H})[\hat{\mathbf{x}}_{k-1}';\hat{\mathbf{x}}_k]=\mathbf{H}^\top\mathbf{W}^{-1}\mathbf{z}$ 的 $(2,2)$ 块为 $\boldsymbol{\Sigma}_{\mathbf{w}}^{-1}+\mathbf{C}_k^\top\boldsymbol{\Sigma}_{\mathbf{v}}^{-1}\mathbf{C}_k$，$(1,1)$ 块为 $\hat{\mathbf{P}}_{k-1}^{-1}+\mathbf{A}_{k-1}^\top\boldsymbol{\Sigma}_{\mathbf{w}}^{-1}\mathbf{A}_{k-1}$。用舒尔补消去不需要的 $\hat{\mathbf{x}}_{k-1}'$（被未来数据更新的过去），$\hat{\mathbf{x}}_k$ 的系数化为
\[
  \underbrace{\big(\boldsymbol{\Sigma}_{\mathbf{w}}+\mathbf{A}_{k-1}\hat{\mathbf{P}}_{k-1}\mathbf{A}_{k-1}^\top\big)^{-1}}_{=\,\check{\mathbf{P}}_k^{-1}\ (\text{SMW})}+\mathbf{C}_k^\top\boldsymbol{\Sigma}_{\mathbf{v}}^{-1}\mathbf{C}_k,
\]
即信息形式 \cref{eq:kf-info-P}，再经矩阵求逆引理还原为协方差形式 \cref{eq:kf-predict,eq:kf-gain,eq:kf-correct}。前向递推这一过程即 KF；附加后向修正即 RTS。\hfill$\square$
\end{derivation}

\begin{derivation}[另一条路线：把批量解直接因式分解为前向 $\times$ 后向（\cref{eq:eskf-fwd-bwd}）]\label{der:factor-batch}
\cref{thm:kf-map} 走“小三块 MAP”。这里给 Barfoot\cite{barfoot2024state} 等价但更显式的\emph{因式分解}路线，直接坐实 \cref{eq:eskf-fwd-bwd} 的“前向 $\times$ 后向”。把按时刻重排的批量 $\mathbf{H}$ 分成 $12$ 块（仅 $6$ 块非零），按“$0\!:\!k\!-\!1$ / $k$ / $k\!+\!1\!:\!K$”三段把正规方程写成块三对角
\[
  \begin{bmatrix}\mathbf{L}_{11}&\mathbf{L}_{12}&\\ \mathbf{L}_{12}^\top&\mathbf{L}_{22}&\mathbf{L}_{32}^\top\\ &\mathbf{L}_{32}&\mathbf{L}_{33}\end{bmatrix}
  \begin{bmatrix}\hat{\mathbf{x}}_{0:k-1}\\\hat{\mathbf{x}}_k\\\hat{\mathbf{x}}_{k+1:K}\end{bmatrix}
  =\begin{bmatrix}\mathbf{r}_1\\\mathbf{r}_2\\\mathbf{r}_3\end{bmatrix}.
\]
左乘行操作矩阵 $\big[\,\mathbf{I};\ -\mathbf{L}_{12}^\top\mathbf{L}_{11}^{-1}\ \ \mathbf{I}\ \ -\mathbf{L}_{32}^\top\mathbf{L}_{33}^{-1};\ \mathbf{I}\,\big]$（初等行变换，不改变解），把 $\hat{\mathbf{x}}_{0:k-1}$ 与 $\hat{\mathbf{x}}_{k+1:K}$ 同时\emph{边缘化}，孤立出 $\hat{\mathbf{x}}_k$：
\[
  \underbrace{\big(\mathbf{L}_{22}-\mathbf{L}_{12}^\top\mathbf{L}_{11}^{-1}\mathbf{L}_{12}-\mathbf{L}_{32}^\top\mathbf{L}_{33}^{-1}\mathbf{L}_{32}\big)}_{\hat{\mathbf{P}}_k^{-1}}\hat{\mathbf{x}}_k
  =\underbrace{\big(\mathbf{r}_2-\mathbf{L}_{12}^\top\mathbf{L}_{11}^{-1}\mathbf{r}_1-\mathbf{L}_{32}^\top\mathbf{L}_{33}^{-1}\mathbf{r}_3\big)}_{\mathbf{q}_k}.
\]
回代各 $\mathbf{L}_{ij}$ 块、对每段用 SMW，$\hat{\mathbf{P}}_k^{-1}$ 与 $\mathbf{q}_k$ \emph{恰好}各裂成只依赖 $\le k$ 数据的“前向”项与只依赖 $>k$ 数据的“后向”项：
\[
  \hat{\mathbf{P}}_k^{-1}=\hat{\mathbf{P}}_{k,f}^{-1}+\hat{\mathbf{P}}_{k,b}^{-1},\qquad
  \hat{\mathbf{P}}_k^{-1}\hat{\mathbf{x}}_k=\hat{\mathbf{P}}_{k,f}^{-1}\hat{\mathbf{x}}_{k,f}+\hat{\mathbf{P}}_{k,b}^{-1}\hat{\mathbf{x}}_{k,b},
\]
正是\emph{两个高斯 PDF 的归一化乘积}（\cref{ch:slam_est} 的“独立高斯相乘 $=$ 信息相加”），对应 \cref{eq:eskf-fwd-bwd} 的 $p(\mathbf{x}_k\mid\le k)\,p(\mathbf{x}_k\mid>k)$。前向估计器 $\{\hat{\mathbf{x}}_{k,f},\hat{\mathbf{P}}_{k,f}\}$ 递归化即 KF；后向估计器 $\{\hat{\mathbf{x}}_{k,b},\hat{\mathbf{P}}_{k,b}\}$ \emph{只}用未来数据（注意它\emph{不是}在修正前向，与 RTS 后向\emph{修正}不同），二者相乘即完整平滑后验。这条路线比 \cref{thm:kf-map} 更直白地暴露了“滤波 $=$ 前向因子”。\hfill$\square$
\end{derivation}

\begin{derivation}[KF 无偏与一致的数学归纳证（\cref{thm:kf} 的 BLUE/CRLB 根基）]\label{der:kf-unbiased-consistent}
由误差 $\check{\mathbf{e}}_k=\check{\mathbf{x}}_k-\mathbf{x}_k$、$\hat{\mathbf{e}}_k=\hat{\mathbf{x}}_k-\mathbf{x}_k$，KF 递推减真实动力学得\emph{误差动力学}（即 \cref{thm:bias-overconfident} 令偏置 $\bar{\mathbf{u}}=\bar{\mathbf{z}}=\mathbf{0}$ 的洁净特例）
\[
  \check{\mathbf{e}}_k=\mathbf{A}_{k-1}\hat{\mathbf{e}}_{k-1}-\mathbf{w}_k,\qquad
  \hat{\mathbf{e}}_k=(\mathbf{I}-\mathbf{K}_k\mathbf{C}_k)\check{\mathbf{e}}_k+\mathbf{K}_k\mathbf{v}_k.
\]
\textbf{无偏}（$\mathbb{E}[\hat{\mathbf{e}}_k]=\mathbf{0}$），归纳：$k=0$ 由初始假设 $\mathbb{E}[\hat{\mathbf{e}}_0]=\mathbf{0}$ 成立；设 $\mathbb{E}[\hat{\mathbf{e}}_{k-1}]=\mathbf{0}$，取期望（$\mathbb{E}[\mathbf{w}_k]=\mathbb{E}[\mathbf{v}_k]=\mathbf{0}$）
\[
  \mathbb{E}[\check{\mathbf{e}}_k]=\mathbf{A}_{k-1}\underbrace{\mathbb{E}[\hat{\mathbf{e}}_{k-1}]}_{\mathbf{0}}-\underbrace{\mathbb{E}[\mathbf{w}_k]}_{\mathbf{0}}=\mathbf{0},\quad
  \mathbb{E}[\hat{\mathbf{e}}_k]=(\mathbf{I}-\mathbf{K}_k\mathbf{C}_k)\underbrace{\mathbb{E}[\check{\mathbf{e}}_k]}_{\mathbf{0}}+\mathbf{K}_k\underbrace{\mathbb{E}[\mathbf{v}_k]}_{\mathbf{0}}=\mathbf{0},
\]
故对一切 $k$ 无偏。\textbf{一致}（$\mathbb{E}[\hat{\mathbf{e}}_k\hat{\mathbf{e}}_k^\top]=\hat{\mathbf{P}}_k$），归纳：$k=0$ 由 $\mathbb{E}[\hat{\mathbf{e}}_0\hat{\mathbf{e}}_0^\top]=\hat{\mathbf{P}}_0$ 成立；设 $\mathbb{E}[\hat{\mathbf{e}}_{k-1}\hat{\mathbf{e}}_{k-1}^\top]=\hat{\mathbf{P}}_{k-1}$。预测误差外积（交叉项因 $\mathbf{w}_k$ 与 $\hat{\mathbf{e}}_{k-1}$ 独立而为零）
\[
  \mathbb{E}[\check{\mathbf{e}}_k\check{\mathbf{e}}_k^\top]=\mathbf{A}_{k-1}\underbrace{\mathbb{E}[\hat{\mathbf{e}}_{k-1}\hat{\mathbf{e}}_{k-1}^\top]}_{\hat{\mathbf{P}}_{k-1}}\mathbf{A}_{k-1}^\top+\underbrace{\mathbb{E}[\mathbf{w}_k\mathbf{w}_k^\top]}_{\boldsymbol{\Sigma}_{\mathbf{w}}}=\check{\mathbf{P}}_k,
\]
恰为 \cref{eq:kf-predict}。校正误差外积（交叉项同理为零）给 Joseph 形，再用 $\mathbf{K}_k=\hat{\mathbf{P}}_k\mathbf{C}_k^\top\boldsymbol{\Sigma}_{\mathbf{v}}^{-1}$（即 $\hat{\mathbf{P}}_k\mathbf{C}_k^\top=\mathbf{K}_k\boldsymbol{\Sigma}_{\mathbf{v}}$）坍缩回规范形：
\begin{align*}
  \mathbb{E}[\hat{\mathbf{e}}_k\hat{\mathbf{e}}_k^\top]
  &=(\mathbf{I}-\mathbf{K}_k\mathbf{C}_k)\check{\mathbf{P}}_k(\mathbf{I}-\mathbf{K}_k\mathbf{C}_k)^\top+\mathbf{K}_k\boldsymbol{\Sigma}_{\mathbf{v}}\mathbf{K}_k^\top\\
  &=(\mathbf{I}-\mathbf{K}_k\mathbf{C}_k)\check{\mathbf{P}}_k-\underbrace{(\mathbf{I}-\mathbf{K}_k\mathbf{C}_k)\check{\mathbf{P}}_k\mathbf{C}_k^\top}_{=\,\hat{\mathbf{P}}_k\mathbf{C}_k^\top}\mathbf{K}_k^\top+\mathbf{K}_k\boldsymbol{\Sigma}_{\mathbf{v}}\mathbf{K}_k^\top
  =\hat{\mathbf{P}}_k\underbrace{-\hat{\mathbf{P}}_k\mathbf{C}_k^\top\mathbf{K}_k^\top+\mathbf{K}_k\boldsymbol{\Sigma}_{\mathbf{v}}\mathbf{K}_k^\top}_{=\,\mathbf{0}}=\hat{\mathbf{P}}_k,
\end{align*}
恰为 \cref{eq:kf-correct}。故对一切 $k$ 一致——KF 报告的 $\hat{\mathbf{P}}_k$ \emph{正是}真实误差协方差，这是“KF 最优 / BLUE / 达 CRLB”的根基。（此处期望对随机变量\emph{全体实现}取集合平均、非时间平均：单条轨迹的误差并不为零。）\hfill$\square$
\end{derivation}

\begin{derivation}[RTS 规范形：从块三对角 Cholesky 后向回代到 \cref{eq:rts,eq:rts-cov}（多次 SMW）]
前一条已把\emph{前向}消元化简为 KF。\emph{后向}回代给出平滑量，化成规范 RTS 需\emph{多次}用矩阵求逆引理（SMW）。记前向滤波量 $\hat{\mathbf{P}}_{k-1,f}$、$\check{\mathbf{P}}_{k,f}=\mathbf{A}_{k-1}\hat{\mathbf{P}}_{k-1,f}\mathbf{A}_{k-1}^\top+\boldsymbol{\Sigma}_{\mathbf{w}}$，把块三对角信息矩阵 $\mathbf{H}^\top\mathbf{W}^{-1}\mathbf{H}$ 的下双对角 Cholesky 因子记为 $\mathbf{L}$（对角块 $\mathbf{L}_{k-1}$、次对角块 $\mathbf{L}_{k,k-1}$，构造见 \cref{ch:slam_est}），并定义\emph{平滑增益} $\mathbf{G}_{k-1}=\hat{\mathbf{P}}_{k-1,f}\mathbf{A}_{k-1}^\top\check{\mathbf{P}}_{k,f}^{-1}$。

\textbf{均值（规范）.}\;后向回代式 $\hat{\mathbf{x}}_{k-1}=(\mathbf{L}_{k-1}\mathbf{L}_{k-1}^\top)^{-1}\mathbf{L}_{k-1}(-\mathbf{L}_{k,k-1}^\top\hat{\mathbf{x}}_k+\mathbf{d}_{k-1})$ 代入因子块、用一次 SMW
\[
  \big(\hat{\mathbf{P}}_{k-1,f}^{-1}+\mathbf{A}_{k-1}^\top\boldsymbol{\Sigma}_{\mathbf{w}}^{-1}\mathbf{A}_{k-1}\big)^{-1}\mathbf{A}_{k-1}^\top\boldsymbol{\Sigma}_{\mathbf{w}}^{-1}
  =\hat{\mathbf{P}}_{k-1,f}\mathbf{A}_{k-1}^\top\check{\mathbf{P}}_{k,f}^{-1}=\mathbf{G}_{k-1},
\]
整理出规范均值修正 $\hat{\mathbf{x}}_{k-1}=\hat{\mathbf{x}}_{k-1,f}+\mathbf{G}_{k-1}(\hat{\mathbf{x}}_k-\check{\mathbf{x}}_{k,f})$，即 \cref{eq:rts}。

\textbf{协方差（规范）.}\;后向协方差回代式 $\hat{\mathbf{P}}_{k-1}=\mathbf{L}_{k-1}^{-\top}\mathbf{L}_{k-1}^{-1}+(\mathbf{L}_{k,k-1}\mathbf{L}_{k-1}^{-1})^\top\hat{\mathbf{P}}_k(\mathbf{L}_{k,k-1}\mathbf{L}_{k-1}^{-1})$ 需两条 SMW：
\begin{align*}
  \mathbf{L}_{k,k-1}\mathbf{L}_{k-1}^{-1}&=-\big(\hat{\mathbf{P}}_{k-1,f}\mathbf{A}_{k-1}^\top\check{\mathbf{P}}_{k,f}^{-1}\big)^\top=-\mathbf{G}_{k-1}^\top,\\
  \mathbf{L}_{k-1}^{-\top}\mathbf{L}_{k-1}^{-1}&=\big(\hat{\mathbf{P}}_{k-1,f}^{-1}+\mathbf{A}_{k-1}^\top\boldsymbol{\Sigma}_{\mathbf{w}}^{-1}\mathbf{A}_{k-1}\big)^{-1}
  \overset{\text{SMW}}{=}\hat{\mathbf{P}}_{k-1,f}-\mathbf{G}_{k-1}\check{\mathbf{P}}_{k,f}\mathbf{G}_{k-1}^\top.
\end{align*}
代入即
\[
  \hat{\mathbf{P}}_{k-1}=\hat{\mathbf{P}}_{k-1,f}-\mathbf{G}_{k-1}\check{\mathbf{P}}_{k,f}\mathbf{G}_{k-1}^\top+\mathbf{G}_{k-1}\hat{\mathbf{P}}_k\mathbf{G}_{k-1}^\top
  =\hat{\mathbf{P}}_{k-1,f}+\mathbf{G}_{k-1}\big(\hat{\mathbf{P}}_k-\check{\mathbf{P}}_{k,f}\big)\mathbf{G}_{k-1}^\top,
\]
以 $\hat{\mathbf{P}}_K=\hat{\mathbf{P}}_{K,f}$ 起向后迭代，即 \cref{eq:rts-cov}。主对角上下的互协方差块由 $\hat{\mathbf{P}}_{k,k-1}=\hat{\mathbf{P}}_k\mathbf{G}_{k-1}^\top$ 给出（插值额外时刻时需要）\cite{barfoot2024state}。它给每步\emph{平滑}边缘协方差，恒不大于对应滤波协方差（未来数据只增信息）。\hfill$\square$
\end{derivation}

\begin{derivation}[广义高斯滤波 $\to$ 信息形式批量（连接滤波与优化）]
Barfoot\cite{barfoot2024state} 把整轨迹联合密度按舒尔补条件化得后验 $\mathbf{K}=\check{\mathbf{P}}\mathbf{G}^\top(\mathbf{G}\check{\mathbf{P}}\mathbf{G}^\top+\mathbf{R}')^{-1}$、$\hat{\mathbf{P}}=(\mathbf{I}-\mathbf{K}\mathbf{G})\check{\mathbf{P}}$、$\hat{\mathbf{x}}=\check{\mathbf{x}}+\mathbf{K}(\mathbf{y}-\mathbf{y}_{\text{op}}-\mathbf{G}(\check{\mathbf{x}}-\mathbf{x}_{\text{op}}))$，用 SMW 重排为信息形式
\[
  \underbrace{(\mathbf{F}^{-\top}\mathbf{Q}'^{-1}\mathbf{F}^{-1}+\mathbf{G}^\top\mathbf{R}'^{-1}\mathbf{G})}_{\text{块三对角}}\,\delta\mathbf{x}^*=\mathbf{F}^{-\top}\mathbf{Q}'^{-1}(\boldsymbol{\nu}-\mathbf{F}^{-1}\mathbf{x}_{\text{op}})+\mathbf{G}^\top\mathbf{R}'^{-1}(\mathbf{y}-\mathbf{y}_{\text{op}}),
\]
即 $(\mathbf{H}^\top\mathbf{W}^{-1}\mathbf{H})\delta\mathbf{x}^*=\mathbf{H}^\top\mathbf{W}^{-1}\mathbf{e}(\mathbf{x}_{\text{op}})$——\emph{与非线性批量 GN 的正规方程完全相同}。\textbf{MAP 与贝叶斯之别仅在于从 SMW 哪一侧起步}：迭代地在最佳估计均值处重线性化，使两者“均值”一致（同 IEKF $=$ MAP 现象）。这把“滤波/平滑 $=$ 批量 MAP 的稀疏递归”坐实在代数层面。\hfill$\square$
\end{derivation}

\begin{derivation}[ESKF 复位雅可比的姿态块（\cref{eq:eskf-reset}）]
复位要求真值姿态不变 $\mathbf{q}_t^+=\mathbf{q}_t$，即 $\mathbf{q}^+\otimes\delta\mathbf{q}^+=\mathbf{q}\otimes\delta\mathbf{q}$；注入后名义 $\mathbf{q}^+=\mathbf{q}\otimes\widehat{\delta\mathbf{q}}$。合并并左消：
\[
  \delta\mathbf{q}^+=(\mathbf{q}^+)^*\otimes\mathbf{q}\otimes\delta\mathbf{q}=\widehat{\delta\mathbf{q}}^*\otimes\delta\mathbf{q}=[\widehat{\delta\mathbf{q}}^*]_L\,\delta\mathbf{q}.
\]
取 $\widehat{\delta\mathbf{q}}^*\approx[1;-\tfrac12\widehat{\delta\boldsymbol{\theta}}]$、$\delta\mathbf{q}\approx[1;\tfrac12\delta\boldsymbol{\theta}]$，展开取向量行（弃二阶项）：
\[
  \delta\boldsymbol{\theta}^+=-\widehat{\delta\boldsymbol{\theta}}+\Big(\mathbf{I}-\tfrac12[\widehat{\delta\boldsymbol{\theta}}]_\times\Big)\delta\boldsymbol{\theta}+O(\|\delta\boldsymbol{\theta}\|^2),
\]
故 $\partial\delta\boldsymbol{\theta}^+/\partial\delta\boldsymbol{\theta}=\mathbf{I}-\tfrac12[\widehat{\delta\boldsymbol{\theta}}]_\times$，即 \cref{eq:eskf-reset} 的 $\mathbf{G}$ 姿态块\cite{sola2017quaternion}（全局角误差版差一个符号 $\mathbf{I}+\tfrac12[\widehat{\delta\boldsymbol{\theta}}]_\times$，\cref{tab:eskf-pert}）。多数实现取 $\widehat{\delta\boldsymbol{\theta}}\approx\mathbf{0}\Rightarrow\mathbf{G}\approx\mathbf{I}$，复位平凡。\hfill$\square$
\end{derivation}

\begin{derivation}[（选读）ML 偏差的近似闭式]
对纯测量（无先验）的极大似然 $J(\mathbf{x})=\tfrac12\sum_k(\mathbf{z}_k-\mathbf{h}_k(\mathbf{x}))^\top\boldsymbol{\Sigma}_{\mathbf{v},k}^{-1}(\cdot)$，设估计偏差 $\delta\mathbf{x}=\hat{\mathbf{x}}-\mathbf{x}$ 对噪声至多二次依赖，代入收敛最优条件并对噪声取期望，可得系统偏差\cite{barfoot2024state}
\[
  \mathbb{E}[\delta\mathbf{x}]=-\tfrac12\mathbf{W}(\mathbf{x})^{-1}\sum_k\mathbf{H}_k^\top\boldsymbol{\Sigma}_{\mathbf{v},k}^{-1}\sum_j\mathbf{1}_j\operatorname{tr}\!\big(\boldsymbol{\mathcal{H}}_{jk}\mathbf{W}(\mathbf{x})^{-1}\big),\quad
  \mathbf{W}(\mathbf{x})=\sum_k\mathbf{H}_k^\top\boldsymbol{\Sigma}_{\mathbf{v},k}^{-1}\mathbf{H}_k,
\]
其中 $\boldsymbol{\mathcal{H}}_{jk}=\partial^2 h_{jk}/\partial\mathbf{x}\partial\mathbf{x}^\top$ 是观测分量的 Hessian、$\mathbf{1}_j$ 是单位阵第 $j$ 列。可在收敛后以 $\hat{\mathbf{x}}$ 代 $\mathbf{x}$ 算该式并\emph{去偏} $\hat{\mathbf{x}}\leftarrow\hat{\mathbf{x}}-\mathbb{E}[\delta\mathbf{x}]$（仅在轻非线性时好用）。这量化了 \cref{sec:ekf} 立体相机例中“MAP/ML 有偏”的来源——偏差正比于观测模型的二阶曲率。\hfill$\square$
\end{derivation}
